% Statistiques à deux variables — Mathématiques Tle Tech — Cours signé OnBuch+
% Programme officiel MINESEC. Compiler avec : tectonic N09-statistiques-deux-variables.tex
\newcommand{\DOCMATIERE}{Mathématiques}
\newcommand{\DOCNIVEAU}{Tle Tech}
\newcommand{\DOCMODULE}{Mathématiques – classe de Terminale technique}
\newcommand{\DOCLECON}{N09}
\newcommand{\DOCTITRE}{Statistiques à deux variables}
% OnBuch+ — préambule « cours signés » ENRICHI (compilé avec tectonic / XeLaTeX)
% Système visuel « Pop » d'OnBuch+ : fond crème (jamais blanc), bordures encre
% épaisses, ombres « dures » décalées, titres Archivo Black en capitales.
% Version MATHÉMATIQUES : graphes (pgfplots/tikz), boîtes pédagogiques
% (définition, propriété, méthode, attention, le savais-tu, exercice résolu,
% exercice, TP, bilan), page de garde et signature OnBuch+ ancrée en bas.
%
% Macros attendues dans main.tex AVANT \input{preamble} :
%   \DOCMATIERE  \DOCNIVEAU  \DOCMODULE  \DOCLECON (numéro)  \DOCTITRE
\documentclass[11pt]{article}

\usepackage{fontspec}
\usepackage[french]{babel}
\usepackage[a4paper,margin=2cm,top=3.4cm,bottom=2.8cm,headheight=1.4cm,footskip=1.3cm]{geometry}
\usepackage{amsmath,amssymb,amsfonts}
\usepackage{mathtools} % \xmapsto, \coloneqq... souvent utilisés par les modèles
\usepackage{cancel} % simplifications barrées
% \abs{z} (module, valeur absolue) et \norm{u} (norme), utilisés par les modèles
\providecommand{\abs}{}\let\abs\relax\DeclarePairedDelimiter{\abs}{\lvert}{\rvert}
\providecommand{\norm}{}\let\norm\relax\DeclarePairedDelimiter{\norm}{\lVert}{\rVert}
\usepackage[table]{xcolor} % [table] : fournit aussi \rowcolors (alternance de lignes)
\usepackage{pagecolor}
\usepackage{tikz}
\usetikzlibrary{arrows.meta,positioning,calc,shapes.geometric,shapes.misc,shapes.arrows,decorations.pathmorphing,decorations.pathreplacing,decorations.markings,patterns,shadows,mindmap,trees,fit,backgrounds,matrix,intersections,angles,quotes}
\usepackage{pgfplots}
\usepackage{pgf-pie} % diagrammes circulaires \pie (statistiques)
\pgfplotsset{compat=1.18}
\usepgfplotslibrary{fillbetween,groupplots} % aires entre courbes (calcul intégral)
\usepackage{enumitem}
\usepackage{fancyhdr}
% [most] charge toutes les bibliothèques tcolorbox (listings, minted, raster,
% documentation...) et gonfle inutilement la mémoire TeX sur un document long
% (~300 boîtes) : on ne charge que ce qui est réellement utilisé.
\usepackage[skins,breakable]{tcolorbox}
\usepackage{graphicx}
\usepackage{colortbl}
\usepackage{booktabs}
\usepackage{tabularx}
\usepackage{diagbox} % case d'en-tête diagonale des tableaux à double entrée
% Colonnes extensibles centrée / gauche / droite (C, L, R), souvent utilisées par les modèles
\newcolumntype{C}{>{\centering\arraybackslash}X}
\newcolumntype{L}{>{\raggedright\arraybackslash}X}
\newcolumntype{R}{>{\raggedleft\arraybackslash}X}
\usepackage{array}
\usepackage{multicol}
\usepackage{siunitx}
% parse-numbers=false : accepte aussi \SI{2,5 \times 10^{-3}}{...}
\sisetup{locale=FR,output-decimal-marker={,},group-minimum-digits=4,parse-numbers=false}
\DeclareSIUnit\ohm{\ensuremath{\symup{\Omega}}} % U+2126 absent de la police
\usepackage{qrcode}
% \degree : commande fréquemment utilisée par les modèles pour un angle ou une
% température en dehors de \SI{}{} (non fournie par les paquets chargés).
\providecommand{\degree}{^{\circ}}
% \qed (fin de démonstration), utilisé par les modèles sans amsthm
\providecommand{\qed}{\hfill$\square$}
% \barre : double barre des valeurs interdites dans un tableau de variations (tkz-tab)
\providecommand{\barre}{\ensuremath{\|}}
% environnement proof (amsthm non chargé) utilisé par les modèles
\ifdefined\proof\else\newenvironment{proof}[1][Démonstration]{\par\noindent\textit{#1.}\ }{\qed\par}\fi
\usepackage{lastpage}
\usepackage{hyperref}
\hypersetup{hidelinks}

% ============================================================
% Palette « Pop » OnBuch+ — mêmes hex que la classe P (plus_ui.dart)
% ============================================================
\definecolor{poporange}{HTML}{F59321}
\definecolor{popdark}{HTML}{E07A0C}
\definecolor{popcream}{HTML}{FFF6E8}
\definecolor{popink}{HTML}{1C1714}
\definecolor{poppurple}{HTML}{7A5AE0}
\definecolor{popgreen}{HTML}{2E9E6A}
\definecolor{poppink}{HTML}{E0457B}
\definecolor{popblue}{HTML}{2F7DE1}
\definecolor{popgold}{HTML}{C98A12}
\definecolor{popmuted}{HTML}{8A7F76}
\definecolor{popline}{HTML}{EADFD2}
\definecolor{popgray}{HTML}{8A7F76}
\colorlet{popgrey}{popgray}
% Alias tolérants (le modèle nomme parfois les couleurs différemment)
\colorlet{popwhite}{white}
\colorlet{popblack}{popink}
\colorlet{popred}{poppink}
\colorlet{popyellow}{popgold}
% Teintes claires pour fonds de tableaux / zones
\colorlet{popblueL}{popblue!10!white}
\colorlet{poporangeL}{poporange!14!white}
\colorlet{popgreenL}{popgreen!12!white}
\colorlet{poppurpleL}{poppurple!12!white}
\colorlet{poppinkL}{poppink!12!white}
\colorlet{popgoldL}{popgold!14!white}

\pagecolor{popcream}

% ============================================================
% Typographie
% ============================================================
\defaultfontfeatures{Ligatures=TeX}
\setmainfont{PlusJakartaSans-Medium.ttf}[Path=../../../fonts/,BoldFont=PlusJakartaSans-Bold.ttf]
% Maths en sans-serif (Fira Math) avec chiffres et lettres droites de Plus
% Jakarta, pour que les formules se fondent dans le texte.
\usepackage[math-style=ISO,bold-style=ISO]{unicode-math}
\setmathfont{FiraMath-Regular.otf}[Path=../../../fonts/]
\setmathfont{PlusJakartaSans-Medium.ttf}[Path=../../../fonts/,range={up/num,up/{Latin,latin},\mathup}]
\usepackage{icomma} % virgule décimale sans espace en mode math
% Caractères mathématiques Unicode absents des polices de texte (Plus Jakarta,
% Archivo Black) : rendus par leur équivalent mathématique, en texte comme en
% formule — sinon ils s'affichent en carré vide (ex. « Arithmétique dans ℤ »).
\usepackage{newunicodechar}
\newunicodechar{ℝ}{\ensuremath{\mathbb{R}}}
\newunicodechar{ℤ}{\ensuremath{\mathbb{Z}}}
\newunicodechar{ℂ}{\ensuremath{\mathbb{C}}}
\newunicodechar{ℕ}{\ensuremath{\mathbb{N}}}
\newunicodechar{ℚ}{\ensuremath{\mathbb{Q}}}
\newunicodechar{𝔻}{\ensuremath{\mathbb{D}}}
\newunicodechar{α}{\ensuremath{\alpha}}
\newunicodechar{β}{\ensuremath{\beta}}
\newunicodechar{γ}{\ensuremath{\gamma}}
\newunicodechar{δ}{\ensuremath{\delta}}
\newunicodechar{ε}{\ensuremath{\varepsilon}}
\newunicodechar{θ}{\ensuremath{\theta}}
\newunicodechar{λ}{\ensuremath{\lambda}}
\newunicodechar{μ}{\ensuremath{\mu}}
\newunicodechar{σ}{\ensuremath{\sigma}}
\newunicodechar{τ}{\ensuremath{\tau}}
\newunicodechar{φ}{\ensuremath{\varphi}}
\newunicodechar{ψ}{\ensuremath{\psi}}
\newunicodechar{ω}{\ensuremath{\omega}}
\newunicodechar{Σ}{\ensuremath{\Sigma}}
% majuscules produites par \MakeUppercase dans les titres (α-aminés -> Α)
\newunicodechar{Α}{\ensuremath{\alpha}}
\newunicodechar{Β}{\ensuremath{\beta}}
\newunicodechar{Γ}{\ensuremath{\Gamma}}
\newunicodechar{Δ}{\ensuremath{\Delta}}
\newunicodechar{Ε}{\ensuremath{\varepsilon}}
\newunicodechar{Θ}{\ensuremath{\Theta}}
\newunicodechar{Λ}{\ensuremath{\Lambda}}
\newunicodechar{Μ}{\ensuremath{\mu}}
\newunicodechar{Τ}{\ensuremath{\tau}}
\newunicodechar{Φ}{\ensuremath{\Phi}}
\newunicodechar{Ψ}{\ensuremath{\Psi}}
\newunicodechar{Ω}{\ensuremath{\Omega}}
\newunicodechar{↦}{\ensuremath{\mapsto}}
\newunicodechar{∈}{\ensuremath{\in}}
\newunicodechar{∉}{\ensuremath{\notin}}
\newunicodechar{∧}{\ensuremath{\wedge}}
\newunicodechar{⇔}{\ensuremath{\Leftrightarrow}}
\newunicodechar{⟺}{\ensuremath{\Longleftrightarrow}}
\newunicodechar{⇒}{\ensuremath{\Rightarrow}}
\newunicodechar{⟹}{\ensuremath{\Longrightarrow}}
\newunicodechar{∘}{\ensuremath{\circ}}
\newunicodechar{∪}{\ensuremath{\cup}}
\newunicodechar{∩}{\ensuremath{\cap}}
\newunicodechar{≡}{\ensuremath{\equiv}}
\newunicodechar{⊂}{\ensuremath{\subset}}
\newunicodechar{⊥}{\ensuremath{\perp}}
\newunicodechar{∃}{\ensuremath{\exists}}
\newunicodechar{∀}{\ensuremath{\forall}}
\newunicodechar{∛}{\ensuremath{\sqrt[3]{\,}}}
\newunicodechar{∜}{\ensuremath{\sqrt[4]{\,}}}
\newunicodechar{⃗}{\ensuremath{{}^{\to}}}
\newunicodechar{ⁿ}{\ifmmode^{n}\else\textsuperscript{\ensuremath{n}}\fi}
\newunicodechar{ˣ}{\ifmmode^{x}\else\textsuperscript{\ensuremath{x}}\fi}
\newunicodechar{ᵇ}{\ifmmode^{b}\else\textsuperscript{\ensuremath{b}}\fi}
\newunicodechar{ʳ}{\ifmmode^{r}\else\textsuperscript{\ensuremath{r}}\fi}
\newunicodechar{ᵘ}{\ifmmode^{u}\else\textsuperscript{\ensuremath{u}}\fi}
\newunicodechar{ʸ}{\ifmmode^{y}\else\textsuperscript{\ensuremath{y}}\fi}
\newunicodechar{ᵃ}{\ifmmode^{a}\else\textsuperscript{\ensuremath{a}}\fi}
\newunicodechar{ᵉ}{\ifmmode^{e}\else\textsuperscript{\ensuremath{e}}\fi}
\newunicodechar{ᵏ}{\ifmmode^{k}\else\textsuperscript{\ensuremath{k}}\fi}
\newunicodechar{ᵖ}{\ifmmode^{p}\else\textsuperscript{\ensuremath{p}}\fi}
\newunicodechar{ᴺ}{\ifmmode^{N}\else\textsuperscript{\ensuremath{N}}\fi}
\newunicodechar{ᶜ}{\ifmmode^{c}\else\textsuperscript{\ensuremath{c}}\fi}
\newunicodechar{ʰ}{\ifmmode^{h}\else\textsuperscript{\ensuremath{h}}\fi}
\newunicodechar{ᵠ}{\ifmmode^{q}\else\textsuperscript{\ensuremath{q}}\fi}
\newunicodechar{ₙ}{\ifmmode_{n}\else\textsubscript{\ensuremath{n}}\fi}
\newunicodechar{ₐ}{\ifmmode_{a}\else\textsubscript{\ensuremath{a}}\fi}
\newunicodechar{ᵢ}{\ifmmode_{i}\else\textsubscript{\ensuremath{i}}\fi}
\newunicodechar{ᵣ}{\ifmmode_{r}\else\textsubscript{\ensuremath{r}}\fi}
\newunicodechar{ₖ}{\ifmmode_{k}\else\textsubscript{\ensuremath{k}}\fi}
\newunicodechar{ᵦ}{\ifmmode_{\beta}\else\textsubscript{\ensuremath{\beta}}\fi}
\newunicodechar{ₓ}{\ifmmode_{x}\else\textsubscript{\ensuremath{x}}\fi}
\newfontfamily\popdisplay{ArchivoBlack-Regular.ttf}[Path=../../../fonts/]
\newfontfamily\poplogo{SpaceGrotesk-Bold.ttf}[Path=../../../fonts/]
\newfontfamily\popsemi{PlusJakartaSans-SemiBold.ttf}[Path=../../../fonts/]
\newfontfamily\popxbold{PlusJakartaSans-ExtraBold.ttf}[Path=../../../fonts/]
\newfontfamily\popxboldls{PlusJakartaSans-ExtraBold.ttf}[Path=../../../fonts/,LetterSpace=3.0]

\setlist[enumerate]{leftmargin=1.7em,itemsep=5pt,topsep=4pt}
\setlist[itemize]{leftmargin=1.7em,itemsep=4pt,topsep=4pt,label={\color{poporange}\textbullet}}
\setlength{\parindent}{0pt}
\setlength{\parskip}{5pt}
\renewcommand{\arraystretch}{1.25}

% pgfplots : style Pop par défaut
\pgfplotsset{
  every axis/.append style={axis line style={popink,line width=1pt},tick style={popink},
    grid style={popline,line width=0.5pt},label style={font=\small},tick label style={font=\footnotesize}},
  popcurve/.style={poporange,line width=1.8pt,no markers,smooth},
}
% Même style disponible dans un \draw TikZ simple (hors pgfplots)
\tikzset{popcurve/.style={poporange,line width=1.8pt}}

% ============================================================
% Pill / squiggle
% ============================================================
\newcommand{\poppill}[3]{%
  \tikz[baseline=-0.55ex]\node[draw=popink,line width=1.2pt,rounded corners=2.6pt,%
    fill=#2,inner xsep=6.5pt,inner ysep=3pt]{%
    \popxboldls\fontsize{7}{7}\selectfont\color{#3}\MakeUppercase{#1}};%
}
\newcommand{\popsquiggle}[1]{%
  \tikz[baseline=-0.4ex]{
    \draw[#1,line width=2.6pt,line cap=round]
      plot [smooth] coordinates {(0,0.09) (0.34,0.01) (0.68,0.21) (1.02,0.01) (1.36,0.10)};
  }%
}
% Mot-clé surligné (marqueur orange sous le texte)
\newcommand{\cle}[1]{{\popxbold\color{popdark}#1}}

% ============================================================
% En-tête / pied
% ============================================================
\pagestyle{fancy}
\fancyhf{}
\renewcommand{\headrulewidth}{0pt}
\renewcommand{\footrulewidth}{0pt}
\lhead{\raisebox{-0.15ex}{\poplogo\fontsize{13}{13}\selectfont\color{popink}OnBuch\color{poporange}+}}
\rhead{\poppill{\DOCMATIERE\ \textbullet\ \DOCNIVEAU\ \textbullet\ Leçon \DOCLECON}{white}{popink}}
\lfoot{\popxbold\fontsize{7.6}{7.6}\selectfont\color{popmuted}\MakeUppercase{Cours signé OnBuch+}\ \textbullet\ {\popsemi\DOCTITRE}}
\rfoot{\tikz[baseline=-0.6ex]\node[rounded corners=8.5pt,fill=popink,minimum height=17pt,minimum width=17pt,inner xsep=5pt,inner sep=0]{\popxbold\fontsize{8}{8}\selectfont\color{popcream}\thepage\,/\,\pageref*{LastPage}};}

% ============================================================
% Titres
% ============================================================
\newcommand{\chaptitre}[2]{%
  \begin{tcolorbox}[enhanced,colback=poporange,colframe=popink,boxrule=2.6pt,arc=11pt,
      drop shadow=popink,left=15pt,right=15pt,top=12pt,bottom=11pt,before skip=3pt,after skip=18pt]
    \poppill{#2}{popink}{popcream}\par\vspace{9pt}
    {\raggedright\hyphenpenalty=10000\exhyphenpenalty=10000\popdisplay\fontsize{22}{24}\selectfont\color{popcream}\MakeUppercase{#1}\par}
  \end{tcolorbox}
}
% Section numérotée : pastille orange avec le numéro + titre Archivo
\newcounter{popsec}
\newcommand{\coursec}[1]{%
  \stepcounter{popsec}\setcounter{popsub}{0}%
  \par\vspace{16pt}\noindent
  \begin{minipage}{\linewidth}
  \tikz[baseline=-0.7ex]\node[draw=popink,line width=1.6pt,fill=poporange,rounded corners=4pt,minimum size=22pt,inner sep=0]{\popdisplay\fontsize{12}{12}\selectfont\color{popcream}\thepopsec};%
  \hspace{8pt}{\popdisplay\fontsize{15}{17}\selectfont\color{popink}\MakeUppercase{#1}}\par
  \vspace{4pt}\noindent{\color{poporange}\rule{36pt}{3.6pt}}
  \end{minipage}\par\nopagebreak\vspace{8pt}
}
\newcounter{popsub}
\newcommand{\courssub}[1]{%
  \stepcounter{popsub}%
  \par\vspace{9pt}\noindent{\popxbold\fontsize{11.5}{13}\selectfont\color{popink}{\color{poporange}\thepopsec.\thepopsub}\hspace{6pt}#1}\par\nopagebreak\vspace{4pt}
}

% ============================================================
% Boîtes pédagogiques — même recette : fond blanc, bordure épaisse colorée,
% ombre dure encre, badge plein. Titre optionnel [#1].
% ============================================================
\tcbset{popbase/.style={breakable,enhanced,colback=white,boxrule=2.3pt,arc=9pt,
  shadow={3pt}{-3pt}{0pt}{popink},left=14pt,right=14pt,top=11pt,bottom=11pt,before skip=11pt,after skip=15pt,
  fontupper=\popsemi\fontsize{10.6}{14.5}\selectfont\color{popink},
  fontlower=\popsemi\fontsize{10.6}{14.5}\selectfont\color{popink}}}
% Coche dessinée (le glyphe ✓ n'existe pas dans les polices du document)
\newcommand{\popcheck}[1]{\tikz[baseline=-0.5ex]\draw[#1,line width=1.6pt,line cap=round,line join=round](-0.13,0.0)--(-0.04,-0.09)--(0.13,0.1);}
\newcommand{\popboxhead}[3]{\poppill{#1}{#2}{white}\ifx\relax#3\relax\else\hspace{6pt}{\popxbold\fontsize{10.6}{12}\selectfont\color{popink}#3}\fi\par\vspace{7pt}}

\newtcolorbox{aretenir}[1][]{popbase,colframe=popgreen,before upper={\popboxhead{À retenir}{popgreen}{#1}}}
\newtcolorbox{exemplebox}[1][]{popbase,colframe=poporange,before upper={\popboxhead{Exemple}{poporange}{#1}}}
\newtcolorbox{definition}[1][]{popbase,colframe=popblue,before upper={\popboxhead{Définition}{popblue}{#1}}}
\newtcolorbox{propriete}[1][]{popbase,colframe=poppurple,before upper={\popboxhead{Propriété}{poppurple}{#1}}}
\newtcolorbox{methode}[1][]{popbase,colframe=popgold,before upper={\popboxhead{Méthode}{popgold}{#1}}}
\newtcolorbox{attention}[1][]{popbase,colframe=poppink,before upper={\popboxhead{Attention}{poppink}{#1}}}
\newtcolorbox{experience}[1][]{popbase,colframe=popblue,colback=popblueL,before upper={\popboxhead{Expérience}{popblue}{#1}}}
% « Le savais-tu ? » : carte violette pleine (contexte, culture, Cameroun)
\newtcolorbox{savaistu}[1][]{popbase,colback=poppurple,colframe=popink,
  fontupper=\popsemi\fontsize{10.4}{14.2}\selectfont\color{white},
  before upper={\poppill{Le savais-tu\,?}{white}{poppurple}\ifx\relax#1\relax\else\hspace{6pt}{\popxbold\color{white}#1}\fi\par\vspace{7pt}}}
% Exercice résolu : énoncé en haut, \tcblower puis la solution sur fond crème
\newcounter{popexo}\newcounter{popexr}
\newtcolorbox{exoresolu}[1][]{popbase,colframe=poporange,colbacklower=popcream,
  segmentation style={popink,line width=1.2pt,dashed},
  before upper={\stepcounter{popexr}\popboxhead{Exercice résolu \thepopexr}{poporange}{#1}},
  before lower={\poppill{Solution}{popink}{popcream}\par\vspace{6pt}}}
% Exercice (non résolu en place) — la correction est dans la partie Corrigés
\newtcolorbox{exercice}[1][]{popbase,colframe=popink,
  before upper={\stepcounter{popexo}\popboxhead{Exercice \thepopexo}{popink}{#1}}}
% Corrigé
\newtcolorbox{corrige}[1][]{popbase,colframe=popgreen,colback=popgreenL,
  before upper={\popboxhead{Corrigé}{popgreen}{#1}}}
% Objectifs / prérequis / bilan
\newtcolorbox{objectifs}{popbase,colframe=popink,colback=poporangeL,
  before upper={\popboxhead{Objectifs de la leçon}{popink}{}}}
\newtcolorbox{prerequis}{popbase,colframe=popink,colback=popblueL,
  before upper={\popboxhead{Prérequis}{popblue}{}}}
\newtcolorbox{bilan}{popbase,colframe=popink,colback=white,boxrule=2.8pt,
  before upper={\popboxhead{Fiche bilan}{poporange}{L'essentiel à connaître pour l'examen}}}
% Figure encadrée avec légende
\newtcolorbox{popfigure}[1][]{enhanced,breakable=false,colback=white,colframe=popline,boxrule=1.4pt,arc=8pt,
  left=10pt,right=10pt,top=10pt,bottom=8pt,before skip=10pt,after skip=12pt,halign=center,
  lower separated=false,halign lower=center,
  fontlower=\popsemi\fontsize{9}{11.5}\selectfont\color{popmuted},
  #1}
\newcommand{\legende}[1]{\par\vspace{6pt}{\popsemi\fontsize{9}{11.5}\selectfont\color{popmuted}#1}}

% ============================================================
% Signature OnBuch+
% ============================================================
\newcommand{\poplogoinline}[1]{{\poplogo\fontsize{#1}{#1}\selectfont\color{popink}OnBuch\color{poporange}+}}
\newcommand{\signatureonbuch}{%
  \begin{tcolorbox}[enhanced,colback=popink,colframe=popink,boxrule=2.2pt,arc=10pt,
      drop shadow=poporange,left=14pt,right=14pt,top=10pt,bottom=10pt,before skip=0pt,after skip=0pt]
    \tikz[baseline=-0.7ex]\node[circle,fill=popgreen,draw=popcream,line width=1.3pt,minimum size=22pt,inner sep=0]{\popcheck{white}};%
    \hspace{9pt}\begin{minipage}[c]{0.62\linewidth}
      {\popxbold\fontsize{12}{13}\selectfont\color{popcream}Signé\ \textbullet\ {\poplogo OnBuch\color{poporange}+}}\par\vspace{2pt}
      {\raggedright\popsemi\fontsize{8}{10.5}\selectfont\color{popline}\MakeUppercase{Équipe pédagogique OnBuch+}\\\MakeUppercase{Conforme au programme officiel MINESEC}\par}
    \end{minipage}\hfill
    {\poplogo\fontsize{22}{22}\selectfont\color{popcream}OnBuch\color{poporange}+}
  \end{tcolorbox}%
}

% ============================================================
% Page de garde
% #1 titre · #2 matière · #3 niveau · #4 module · #5 n° leçon · #6 durée ·
% #7 description / objectifs (texte ou itemize)
% ============================================================
\newcommand{\pagegarde}[7]{%
  \thispagestyle{empty}
  \newgeometry{a4paper,margin=1.7cm,top=1.6cm,bottom=1.4cm}%
  \begin{tikzpicture}[remember picture,overlay]
    \fill[poporange!18] ([xshift=-3.2cm,yshift=-3.6cm]current page.north east) circle (4.6cm);
    \fill[poppurple!14] ([xshift=2.4cm,yshift=7.6cm]current page.south west) circle (3.8cm);
  \end{tikzpicture}%
  {\poplogo\fontsize{30}{30}\selectfont\color{popink}OnBuch\color{poporange}+}\hfill
  \raisebox{0.6ex}{\poppill{Cours signé \textbullet\ Premium}{popink}{popcream}}\par
  \vspace{1pt}\popsquiggle{poppurple}\par
  \vspace{8pt}
  \begin{tcolorbox}[enhanced,colback=poporange,colframe=popink,boxrule=2.8pt,arc=13pt,
      drop shadow=popink,left=18pt,right=18pt,top=12pt,bottom=13pt,after skip=10pt]
    \poppill{#2 \textbullet\ #3}{popink}{popcream}\hspace{5pt}\poppill{#4}{white}{popink}\par\vspace{8pt}
    {\popdisplay\fontsize{14}{14}\selectfont\color{popink}LEÇON #5}\par\vspace{4pt}
    {\raggedright\hyphenpenalty=10000\exhyphenpenalty=10000\popdisplay\fontsize{25}{27}\selectfont\color{popcream}\MakeUppercase{#1}\par}
    \vspace{8pt}
    {\popxbold\fontsize{9.5}{9.5}\selectfont\color{popink}\MakeUppercase{Durée conseillée : #6}}
  \end{tcolorbox}
  \begin{tcolorbox}[enhanced,colback=white,colframe=popink,boxrule=2.2pt,arc=10pt,
      drop shadow=popink,left=16pt,right=16pt,top=10pt,bottom=10pt,after skip=8pt]
    \poppill{Dans cette leçon}{poppurple}{white}\par\vspace{5pt}
    \popsemi\fontsize{9.3}{12.6}\selectfont\color{popink}#7
  \end{tcolorbox}
  \noindent\begin{minipage}[t]{0.487\linewidth}
    \begin{tcolorbox}[enhanced,colback=popgreen,colframe=popink,boxrule=2.2pt,arc=10pt,
        drop shadow=popink,left=11pt,right=11pt,top=10pt,bottom=10pt,height=3.1cm,valign=center]
      \tcbox[colback=white,colframe=popink,boxrule=1.3pt,arc=5pt,boxsep=0pt,nobeforeafter,
        left=3pt,right=3pt,top=3pt,bottom=3pt,valign=center]{\qrcode[height=1.9cm]{https://play.google.com/store/apps/details?id=com.luvvix.onbuch&hl=fr}}%
      \hspace{8pt}\begin{minipage}[c]{0.5\linewidth}
        {\popdisplay\fontsize{11}{12.5}\selectfont\color{white}\MakeUppercase{Télécharge OnBuch}}\par\vspace{4pt}
        {\raggedright\popsemi\fontsize{8.4}{11}\selectfont\color{white}Gratuit \textbullet\ Google Play\\Tuteur IA, annales, concours, résultats\par}
      \end{minipage}
    \end{tcolorbox}
  \end{minipage}\hfill
  \begin{minipage}[t]{0.487\linewidth}
    \begin{tcolorbox}[enhanced,colback=poppurple,colframe=popink,boxrule=2.2pt,arc=10pt,
        drop shadow=popink,left=11pt,right=11pt,top=10pt,bottom=10pt,height=3.1cm,valign=center]
      \tikz[baseline=-0.7ex]\node[circle,fill=white,draw=popink,line width=1.3pt,minimum size=1.6cm,inner sep=0]{\popdisplay\fontsize{24}{24}\selectfont\color{poppurple}+};%
      \hspace{8pt}\begin{minipage}[c]{0.6\linewidth}
        {\popdisplay\fontsize{11}{12.5}\selectfont\color{white}\MakeUppercase{Passe à OnBuch+}}\par\vspace{4pt}
        {\raggedright\popsemi\fontsize{8.4}{11}\selectfont\color{white}Tous les cours signés, Tuteur IA illimité, fiches PDF — onglet OnBuch+ de l'app\par}
      \end{minipage}
    \end{tcolorbox}
  \end{minipage}
  \vspace{8pt}
  \popcontenu
  \vfill
  \signatureonbuch
  \restoregeometry
  \newpage
}

% Rangée « Ce cours contient » (page de garde)
\newcommand{\poptile}[3]{% #1 couleur #2 gros texte #3 légende
  \begin{tcolorbox}[enhanced,colback=#1,colframe=popink,boxrule=2pt,arc=9pt,shadow={2.5pt}{-2.5pt}{0pt}{popink},
      left=6pt,right=6pt,top=9pt,bottom=9pt,halign=center,valign=center,height=1.75cm,width=\linewidth,nobeforeafter]
    {\popdisplay\fontsize{12.5}{13}\selectfont\color{white}\MakeUppercase{#2}}\par\vspace{3pt}
    {\centering\popsemi\fontsize{7.8}{9.5}\selectfont\color{white}#3\par}
  \end{tcolorbox}}
\newcommand{\popcontenu}{%
  \par\noindent{\popxboldls\fontsize{7.5}{7.5}\selectfont\color{popmuted}CE COURS SIGNÉ CONTIENT}\par\vspace{7pt}
  \noindent\begin{minipage}[t]{0.235\linewidth}\poptile{popblue}{Cours}{complet, illustré, pas à pas}\end{minipage}\hfill
  \begin{minipage}[t]{0.235\linewidth}\poptile{poporange}{Méthodes}{et exercices résolus}\end{minipage}\hfill
  \begin{minipage}[t]{0.235\linewidth}\poptile{poppink}{Exercices}{type examen + corrigés}\end{minipage}\hfill
  \begin{minipage}[t]{0.235\linewidth}\poptile{popgreen}{Fiche bilan}{l'essentiel à retenir}\end{minipage}\par}

% Sommaire
\newtcolorbox{sommaire}{enhanced,colback=white,colframe=popink,boxrule=2.2pt,arc=10pt,
  drop shadow=popink,left=18pt,right=18pt,top=13pt,bottom=13pt,before skip=6pt,after skip=20pt,
  before upper={\poppill{Sommaire}{poppurple}{white}\par\vspace{9pt}\popsemi\fontsize{10.6}{14.5}\selectfont\color{popink}}}

% Signature de fin de cours — poussée tout en bas de la dernière page
\newcommand{\findecours}{%
  \par\vfill\nopagebreak
  \begin{center}\popsquiggle{poporange}\end{center}\vspace{4pt}
  \signatureonbuch
}
\AtBeginDocument{\def\llbracket{\mathopen{[\mkern-3.5mu[}}\def\rrbracket{\mathclose{]\mkern-3.5mu]}}} % intervalles d'entiers (glyphe ⟦ absent de la police)
% Glyphes absents de Fira Math : redéfinis à partir de glyphes disponibles
\AtBeginDocument{%
  \def\setminus{\mathbin{\backslash}}\let\smallsetminus\setminus
  \def\vdots{\mathord{\vbox{\baselineskip3.2pt\lineskiplimit0pt\kern4pt\hbox{.}\hbox{.}\hbox{.}}}}%
  \def\ddots{\mathinner{\mkern1mu\raise6.5pt\hbox{.}\mkern2mu\raise3.6pt\hbox{.}\mkern2mu\raise0.7pt\hbox{.}\mkern1mu}}%
  \def\triangle{\mathord{\scalebox{0.9}{$\symup{\Delta}$}}}%
  \def\nabla{\mathord{\raisebox{\depth}{\scalebox{1}[-1]{$\symup{\Delta}$}}}}%
  \def\checkmark{\popcheck{popdark}}%
  \def\star{\mathbin{\ast}}\def\bigstar{\mathord{\ast}}%
  \def\diamond{\mathbin{\scalebox{0.6}{$\Diamond$}}}%
  \def\bigcup{\mathop{\mathchoice{\raisebox{-0.3em}{\scalebox{1.7}{$\cup$}}}{\scalebox{1.25}{$\cup$}}{\cup}{\cup}}\displaylimits}%
  \def\bigcap{\mathop{\mathchoice{\raisebox{-0.3em}{\scalebox{1.7}{$\cap$}}}{\scalebox{1.25}{$\cap$}}{\cap}{\cap}}\displaylimits}%
  \def\lesseqqgtr{\mathrel{\lessgtr}}\def\bigoplus{\mathop{\oplus}\displaylimits}%
}
\providecommand{\ssi}{si et seulement si} % abréviation fréquente
\AtBeginDocument{\providecommand{\wideparen}{\overparen}} % arc de cercle

%% FIGFIX-BEGIN v5 — correctifs globaux des figures (docs/onbuch-generation/tools/figfix)
\usepackage{adjustbox}\usepackage{etoolbox}\tcbuselibrary{raster}
% toute figure TikZ/pgfplots plus large que la ligne est réduite à la largeur disponible
\BeforeBeginEnvironment{tikzpicture}{\begin{adjustbox}{max width=\linewidth}}
\AfterEndEnvironment{tikzpicture}{\end{adjustbox}}
% légendes pgfplots lisibles par-dessus les courbes
\pgfplotsset{every axis/.append style={legend style={fill=white,fill opacity=0.92,text opacity=1,draw=black!15,inner sep=3pt}}}
% étiquettes d'arêtes (syntaxe edge["..."]) avec fond blanc
\tikzset{every edge quotes/.append style={fill=white,inner sep=1.2pt}}
%% FIGFIX-END
\begin{document}

\pagegarde{Statistiques à deux variables}{Mathématiques}{Tle Tech}{Mathématiques – classe de Terminale technique}{N09}{2 séances (fiche de progression harmonisée)}{Cette leçon structure l'analyse des séries statistiques à deux variables : organisation en tableaux à double entrée, étude de la liaison via le nuage de points et le coefficient de corrélation, puis modélisation par la droite de régression des moindres carrés pour prévoir et interpréter.
\vspace{4pt}
\begin{multicols}{2}\small
\begin{itemize}
\item À la fin de cette leçon, je sais construire et lire un tableau à double entrée effectifs/effectifs cumulés.
\item Je sais représenter un nuage de points dans un repère orthonormé et en déduire visuellement l'existence et le sens d'une liaison.
\item Je sais calculer et interpréter le coefficient de corrélation linéaire de Pearson (r).
\item Je sais déterminer l'équation de la droite de régression linéaire (moindres carrés) de y en x.
\item Je sais utiliser cette droite pour faire des prévisions (interpolation/extrapolation) en précisant la fiabilité.
\item Je sais justifier le choix du modèle linéaire par l'analyse des résidus et la valeur de r².
\end{itemize}\end{multicols}}

\begin{sommaire}
\begin{multicols}{2}
\begin{enumerate}
\item 1. Organisation des données : Tableau à double entrée
\item 2. Visualisation et mesure de la liaison : Nuage de points \& Coefficient de corrélation
\item 3. Modélisation : Droite de régression des moindres carrés (y en x)
\item 4. Validation du modèle, Prévision et Communication
\item 5. Complément utile : Ajustement non linéaire par linéarisation (Hors programme strict, culture Terminale)
\item Activité d'intégration
\item Méthodes et exercices résolus
\item Exercices
\item Corrigés des exercices
\item Fiche bilan
\end{enumerate}
\end{multicols}
\end{sommaire}

\begin{objectifs}
\begin{itemize}
\item À la fin de cette leçon, je sais construire et lire un tableau à double entrée effectifs/effectifs cumulés.
\item Je sais représenter un nuage de points dans un repère orthonormé et en déduire visuellement l'existence et le sens d'une liaison.
\item Je sais calculer et interpréter le coefficient de corrélation linéaire de Pearson (r).
\item Je sais déterminer l'équation de la droite de régression linéaire (moindres carrés) de y en x.
\item Je sais utiliser cette droite pour faire des prévisions (interpolation/extrapolation) en précisant la fiabilité.
\item Je sais justifier le choix du modèle linéaire par l'analyse des résidus et la valeur de r².
\item Je sais rédiger une conclusion statistique complète et contextualisée pour une situation réelle (agriculture, commerce, démographie).
\item Je sais identifier les limites du modèle (extrapolation abusive, confusion corrélation/causalité).
\end{itemize}
\end{objectifs}

\begin{prerequis}
\begin{itemize}
\item Moyenne, variance, écart-type d'une série statistique à une variable (1ère/Terminale).
\item Équation réduite d'une droite, passage par deux points, interprétation du coefficient directeur (2nde).
\item Calcul de covariance et utilisation de la calculatrice (mode STAT) pour les sommes ∑x, ∑y, ∑xy, ∑x², ∑y².
\item Notion de pourcentage et taux d'évolution (4ème/2nde).
\item Lecture de coordonnées dans un repère orthonormé.
\end{itemize}
\end{prerequis}

\newpage

\coursec{Organisation des données : Tableau à double entrée}

\textbf{Situation d'accroche.} À la direction régionale de l'agriculture du Littoral, à Douala, on cherche à prévoir la production de plantains (en tonnes) en fonction des précipitations annuelles (en millimètres) observées sur les dix dernières années. Les registres de l'antenne agricole contiennent deux colonnes de chiffres : pour chaque année, la pluie tombée et la récolte obtenue. Dix lignes, vingt nombres, et aucune tendance visible à l'œil nu. Comment organiser ces informations pour mettre en évidence une tendance ? Comment quantifier le lien entre pluie et rendement, et proposer une prévision fiable pour la campagne prochaine ?

\textbf{Problématique.} Avant de calculer quoi que ce soit, il faut \emph{ranger} les données. C'est le rôle du \cle{tableau à double entrée} : un outil de classement qui croise deux caractères quantitatifs et qui constitue le point de départ de toute l'étude statistique à deux variables.

\courssub{Série statistique à deux variables}

En Première, on étudiait une \cle{série statistique à une variable} : on observait un seul caractère (par exemple la taille des élèves d'une classe) sur une population. En Terminale, on s'intéresse simultanément à \emph{deux} caractères mesurés sur les \emph{mêmes} individus, car on soupçonne qu'ils sont liés : la pluie influence-t-elle la récolte ? Le prix d'un article influence-t-il les ventes ? Le nombre d'heures de fonctionnement d'une machine influence-t-il le nombre de pannes ?

\begin{definition}[Série statistique à deux variables]
Soit une \cle{population} de $n$ individus. Une \cle{série statistique à deux variables} est la donnée, pour chaque individu, des valeurs de deux caractères \cle{quantitatifs} $X$ et $Y$. On la note $(X, Y)$ ou $(x_i ; y_i)$ pour l'individu numéro $i$, avec $i$ allant de $1$ à $n$.
\end{definition}

Dans notre situation : la population est l'ensemble des \cle{dix campagnes agricoles} (les dix dernières années) ; le caractère $X$ est la \cle{précipitation annuelle} (en mm) ; le caractère $Y$ est la \cle{production de plantains} (en tonnes). Chaque année fournit un couple $(x_i ; y_i)$.

\begin{attention}[Deux caractères, un seul individu]
Les deux valeurs d'un couple $(x_i ; y_i)$ doivent porter sur le \textbf{même individu} (ici, la même année). On ne peut pas mélanger la pluie de 2015 avec la récolte de 2018 : le couple perdrait tout sens. C'est ce couplage qui permettra ensuite d'étudier une éventuelle liaison.
\end{attention}

\courssub{Le tableau de croisement (double entrée)}

\textbf{Intuition.} Quand les valeurs de $X$ et de $Y$ sont nombreuses et étalées, on les regroupe en \cle{classes} (intervalles), exactement comme on le faisait pour une seule variable. Mais au lieu d'une seule ligne de classes, on croise les classes de $X$ (en lignes) avec les classes de $Y$ (en colonnes) : chaque couple $(x_i ; y_i)$ tombe alors dans une \emph{case} du tableau, et on compte combien d'individus tombent dans chaque case.

\begin{definition}[Tableau de croisement]
Un \cle{tableau de croisement} (ou \cle{tableau à double entrée}, ou tableau de contingence) est un tableau dont :
\begin{itemize}
\item les \textbf{lignes} correspondent aux classes (ou valeurs) du caractère $X$ ;
\item les \textbf{colonnes} correspondent aux classes (ou valeurs) du caractère $Y$ ;
\item la case située à l'intersection de la ligne $i$ et de la colonne $j$ contient l'\cle{effectif partiel} $n_{ij}$ : le nombre d'individus pour lesquels $X$ appartient à la classe $i$ \textbf{et} $Y$ appartient à la classe $j$.
\end{itemize}
\end{definition}

\begin{definition}[Effectifs marginaux et effectif total]
On appelle :
\begin{itemize}
\item \cle{effectif marginal en ligne} $n_{i.}$ : la somme des effectifs de la ligne $i$ (nombre d'individus de la classe $i$ de $X$, toutes valeurs de $Y$ confondues) ;
\item \cle{effectif marginal en colonne} $n_{.j}$ : la somme des effectifs de la colonne $j$ ;
\item \cle{effectif total} $n$ : le nombre total d'individus.
\end{itemize}
On a les relations fondamentales :
\[ n_{i.} = \sum_{j} n_{ij} \qquad n_{.j} = \sum_{i} n_{ij} \qquad n = \sum_{i}\sum_{j} n_{ij} = \sum_{i} n_{i.} = \sum_{j} n_{.j} \]
\end{definition}

\begin{aretenir}[La triple vérification]
La somme de tous les effectifs partiels $n_{ij}$, la somme des effectifs marginaux en ligne $n_{i.}$ et la somme des effectifs marginaux en colonne $n_{.j}$ donnent \textbf{toutes les trois} l'effectif total $n$. C'est un contrôle de cohérence indispensable avant tout calcul.
\end{aretenir}

\begin{methode}[Construire un tableau à double entrée à partir de données brutes]
\begin{enumerate}
\item \textbf{Repérer} les deux caractères $X$ et $Y$ et l'effectif total $n$.
\item \textbf{Choisir les classes} de $X$ et de $Y$ : intervalles de même amplitude, sans recouvrement, couvrant toutes les valeurs observées (3 à 6 classes en général).
\item \textbf{Compter} : pour chaque individu, repérer la case (ligne, colonne) correspondant à son couple $(x_i ; y_i)$ et ajouter $1$ à $n_{ij}$.
\item \textbf{Calculer les marges} : sommer chaque ligne ($n_{i.}$) et chaque colonne ($n_{.j}$), puis vérifier que les deux totaux valent $n$.
\item \textbf{Compléter} par les fréquences et les moyennes conditionnelles si l'énoncé le demande.
\end{enumerate}
\end{methode}

\courssub{Exemple chiffré : pluie et production de plantains}

\begin{exemplebox}[Données brutes sur dix campagnes]
Les dix couples (précipitations en mm ; production en tonnes) relevés sont :
\begin{center}
(1200 ; 8) \quad (1350 ; 10) \quad (1450 ; 12) \quad (1600 ; 15) \quad (1700 ; 16)\\
(1750 ; 17) \quad (1900 ; 18) \quad (2000 ; 20) \quad (2100 ; 19) \quad (2200 ; 22)
\end{center}
On choisit 4 classes de pluie (amplitude 300) : $[1150 ; 1450[$, $[1450 ; 1750[$, $[1750 ; 2050[$, $[2050 ; 2350[$ ; et 3 classes de production (amplitude 6) : $[5 ; 11[$, $[11 ; 17[$, $[17 ; 23[$.
\end{exemplebox}

\begin{methode}[Comptage pas à pas (exemple)]
\begin{enumerate}
\item $(1200 ; 8)$ : Pluie $\in [1150 ; 1450[$, Prod $\in [5 ; 11[$ $\rightarrow$ Case (1,1) $+1$.
\item $(1350 ; 10)$ : Pluie $\in [1150 ; 1450[$, Prod $\in [5 ; 11[$ $\rightarrow$ Case (1,1) $+1$.
\item $(1450 ; 12)$ : Pluie $\in [1450 ; 1750[$, Prod $\in [11 ; 17[$ $\rightarrow$ Case (2,2) $+1$.
\item $(1600 ; 15)$ : Pluie $\in [1450 ; 1750[$, Prod $\in [11 ; 17[$ $\rightarrow$ Case (2,2) $+1$.
\item $(1700 ; 16)$ : Pluie $\in [1450 ; 1750[$, Prod $\in [11 ; 17[$ $\rightarrow$ Case (2,2) $+1$.
\item $(1750 ; 17)$ : Pluie $\in [1750 ; 2050[$, Prod $\in [17 ; 23[$ $\rightarrow$ Case (3,3) $+1$.
\item $(1900 ; 18)$ : Pluie $\in [1750 ; 2050[$, Prod $\in [17 ; 23[$ $\rightarrow$ Case (3,3) $+1$.
\item $(2000 ; 20)$ : Pluie $\in [1750 ; 2050[$, Prod $\in [17 ; 23[$ $\rightarrow$ Case (3,3) $+1$.
\item $(2100 ; 19)$ : Pluie $\in [2050 ; 2350[$, Prod $\in [17 ; 23[$ $\rightarrow$ Case (4,3) $+1$.
\item $(2200 ; 22)$ : Pluie $\in [2050 ; 2350[$, Prod $\in [17 ; 23[$ $\rightarrow$ Case (4,3) $+1$.
\end{enumerate}
\end{methode}

On obtient le tableau complet ci-dessous, avec les marges et les moyennes conditionnelles de production (calculées à partir des valeurs brutes de chaque ligne).

\begin{center}
\begin{tabularx}{\linewidth}{|l|X|X|X|c|c|}
\hline
\rowcolor{popblueL} Pluie $\backslash$ Production & $[5 ; 11[$ & $[11 ; 17[$ & $[17 ; 23[$ & $n_{i.}$ & $\bar{y}_{/i}$ (t) \\
\hline
$[1150 ; 1450[$ & 2 & 0 & 0 & 2 & 9 \\
\hline
$[1450 ; 1750[$ & 0 & 3 & 0 & 3 & 16 \\
\hline
$[1750 ; 2050[$ & 0 & 0 & 3 & 3 & 18{,}3 \\
\hline
$[2050 ; 2350[$ & 0 & 0 & 2 & 2 & 20{,}5 \\
\hline
\rowcolor{popgoldL} $n_{.j}$ & 2 & 3 & 5 & \textbf{10} & \\
\hline
\end{tabularx}
\end{center}

\textbf{Vérification de cohérence.} Somme des lignes : $2+3+3+2 = 10$. Somme des colonnes : $2+3+5 = 10$. Total $n=10$. Le contrôle est validé.

\courssub{Fréquences partielles et marginales}

\begin{definition}[Fréquences]
La \cle{fréquence partielle} de la case $(i,j)$ est $f_{ij} = \dfrac{n_{ij}}{n}$. La \cle{fréquence marginale} de la ligne $i$ est $f_{i.} = \dfrac{n_{i.}}{n}$, et celle de la colonne $j$ est $f_{.j} = \dfrac{n_{.j}}{n}$. La somme de toutes les fréquences partielles vaut $1$.
\end{definition}

Par exemple, la case (ligne $[1750 ; 2050[$, colonne $[17 ; 23[$) a pour fréquence $f_{33} = \dfrac{3}{10} = 0{,}3$ : $30\,\%$ des campagnes ont connu à la fois une pluie entre \SI{1750}{mm} et \SI{2050}{mm} et une production entre 17 et 23 tonnes.

\courssub{Distributions conditionnelles et moyennes conditionnelles}

\textbf{Intuition.} Pour savoir si la pluie influence la production, on « fige » une classe de pluie et on regarde comment la production se répartit \emph{à l'intérieur de cette ligne}. C'est cela, une distribution conditionnelle : on ne regarde plus toute la population, mais seulement un sous-groupe.

\begin{definition}[Distribution conditionnelle]
La \cle{distribution conditionnelle} de $Y$ sachant que $X$ appartient à la classe $i$ est la répartition des effectifs de la ligne $i$ entre les colonnes, exprimée en fréquences conditionnelles :
\[ f_{j/i} = \dfrac{n_{ij}}{n_{i.}} \]
On définit de même la distribution conditionnelle de $X$ sachant la classe $j$ de $Y$ : $f_{i/j} = \dfrac{n_{ij}}{n_{.j}}$.
\end{definition}

\begin{attention}[Le bon dénominateur]
Dans une fréquence conditionnelle $f_{j/i}$, on divise par l'effectif \textbf{marginal de la ligne} $n_{i.}$, et non par l'effectif total $n$. Confondre $f_{ij} = \dfrac{n_{ij}}{n}$ et $f_{j/i} = \dfrac{n_{ij}}{n_{i.}}$ est l'erreur la plus fréquente au Baccalauréat technique.
\end{attention}

\begin{exoresolu}[{Moyenne conditionnelle pour la classe de pluie $[1450 ; 1750[$}]
Énoncé. Calculer la moyenne de production des campagnes dont la pluie est tombée dans la classe $[1450 ; 1750[$ mm, puis interpréter.
\tcblower
Solution détaillée. La ligne $[1450 ; 1750[$ contient trois campagnes (effectif marginal $n_{2.}=3$), toutes dans la colonne $[11 ; 17[$. Les productions réelles correspondantes (données brutes) sont 15, 16 et 17 tonnes. La moyenne conditionnelle exacte est :
\[ \bar{y}_{/[1450;1750[} = \dfrac{15 + 16 + 17}{3} = \dfrac{48}{3} = 16 \text{ tonnes}. \]
Si l'on ne disposait que du tableau groupé, on utiliserait le centre de la classe $[11 ; 17[$, soit $14$, et l'on obtiendrait $\bar{y} = \dfrac{3 \times 14}{3} = 14$ tonnes : c'est une valeur approchée (perte d'information due au regroupement).
Interprétation : lorsque la pluie annuelle est comprise entre \SI{1450}{mm} et \SI{1750}{mm}, la production moyenne de plantains est de 16 tonnes. Comparée à la moyenne conditionnelle de la classe $[1150 ; 1450[$ (9 tonnes), elle suggère que \textbf{plus il pleut, plus la production augmente} : c'est un premier indicateur de liaison entre $X$ et $Y$.
\end{exoresolu}

L'évolution des moyennes conditionnelles (9 ; 16 ; 18,3 ; 20,5) lorsque la pluie augmente est le \cle{premier signal} d'une liaison croissante entre les deux variables. On la visualise avec des barres empilées montrant, pour trois classes de pluie, la répartition conditionnelle de la production.

\begin{popfigure}
\begin{tikzpicture}
\begin{axis}[
  ybar stacked, bar width=0.9cm, width=0.85\linewidth, height=6.5cm,
  ymin=0, ymax=3.5,
  symbolic x coords={$[1150;1450[$,$[1450;1750[$,$[1750;2050[$,$[2050;2350[$},
  xtick=data, x tick label style={font=\small, rotate=45, anchor=east},
  ylabel={Effectif conditionnel}, ylabel style={font=\small},
  legend style={at={(1.02,1)}, anchor=north west, font=\small},
  axis lines=left, enlarge x limits=0.2]
\addplot+[ybar, fill=popblue] coordinates {($[1150;1450[$,2) ($[1450;1750[$,0) ($[1750;2050[$,0) ($[2050;2350[$,0)};
\addplot+[ybar, fill=poporange] coordinates {($[1150;1450[$,0) ($[1450;1750[$,3) ($[1750;2050[$,0) ($[2050;2350[$,0)};
\addplot+[ybar, fill=popgreen] coordinates {($[1150;1450[$,0) ($[1450;1750[$,0) ($[1750;2050[$,3) ($[2050;2350[$,2)};
\legend{$[5;11[$ t, $[11;17[$ t, $[17;23[$ t}
\end{axis}
\end{tikzpicture}
\legende{Distributions conditionnelles de la production pour quatre classes de précipitations : la production « monte » en même temps que la pluie.}
\end{popfigure}

\courssub{Du tableau brut au tableau de valeurs $(x_i, y_i)$}

Le tableau à double entrée est idéal pour \emph{résumer} de nombreuses données. Mais pour tracer le \cle{nuage de points} (section suivante), on a besoin des couples $(x_i ; y_i)$. Deux cas se présentent :

\begin{itemize}
\item \textbf{Données brutes disponibles et effectif faible} (comme nos 10 campagnes) : on garde les couples exacts, par exemple $(1600 ; 15)$. Chaque couple deviendra un point du nuage.
\item \textbf{Données groupées en classes} : on remplace chaque classe par son \cle{centre} ; la case $(i,j)$ fournit alors le couple (centre de la classe $i$ de $X$ ; centre de la classe $j$ de $Y$), compté $n_{ij}$ fois. On perd en précision, mais on peut quand même construire le nuage et calculer les indicateurs.
\end{itemize}

\begin{attention}[Effectifs partiels ou marginaux ?]
Erreur fréquente : confondre $n_{ij}$ (une case) et $n_{i.}$ (une ligne entière). Retenir : l'effectif \textbf{partiel} concerne \emph{une seule case} (un couple de classes), l'effectif \textbf{marginal} concerne \emph{toute une marge} (ligne ou colonne). Et la somme des partiels égale la somme des marginaux, égale à $n$.
\end{attention}

\begin{savaistu}[Et le test du Chi-2 ?]
Le tableau de contingence que tu viens de construire est aussi l'outil de base du célèbre \textbf{test du $\chi^2$ (Chi-2) d'indépendance}, utilisé en médecine, en agronomie et en sciences sociales pour décider si deux caractères sont liés ou indépendants. Ce test se fait en comparant les effectifs observés $n_{ij}$ aux effectifs théoriques $\dfrac{n_{i.} \times n_{.j}}{n}$ que l'on aurait si $X$ et $Y$ étaient parfaitement indépendants. En Terminale technique, on se contente d'\emph{décrire} la liaison à l'aide du tableau et des moyennes conditionnelles ; la mesure précise de la liaison (coefficient de corrélation) arrive dans la section suivante.
\end{savaistu}

\begin{exercice}[À toi de jouer]
Un atelier de Douala relève, pour 8 machines, le couple (âge en années ; nombre de pannes annuel) : $(1 ; 2)$, $(2 ; 3)$, $(3 ; 3)$, $(4 ; 5)$, $(5 ; 6)$, $(6 ; 7)$, $(7 ; 8)$, $(8 ; 9)$.
\begin{enumerate}
\item Construire le tableau à double entrée avec les classes d'âge $[0 ; 3[$, $[3 ; 6[$, $[6 ; 9[$ et les classes de pannes $[0 ; 4[$, $[4 ; 7[$, $[7 ; 10[$.
\item Calculer les effectifs marginaux et vérifier que leur somme vaut $8$.
\item Calculer la moyenne conditionnelle du nombre de pannes pour la classe d'âge $[6 ; 9[$. Que suggère-t-elle ?
\end{enumerate}
\end{exercice}

\begin{corrige}[Exercice : atelier de Douala]
1. En comptant case par case : ligne $[0 ; 3[$ : $(1;2)$ et $(2;3)$ donnent 2 dans la colonne $[0 ; 4[$ ; ligne $[3 ; 6[$ : $(3;3)$ en colonne $[0;4[$, $(4;5)$ et $(5;6)$ en colonne $[4;7[$ ; ligne $[6 ; 9[$ : $(6;7)$, $(7;8)$, $(8;9)$ en colonne $[7;10[$ (attention : 7 appartient à $[7;10[$).
\begin{center}
\begin{tabular}{|l|c|c|c|c|}
\hline
\rowcolor{popblueL} Âge $\backslash$ Pannes & $[0;4[$ & $[4;7[$ & $[7;10[$ & $n_{i.}$ \\
\hline
$[0 ; 3[$ & 2 & 0 & 0 & 2 \\
\hline
$[3 ; 6[$ & 1 & 2 & 0 & 3 \\
\hline
$[6 ; 9[$ & 0 & 0 & 3 & 3 \\
\hline
\rowcolor{popgoldL} $n_{.j}$ & 3 & 2 & 3 & \textbf{8} \\
\hline
\end{tabular}
\end{center}
2. Somme des lignes : $2 + 3 + 3 = 8$ ; somme des colonnes : $3 + 2 + 3 = 8$. Le contrôle est validé.
3. Pour la classe $[6 ; 9[$ : les pannes sont 7, 8, 9. $\bar{y} = \dfrac{7 + 8 + 9}{3} = 8$ pannes. Les machines les plus âgées tombent en panne bien plus souvent (8 pannes contre 2,5 pour la classe $[0 ; 3[$) : l'âge et les pannes semblent liés positivement.
\end{corrige}

\newpage
\coursec{Visualisation et mesure de la liaison : Nuage de points \& Coefficient de corrélation}

\courssub{Le nuage de points}

\begin{definition}[Nuage de points]
Soit une série statistique à deux variables $(x_i ; y_i)_{1 \le i \le n}$. Le \cle{nuage de points} est l'ensemble des points $M_i(x_i ; y_i)$ représentés dans un repère orthonormé $(O ; \vec{i}, \vec{j})$ où l'axe des abscisses porte la variable $X$ et l'axe des ordonnées la variable $Y$.
\end{definition}

\begin{methode}[Tracer un nuage de points]
\begin{enumerate}
\item Choisir une échelle adaptée sur chaque axe (origines non nulles si nécessaire, graduation régulière).
\item Repérer chaque couple $(x_i ; y_i)$ et placer un point.
\item \textbf{Ne pas relier les points} : ce n'est pas une fonction, c'est une observation statistique.
\end{enumerate}
\end{methode}

\begin{exemplebox}[Nuage de points : Pluie vs Production]
On utilise les 10 couples bruts de l'exemple précédent.
\end{exemplebox}

\begin{popfigure}
\begin{tikzpicture}
\begin{axis}[
  width=0.9\linewidth, height=7cm,
  xlabel={Précipitations (mm)}, ylabel={Production (tonnes)},
  xmin=1100, xmax=2300, ymin=5, ymax=24,
  xtick={1200,1400,1600,1800,2000,2200}, ytick={5,10,15,20},
  grid=major, grid style={popline}, axis lines=middle,
  scatter, only marks, mark=*, mark size=2.5, popblue
]
\addplot coordinates {
  (1200,8) (1350,10) (1450,12) (1600,15) (1700,16)
  (1750,17) (1900,18) (2000,20) (2100,19) (2200,22)
};
\end{axis}
\end{tikzpicture}
\legende{Nuage de points des 10 campagnes. On devine une forme allongée selon une droite croissante : liaison linéaire positive.}
\end{popfigure}

\begin{aretenir}[Lecture du nuage de points]
\begin{itemize}
\item \textbf{Forme allongée selon une droite croissante} $\rightarrow$ liaison \textbf{linéaire positive} ($Y$ augmente quand $X$ augmente).
\item \textbf{Forme allongée selon une droite décroissante} $\rightarrow$ liaison \textbf{linéaire négative} ($Y$ diminue quand $X$ augmente).
\item \textbf{Forme circulaire / nuage diffus sans direction} $\rightarrow$ \textbf{pas de liaison linéaire} (indépendance ou liaison non linéaire).
\item \textbf{Points isolés (valeurs aberrantes)} $\rightarrow$ peuvent fausser l'analyse, à commenter.
\end{itemize}
\end{aretenir}

\courssub{Coefficient de corrélation linéaire (de Pearson)}

Le nuage donne une impression visuelle. Le \cle{coefficient de corrélation linéaire} $r$ (ou $r_{XY}$) donne une \emph{mesure numérique} de l'intensité de la liaison linéaire.

\begin{definition}[Coefficient de corrélation linéaire]
\[ r = \frac{\text{cov}(X,Y)}{\sigma_X \sigma_Y} = \frac{\frac{1}{n}\sum_{i=1}^n (x_i - \bar{x})(y_i - \bar{y})}{\sqrt{\frac{1}{n}\sum_{i=1}^n (x_i - \bar{x})^2} \sqrt{\frac{1}{n}\sum_{i=1}^n (y_i - \bar{y})^2}} \]
Formule de calcul pratique (à la calculatrice ou par tableau) :
\[ r = \frac{n\sum x_i y_i - (\sum x_i)(\sum y_i)}{\sqrt{n\sum x_i^2 - (\sum x_i)^2} \sqrt{n\sum y_i^2 - (\sum y_i)^2}} \]
\end{definition}

\begin{propriete}[Propriétés de $r$]
\begin{itemize}
\item $-1 \le r \le 1$ toujours.
\item $r = 1$ : liaison linéaire positive parfaite (tous les points sur une droite croissante).
\item $r = -1$ : liaison linéaire négative parfaite (tous les points sur une droite décroissante).
\item $r = 0$ : \textbf{aucune liaison linéaire} (mais il peut y avoir une liaison non linéaire !).
\item $|r|$ proche de 1 $\rightarrow$ liaison linéaire forte.
\item $|r|$ proche de 0 $\rightarrow$ liaison linéaire faible ou nulle.
\item $r$ ne dépend pas des unités de $X$ et $Y$ (adimensionnel).
\item $r$ est sensible aux valeurs aberrantes (outliers).
\end{itemize}
\end{propriete}

\begin{attention}[Corrélation $\neq$ Causalité]
Un $r$ proche de 1 ne prouve \textbf{pas} que $X$ cause $Y$. Il y a liaison, pas nécessairement causalité. Exemple : le nombre de vendeurs de glaces et le nombre de noyades sont corrélés (cause commune : la chaleur).
\end{attention}

\begin{exoresolu}[Calcul de $r$ pour les 10 campagnes]
Énoncé. Calculer le coefficient de corrélation entre Précipitations ($X$) et Production ($Y$).
\tcblower
Solution détaillée. On construit le tableau de calcul :
\begin{center}
\begin{tabular}{|c|c|c|c|c|c|}
\hline
$x_i$ & $y_i$ & $x_i^2$ & $y_i^2$ & $x_i y_i$ \\
\hline
1200 & 8 & 1\,440\,000 & 64 & 9\,600 \\
1350 & 10 & 1\,822\,500 & 100 & 13\,500 \\
1450 & 12 & 2\,102\,500 & 144 & 17\,400 \\
1600 & 15 & 2\,560\,000 & 225 & 24\,000 \\
1700 & 16 & 2\,890\,000 & 256 & 27\,200 \\
1750 & 17 & 3\,062\,500 & 289 & 29\,750 \\
1900 & 18 & 3\,610\,000 & 324 & 34\,200 \\
2000 & 20 & 4\,000\,000 & 400 & 40\,000 \\
2100 & 19 & 4\,410\,000 & 361 & 39\,900 \\
2200 & 22 & 4\,840\,000 & 484 & 48\,400 \\
\hline
\textbf{Total} & \textbf{17\,250} & \textbf{157} & \textbf{30\,737\,500} & \textbf{2\,647} & \textbf{283\,950} \\
\hline
\end{tabular}
\end{center}
$n=10$, $\sum x = 17250$, $\sum y = 157$, $\sum x^2 = 30\,737\,500$, $\sum y^2 = 2647$, $\sum xy = 283\,950$.
\[ r = \frac{10 \times 283\,950 - 17\,250 \times 157}{\sqrt{10 \times 30\,737\,500 - 17\,250^2} \sqrt{10 \times 2\,647 - 157^2}} \]
Numérateur : $2\,839\,500 - 2\,708\,250 = 131\,250$.
Dénominateur $X$ : $\sqrt{307\,375\,000 - 297\,562\,500} = \sqrt{9\,812\,500} \approx 3132,5$.
Dénominateur $Y$ : $\sqrt{26\,470 - 24\,649} = \sqrt{1\,821} \approx 42,67$.
$r \approx \dfrac{131\,250}{3132,5 \times 42,67} \approx \dfrac{131\,250}{133\,660} \approx 0,982$.
\end{exoresolu}

\begin{aretenir}[Interprétation au Probatoire]
$r \approx 0,98$ : liaison linéaire \textbf{très forte et positive}. Plus les précipitations sont abondantes, plus la production de plantains est élevée. On peut envisager un ajustement linéaire fiable.
\end{aretenir}

\begin{exercice}[Nuage et corrélation]
Pour 6 élèves d'une classe de Terminale, on note le nombre d'heures de révision hebdomadaires ($X$) et la note au devoir de maths sur 20 ($Y$) :
$(2 ; 8)$, $(3 ; 10)$, $(4 ; 11)$, $(5 ; 13)$, $(6 ; 14)$, $(7 ; 15)$.
\begin{enumerate}
\item Tracer le nuage de points dans un repère.
\item Calculer le coefficient de corrélation $r$. Interpréter.
\end{enumerate}
\end{exercice}

\begin{corrige}[Exercice : Révision et notes]
1. Nuage : points alignés grossièrement sur une droite croissante.
2. Tableau : $\sum x = 27$, $\sum y = 71$, $\sum x^2 = 139$, $\sum y^2 = 875$, $\sum xy = 348$, $n=6$.
$r = \frac{6\times 348 - 27\times 71}{\sqrt{6\times 139 - 27^2}\sqrt{6\times 875 - 71^2}} = \frac{2088 - 1917}{\sqrt{834 - 729}\sqrt{5250 - 5041}} = \frac{171}{\sqrt{105}\sqrt{209}} \approx \frac{171}{10,25 \times 14,46} \approx \frac{171}{148,2} \approx 0,978$.
Liaison linéaire très forte positive : plus on révise, meilleure est la note.
\end{corrige}

\newpage
\coursec{Modélisation : Droite de régression des moindres carrés ($y$ en $x$)}

\courssub{Principe de l'ajustement linéaire}

Lorsque le nuage de points suggère une liaison linéaire (forme allongée, $|r|$ proche de 1), on cherche à le \emph{modéliser} par une droite d'équation $y = ax + b$. Cette droite doit être « la plus proche possible » de l'ensemble des points.

\begin{definition}[Méthode des moindres carrés (MMC)]
La \cle{droite de régression de $Y$ en $X$} (ou droite des moindres carrés) est la droite $y = ax + b$ qui minimise la somme des carrés des écarts verticaux (résidus) entre les points observés $M_i(x_i ; y_i)$ et la droite :
\[ S(a,b) = \sum_{i=1}^n (y_i - (ax_i + b))^2 \quad \text{est minimale}. \]
\end{definition}

\begin{propriete}[Coefficients de la droite de régression $y = ax + b$]
La droite solution passe par le \cle{point moyen} $G(\bar{x} ; \bar{y})$.
Son coefficient directeur (pente) est :
\[ a = \frac{\text{cov}(X,Y)}{\text{Var}(X)} = \frac{\frac{1}{n}\sum (x_i - \bar{x})(y_i - \bar{y})}{\frac{1}{n}\sum (x_i - \bar{x})^2} = r \frac{\sigma_Y}{\sigma_X} \]
Son ordonnée à l'origine est :
\[ b = \bar{y} - a\bar{x} \]
Formules de calcul direct (tableau) :
\[ a = \frac{n\sum x_i y_i - (\sum x_i)(\sum y_i)}{n\sum x_i^2 - (\sum x_i)^2} \qquad b = \frac{\sum y_i - a\sum x_i}{n} \]
\end{propriete}

\begin{attention}[Sens de la régression]
La droite de régression de $Y$ en $X$ (prédire $Y$ connaissant $X$) n'est \textbf{pas} la même que la droite de régression de $X$ en $Y$ (prédire $X$ connaissant $Y$), sauf si $r = \pm 1$. Au Probatoire, on demande presque toujours « la droite de régression de $Y$ en $X$ » (équation $y = ax+b$).
\end{attention}

\begin{exoresolu}[Droite de régression Pluie $\rightarrow$ Production]
Énoncé. Déterminer l'équation de la droite de régression de la production ($Y$) en fonction des précipitations ($X$) pour les 10 campagnes. Tracer la droite sur le nuage.
\tcblower
Solution détaillée. On réutilise les sommes de l'exercice précédent :
$n=10$, $\sum x = 17250$, $\sum y = 157$, $\sum x^2 = 30\,737\,500$, $\sum xy = 283\,950$.
Moyennes : $\bar{x} = 1725$, $\bar{y} = 15,7$.
\[ a = \frac{10 \times 283\,950 - 17\,250 \times 157}{10 \times 30\,737\,500 - 17\,250^2} = \frac{131\,250}{9\,812\,500} \approx 0,013376 \text{ t/mm} \]
\[ b = \bar{y} - a\bar{x} = 15,7 - 0,013376 \times 1725 \approx 15,7 - 23,07 \approx -7,37 \text{ t} \]
Équation : $\boxed{y = 0,0134\, x - 7,37}$ (arrondi au $10^{-4}$ pour $a$, au centième pour $b$).
Vérification : $G(1725 ; 15,7)$ appartient à la droite : $0,013376 \times 1725 - 7,37 \approx 15,7$. OK.
\end{exoresolu}

\begin{popfigure}
\begin{tikzpicture}
\begin{axis}[
  width=0.9\linewidth, height=7cm,
  xlabel={Précipitations (mm)}, ylabel={Production (tonnes)},
  xmin=1100, xmax=2300, ymin=5, ymax=24,
  xtick={1200,1400,1600,1800,2000,2200}, ytick={5,10,15,20},
  grid=major, grid style={popline}, axis lines=middle,
]
% Nuage
\addplot[only marks, mark=*, mark size=2.5, popblue] coordinates {
  (1200,8) (1350,10) (1450,12) (1600,15) (1700,16)
  (1750,17) (1900,18) (2000,20) (2100,19) (2200,22)
};
% Droite de régression y = 0.013376 x - 7.37
\addplot[domain=1100:2300, samples=2, popcurve, thick, poporange] {0.013376*x - 7.37};
% Point moyen G
\addplot[only marks, mark=*, mark size=3, popred] coordinates {(1725,15.7)};
\addplot[only marks, mark=*, mark size=3, popred] coordinates {(1725,15.7)};
\node[popred, above right] at (axis cs:1725,15.7) {$G(\bar{x},\bar{y})$};
\legend{Nuage, Régression $y=ax+b$, Point moyen $G$}
\end{axis}
\end{tikzpicture}
\legende{Nuage de points et droite de régression des moindres carrés. La droite passe par $G$ et suit la tendance du nuage.}
\end{popfigure}

\begin{aretenir}[Interprétation des coefficients]
\begin{itemize}
\item $a = 0,0134$ t/mm : \textbf{Pour 1 mm de pluie supplémentaire, la production augmente en moyenne de 0,0134 tonnes (soit 13,4 kg)}. C'est le \cle{taux de variation moyen} modélisé.
\item $b = -7,37$ t : production théorique pour 0 mm de pluie. \textbf{Sans sens physique ici} (extrapolation hors domaine des données : pluie min 1200 mm). Ne pas interpréter $b$ si $0 \notin [\min(x_i), \max(x_i)]$.
\end{itemize}
\end{aretenir}

\begin{exercice}[Droite de révision $\rightarrow$ Note]
Reprenons les données de l'exercice précédent (heures de révision $X$, note $Y$).
\begin{enumerate}
\item Calculer l'équation de la droite de régression de $Y$ en $X$.
\item Interpréter le coefficient directeur $a$.
\item Prédire la note d'un élève qui révise 5,5 h/semaine.
\end{enumerate}
\end{exercice}

\begin{corrige}[Exercice : Droite révision/note]
Données : $n=6$, $\sum x=27$, $\sum y=71$, $\sum x^2=139$, $\sum xy=348$.
$\bar{x}=4,5$, $\bar{y} \approx 11,83$.
$a = \frac{6\times 348 - 27\times 71}{6\times 139 - 27^2} = \frac{171}{105} \approx 1,629$.
$b = \bar{y} - a\bar{x} = 11,833 - 1,629 \times 4,5 \approx 11,833 - 7,33 = 4,50$.
Droite : $y = 1,63 x + 4,50$.
$a=1,63$ : chaque heure de révision supplémentaire rapporte en moyenne 1,63 points.
Prévision pour $x=5,5$ : $y = 1,63 \times 5,5 + 4,50 \approx 8,97 + 4,50 = 13,47 \approx 13,5/20$. Interpolation (5,5 dans [2;7]), prévision fiable.
\end{corrige}

\newpage
\coursec{Validation du modèle, Prévision et Communication}

\courssub{Coefficient de détermination $R^2$}

\begin{definition}[Coefficient de détermination]
Le \cle{coefficient de détermination} est $R^2 = r^2$. Il représente la \textbf{part de la variance de $Y$ expliquée par la liaison linéaire avec $X$}.
\end{definition}

\begin{aretenir}[Interprétation de $R^2$]
\begin{itemize}
\item $R^2 = 1$ : le modèle explique 100\% de la variabilité de $Y$ (points sur la droite).
\item $R^2 = 0,96$ (ex: $r=0,98$) : le modèle explique 96\% de la variabilité de la production. Les 4\% restants sont dus au hasard, aux erreurs de mesure, ou à d'autres facteurs (qualité du sol, variété, main-d'œuvre...).
\item Au Probatoire : $R^2 > 0,9$ $\rightarrow$ modèle \textbf{très bon} ; $0,7 < R^2 < 0,9$ $\rightarrow$ modèle \textbf{correct} ; $R^2 < 0,7$ $\rightarrow$ modèle \textbf{fragile}, prévisions risquées.
\end{itemize}
\end{aretenir}

\courssub{Résidus et validation graphique}

\begin{definition}[Résidu]
Le \cle{résidu} pour l'individu $i$ est l'écart vertical entre la valeur observée et la valeur théorique (prédite par la droite) :
\[ e_i = y_i - \hat{y}_i = y_i - (ax_i + b) \]
\end{definition}

\begin{methode}[Valider la pertinence du modèle linéaire]
\begin{enumerate}
\item \textbf{Calculer $R^2$} : s'il est faible ($<0,7$), le modèle linéaire est inadapté.
\item \textbf{Analyser les résidus} : tracer le nuage des points $(x_i ; e_i)$.
    \begin{itemize}
    \item Résidus dispersés aléatoirement autour de 0 (forme horizontale) $\rightarrow$ \textbf{Modèle validé}.
    \item Résidus formant une courbe (ex: parabole) $\rightarrow$ \textbf{Liaison non linéaire}, le modèle linéaire est mauvais.
    \item Résidus en éventail (écart croissant avec $x$) $\rightarrow$ \textbf{Hétéroscédasticité}, variance non constante, modèle perfectible.
    \end{itemize}
\item \textbf{Vérifier les valeurs aberrantes} : un point très influent (levier fort) peut tirer la droite. Le retirer et recalculer pour voir la stabilité.
\end{enumerate}
\end{methode}

\begin{exemplebox}[Résidus pour l'exemple Pluie/Production]
Calculons $\hat{y}_i = 0,013376 x_i - 7,37$ et $e_i = y_i - \hat{y}_i$ :
\begin{center}
\begin{tabular}{|c|c|c|c|c|}
\hline
$x_i$ & $y_i$ & $\hat{y}_i$ & $e_i$ \\
\hline
1200 & 8 & 8,68 & -0,68 \\
1350 & 10 & 10,69 & -0,69 \\
1450 & 12 & 12,02 & -0,02 \\
1600 & 15 & 14,03 & +0,97 \\
1700 & 16 & 15,37 & +0,63 \\
1750 & 17 & 16,04 & +0,96 \\
1900 & 18 & 18,04 & -0,04 \\
2000 & 20 & 19,38 & +0,62 \\
2100 & 19 & 20,72 & -1,72 \\
2200 & 22 & 22,06 & -0,06 \\
\hline
\end{tabular}
\end{center}
Somme des résidus $\approx 0$ (propriété). Nuage des résidus : dispersion autour de 0, pas de forme courbe nette $\rightarrow$ modèle linéaire acceptable.
\end{exemplebox}

\courssub{Prévision : Interpolation vs Extrapolation}

\begin{definition}[Prévision]
Pour une valeur $x_0$ de $X$, la valeur prévue de $Y$ est $\hat{y}_0 = a x_0 + b$.
\end{definition}

\begin{attention}[Interpolation vs Extrapolation -- Piège du Bac]
\begin{itemize}
\item \textbf{Interpolation} : $x_0 \in [\min(x_i) ; \max(x_i)]$. Prévision \textbf{fiable} (si $R^2$ élevé).
\item \textbf{Extrapolation} : $x_0 \notin [\min(x_i) ; \max(x_i)]$. Prévision \textbf{risquée, non fiable}. La relation linéaire n'est pas garantie hors de l'intervalle d'observation.
\end{itemize}
\textbf{Au Probatoire :} Préciser toujours si la prévision est une interpolation ou une extrapolation, et conclure sur sa fiabilité.
\end{attention}

\begin{exoresolu}[Prévision pour la campagne prochaine]
Énoncé. Les prévisions météo annoncent 1850 mm de pluie pour l'an prochain. Prédire la production. Et si la météo annonce 3000 mm ?
\tcblower
Solution.
1. $x_0 = 1850$. $1850 \in [1200 ; 2200]$ $\rightarrow$ \textbf{Interpolation}.
$\hat{y} = 0,013376 \times 1850 - 7,37 \approx 24,75 - 7,37 = 17,38$ tonnes.
Prévision fiable (modèle validé, $R^2 \approx 0,96$). On attend environ 17,4 tonnes.

2. $x_0 = 3000$. $3000 \notin [1200 ; 2200]$ $\rightarrow$ \textbf{Extrapolation}.
$\hat{y} = 0,013376 \times 3000 - 7,37 \approx 40,13 - 7,37 = 32,76$ tonnes.
\textbf{Attention :} Cette prévision est \textbf{peu fiable}. Au-delà d'un certain seuil, l'excès d'eau peut pourrir les racines (liaison non linéaire réelle). Le modèle linéaire n'a pas vocation à s'appliquer là.
\end{exoresolu}

\courssub{Communication : Rédiger une conclusion complète}

\begin{methode}[Structure de la réponse attendue au Probatoire]
\begin{enumerate}
\item \textbf{Contexte} : Rappeler les variables $X, Y$, l'effectif $n$, la source des données.
\item \textbf{Analyse descriptive} : Tableau à double entrée, moyennes conditionnelles, nuage de points (forme, direction).
\item \textbf{Mesure de liaison} : Coefficient $r$, valeur, interprétation (force, sens), $R^2$.
\item \textbf{Modélisation} : Équation de la droite de régression $y = ax+b$, interprétation de $a$ (et $b$ si pertinent).
\item \textbf{Validation} : Analyse des résidus (aléatoires ?), $R^2$ suffisant ?
\item \textbf{Prévision} : Calcul pour $x_0$ donné, distinction interpolation/extrapolation, fiabilité.
\item \textbf{Conclusion métier} : Phrase finale en langage courant (ex: « Pour la campagne prochaine, avec 1850 mm annoncés, on prévoit une production d'environ 17,4 tonnes de plantains, ce qui permet de planifier le stockage. »).
\end{enumerate}
\end{methode}

\begin{exercice}[Validation et prévision]
Un électricien mesure la résistance $R$ (en $\Omega$) d'un fil de cuivre en fonction de sa longueur $L$ (en m) :
$(10 ; 0,18)$, $(20 ; 0,35)$, $(30 ; 0,53)$, $(40 ; 0,71)$, $(50 ; 0,88)$.
$r \approx 0,999$, $R^2 \approx 0,998$. Droite : $R = 0,0176 L + 0,004$.
\begin{enumerate}
\item Prédire $R$ pour $L = 35$ m. Type de prévision ? Fiable ?
\item Prédire $R$ pour $L = 100$ m. Type de prévision ? Fiable ?
\item L'ordonnée à l'origine $0,004$ a-t-elle un sens physique ?
\end{enumerate}
\end{exercice}

\begin{corrige}[Exercice : Résistance du cuivre]
1. $L=35 \in [10;50]$ : Interpolation. $R = 0,0176 \times 35 + 0,004 = 0,620 \Omega$. Fiable ($R^2 \approx 1$).
2. $L=100 \notin [10;50]$ : Extrapolation. $R = 0,0176 \times 100 + 0,004 = 1,764 \Omega$. Peu fiable (chauffement du fil, changement de résistivité, hors domaine observé).
3. $L=0$ n'est pas dans l'intervalle. Mais physiquement, un fil de longueur 0 a une résistance 0. Le modèle donne 0,004 $\Omega$ (résistance des contacts/bornes). Sens physique partiel (résistance résiduelle), mais c'est une extrapolation vers 0.
\end{corrige}

\newpage
\coursec{Complément utile : Ajustement non linéaire par linéarisation (Culture Terminale)}

\begin{savaistu}[Quand la liaison n'est pas linéaire ?]
Si le nuage de points a une forme \textbf{courbe} (parabolique, exponentielle, hyperbolique...) et que $r$ est faible (ex: $r=0,6$ mais nuage en courbe), le modèle linéaire $y=ax+b$ est inadapté. On peut parfois \textbf{transformer les variables} pour retrouver une liaison linéaire, ajuster une droite, puis revenir au modèle initial. C'est la \cle{linéarisation}.
\end{savaistu}

\courssub{Cas 1 : Modèle exponentiel $y = a \cdot b^x$ ($b>0$)}
Situation typique : croissance bactérienne, décharge d'un condensateur, amortissement, intérêt composé.
\begin{itemize}
\item On prend le logarithme népérien (ou décimal) de $Y$ : $Y' = \ln(Y)$.
\item Le modèle devient $\ln(y) = \ln(a) + x \ln(b)$, soit $Y' = A + B x$ avec $A = \ln(a)$, $B = \ln(b)$.
\item \textbf{Démarche} :
    1. Calculer $y'_i = \ln(y_i)$ pour chaque point.
    2. Faire la régression linéaire de $Y'$ en $X$ : obtenir $A, B$.
    3. Retrouver $a = e^A$, $b = e^B$.
    4. Modèle final : $y = a \cdot b^x$.
\end{itemize}

\courssub{Cas 2 : Modèle puissance $y = a \cdot x^b$ ($x>0$)}
Situation typique : lois d'échelle, résistance aérodynamique, loi de Kepler ($T^2 \propto a^3$).
\begin{itemize}
\item On prend le logarithme de $X$ et $Y$ : $X' = \ln(X)$, $Y' = \ln(Y)$.
\item Le modèle devient $\ln(y) = \ln(a) + b \ln(x)$, soit $Y' = A + b X'$.
\item \textbf{Démarche} :
    1. Calculer $x'_i = \ln(x_i)$, $y'_i = \ln(y_i)$.
    2. Régression linéaire de $Y'$ en $X'$ : obtenir $A, b$.
    3. $a = e^A$.
    4. Modèle final : $y = a \cdot x^b$.
\end{itemize}

\courssub{Cas 3 : Modèle inverse / hyperbolique $y = \frac{a}{x} + b$}
\begin{itemize}
\item Poser $X' = \frac{1}{X}$. Modèle : $y = a X' + b$.
\item Régression de $Y$ en $X'$.
\end{itemize}

\begin{attention}[Valider après linéarisation]
Le $r$ et $R^2$ calculés portent sur les variables \textbf{transformées} ($Y'$ en $X$ ou $Y'$ en $X'$). Ils valident la liaison linéaire \emph{après transformation}. Il faut \textbf{toujours} tracer le nuage des résidus du modèle final $y = f(x)$ sur les données brutes pour valider l'ajustement réel.
\end{attention}

\begin{exemplebox}[Décharge d'un condensateur (RC série)]
On mesure la tension $u$ (V) aux bornes d'un condensateur qui se décharge dans une résistance, en fonction du temps $t$ (ms).
Données : $(0 ; 10,0)$, $(1 ; 6,1)$, $(2 ; 3,7)$, $(3 ; 2,2)$, $(4 ; 1,4)$, $(5 ; 0,8)$.
Nuage : courbe décroissante convexe. $r$ linéaire faible.
Hypothèse : $u = U_0 e^{-t/\tau}$ (modèle exponentiel).
Linéralisation : $u' = \ln(u)$. Couples $(t ; \ln(u))$ : $(0 ; 2,30)$, $(1 ; 1,81)$, $(2 ; 1,31)$, $(3 ; 0,79)$, $(4 ; 0,34)$, $(5 ; -0,22)$.
Régression de $u'$ en $t$ : $u' = -0,50 t + 2,30$ ($r \approx -0,999$).
Retour : $\ln(U_0) = 2,30 \Rightarrow U_0 = e^{2,30} \approx 10,0$ V. $-1/\tau = -0,50 \Rightarrow \tau = 2,0$ ms.
Modèle final : $\boxed{u = 10,0 \cdot e^{-t/2,0}}$.
Vérification : pour $t=2$, $u = 10 e^{-1} \approx 3,68$ V (mesure 3,7 V). Excellent ajustement.
\end{exemplebox}

\begin{exercice}[Croissance bactérienne -- Linéarisation]
Une culture de bactéries voit sa population $N$ (en millions) évoluer avec le temps $t$ (en h) :
$(0 ; 1,0)$, $(1 ; 1,8)$, $(2 ; 3,2)$, $(3 ; 5,8)$, $(4 ; 10,5)$, $(5 ; 19,0)$.
\begin{enumerate}
\item Tracer le nuage $(t ; N)$. Quelle forme ? Quel modèle envisager ?
\item Linéariser par $N' = \ln(N)$. Faire la régression de $N'$ en $t$.
\item Retrouver le modèle exponentiel $N(t) = N_0 e^{kt}$. Donner $N_0$ et $k$.
\item Prédire la population à $t = 6$ h. Fiable ?
\end{enumerate}
\end{exercice}

\begin{corrige}[Exercice : Bactéries]
1. Nuage : courbe croissante convexe $\rightarrow$ croissance exponentielle. Modèle $N = N_0 e^{kt}$.
2. $N' = \ln(N)$ : $(0 ; 0,00)$, $(1 ; 0,59)$, $(2 ; 1,16)$, $(3 ; 1,76)$, $(4 ; 2,35)$, $(5 ; 2,94)$.
Régression $N' = kt + \ln(N_0)$.
Calculs : $n=6$, $\sum t=15$, $\sum N'=8,80$, $\sum t^2=55$, $\sum tN'=31,85$.
$k = \frac{6\times 31,85 - 15\times 8,80}{6\times 55 - 15^2} = \frac{191,1 - 132}{330 - 225} = \frac{59,1}{105} \approx 0,563$.
$\ln(N_0) = \frac{8,80 - 0,563\times 15}{6} = \frac{8,80 - 8,445}{6} \approx 0,059 \Rightarrow N_0 \approx e^{0,059} \approx 1,06$.
3. Modèle : $\boxed{N(t) = 1,06 \cdot e^{0,563 t}}$.
4. $t=6$ : $N(6) = 1,06 \cdot e^{3,378} \approx 1,06 \times 29,3 \approx 31,1$ millions.
$t=6 \notin [0;5]$ $\rightarrow$ Extrapolation. Fiabilité moyenne (ressources limitées, saturation non modélisée).
\end{corrige}

\begin{aretenir}[Résumé de la boîte à outils Terminale Tech -- Statistiques 2 variables]
\begin{center}
\begin{tabularx}{\linewidth}{|l|X|X|}
\hline
\rowcolor{popblueL} \textbf{Étape} & \textbf{Outil / Formule} & \textbf{Point de vigilance (Bac)} \\
\hline
1. Organiser & Tableau à double entrée ($n_{ij}$, marges) & Triple vérification $\sum n_{ij} = \sum n_{i.} = \sum n_{.j} = n$ \\
\hline
2. Visualiser & Nuage de points $M_i(x_i;y_i)$ & Échelles, pas de traits, lecture forme/sens \\
\hline
3. Mesurer liaison & $r = \frac{\text{cov}}{\sigma_X\sigma_Y}$ & $-1\le r\le 1$, $|r|\approx 1$ fort, $r\approx 0$ nul \textbf{linéaire} \\
\hline
4. Modéliser & Droite régression $y=ax+b$ : $a=\frac{\text{cov}}{\text{Var}(X)}$, $b=\bar{y}-a\bar{x}$ & Passe par $G(\bar{x},\bar{y})$. $a$ = variation moyenne de $Y$ pour 1 unité de $X$. \\
\hline
5. Valider & $R^2=r^2$ (part expliquée), Résidus $e_i=y_i-\hat{y}_i$ & $R^2>0,9$ bon. Résidus aléatoires $\checkmark$. Résidus courbes $\times$. \\
\hline
6. Prévoir & $\hat{y}_0 = ax_0+b$ & \textbf{Interpolation} (fiable) vs \textbf{Extrapolation} (risquée). Le dire explicitement. \\
\hline
7. Communiquer & Structure en 7 étapes (Méthode p. 4) & Phrase de conclusion métier obligatoire. \\
\hline
\end{tabularx}
\end{center}
\end{aretenir}

\coursec{Visualisation et mesure de la liaison : Nuage de points \& Coefficient de corrélation}

Dans la section précédente, nous avons appris à \cle{organiser} les données d'une série statistique à deux variables dans un \cle{tableau à double entrée} : effectifs, marges, profils. Ce tableau résume l'information, mais il ne permet pas encore de « voir » la relation entre les deux caractères. Un chef d'atelier qui note chaque semaine le nombre d'heures de fonctionnement d'une machine $X$ et le nombre de pannes $Y$ veut savoir : \emph{les deux grandeurs évoluent-elles ensemble ? Si oui, dans quel sens, et avec quelle force ?} C'est exactement l'objet de cette section : passer du tableau au \cle{graphique} (le nuage de points), puis du graphique à un \cle{nombre} qui mesure la liaison (le coefficient de corrélation).

\courssub{Le nuage de points : voir la liaison}

\begin{definition}[Nuage de points]
Soit une série statistique double $(X, Y)$ de $n$ couples $(x_i, y_i)$. Le \cle{nuage de points} associé est l'ensemble des points $M_i$ de coordonnées $(x_i\,;\,y_i)$ placés dans un repère du plan. On choisit des échelles adaptées aux valeurs observées (le repère n'a pas besoin d'être orthonormé : l'axe des abscisses porte $X$, l'axe des ordonnées porte $Y$).
\end{definition}

L'intuition est simple : chaque individu statistique (une semaine, un client, une parcelle agricole) devient \textbf{un point}. Si les points « montent » globalement de la gauche vers la droite, $Y$ augmente quand $X$ augmente : la liaison est \cle{positive} (croissante). S'ils « descendent », la liaison est \cle{négative} (décroissante). Si les points sont étalés sans direction privilégiée, il n'y a pas de liaison linéaire visible.

\textbf{Lecture visuelle d'un nuage.} On examine trois choses :
\begin{itemize}
\item la \cle{forme} : un nuage allongé en ellipse (en « cigare ») suggère une liaison linéaire ; un nuage en arc de parabole suggère une liaison non linéaire ;
\item la \cle{dispersion} : plus l'ellipse est fine, plus la liaison est forte ; plus elle est épaisse ou ronde, plus elle est faible ;
\item les \cle{valeurs aberrantes} (\emph{outliers}) : un point isolé, très éloigné des autres, peut fausser tous les calculs. Il faut le repérer, vérifier qu'il ne s'agit pas d'une erreur de saisie, et éventuellement l'écarter en le justifiant.
\end{itemize}

\begin{popfigure}
\begin{center}
\begin{tikzpicture}[scale=0.9]
\begin{axis}[width=0.36\linewidth,height=0.32\linewidth,title={$r \approx 0{,}9$ : fort positif},title style={font=\small},xmin=0,xmax=10,ymin=0,ymax=10,xtick=\empty,ytick=\empty,axis lines=left]
\addplot[only marks,mark=*,mark size=1.2pt,popblue] coordinates {(1,1.5) (2,2.4) (3,2.8) (4,4.1) (5,4.6) (6,5.8) (7,6.5) (8,7.6) (9,8.4)};
\end{axis}
\end{tikzpicture}\hfill
\begin{tikzpicture}[scale=0.9]
\begin{axis}[width=0.36\linewidth,height=0.32\linewidth,title={$r \approx -0{,}3$ : faible négatif},title style={font=\small},xmin=0,xmax=10,ymin=0,ymax=10,xtick=\empty,ytick=\empty,axis lines=left]
\addplot[only marks,mark=*,mark size=1.2pt,poporange] coordinates {(1,7.5) (2,4.2) (3,8.1) (4,3.5) (5,6.4) (6,2.8) (7,5.5) (8,3.9) (9,4.6)};
\end{axis}
\end{tikzpicture}\hfill
\begin{tikzpicture}[scale=0.9]
\begin{axis}[width=0.36\linewidth,height=0.32\linewidth,title={$r \approx 0$ : liaison non linéaire},title style={font=\small},xmin=-5,xmax=5,ymin=-1,ymax=10,xtick=\empty,ytick=\empty,axis lines=left]
\addplot[only marks,mark=*,mark size=1.2pt,poppurple] coordinates {(-4,8.2) (-3,4.6) (-2,2.1) (-1,0.6) (0,0.2) (1,0.7) (2,2.3) (3,4.9) (4,8.5)};
\end{axis}
\end{tikzpicture}
\end{center}
\legende{Trois nuages de points. À gauche, une liaison linéaire forte et positive ; au centre, une liaison faible ; à droite, $r \approx 0$ alors qu'existe une liaison parabolique évidente : $r$ ne mesure que la liaison \emph{linéaire}.}
\end{popfigure}

\courssub{La covariance : un premier outil de mesure}

Le graphique donne une impression ; les mathématiques exigent une mesure. Pour savoir si $X$ et $Y$ varient « dans le même sens », l'idée naturelle est de comparer, pour chaque individu, l'écart de $x_i$ à la moyenne $\bar{x}$ et l'écart de $y_i$ à la moyenne $\bar{y}$. Si les deux écarts sont souvent de même signe, le produit $(x_i-\bar{x})(y_i-\bar{y})$ est souvent positif : les grandeurs évoluent ensemble.

\begin{definition}[Covariance]
Soit une série double $(X, Y)$ de $n$ couples $(x_i, y_i)$, de moyennes $\bar{x}$ et $\bar{y}$. La \cle{covariance} de $X$ et $Y$ est le nombre :
\[ \operatorname{cov}(X,Y) = \frac{1}{n}\sum_{i=1}^{n}(x_i - \bar{x})(y_i - \bar{y}) = \frac{1}{n}\sum_{i=1}^{n} x_i y_i - \bar{x}\,\bar{y}. \]
La seconde formule (dite \emph{formule de König-Huygens}) est celle que l'on utilise en pratique : on calcule la moyenne des produits, puis on retire le produit des moyennes.
\end{definition}

\begin{propriete}[Interprétation du signe de la covariance]
\begin{itemize}
\item Si $\operatorname{cov}(X,Y) > 0$ : $X$ et $Y$ ont tendance à varier dans le \textbf{même sens} (liaison positive).
\item Si $\operatorname{cov}(X,Y) < 0$ : $X$ et $Y$ ont tendance à varier en \textbf{sens contraires} (liaison négative).
\item Si $\operatorname{cov}(X,Y) \approx 0$ : pas de tendance linéaire décelable.
\end{itemize}
La covariance dépend des unités choisies (cm, FCFA, heures...) : sa \emph{valeur} n'est donc pas interprétable directement, seul son \emph{signe} l'est. C'est pourquoi on lui préfère le coefficient de corrélation, sans unité.
\end{propriete}

\begin{popfigure}
\begin{center}
\begin{tikzpicture}[scale=0.85]
\draw[->,thick] (-0.3,0) -- (8.5,0) node[below left]{$x$};
\draw[->,thick] (0,-0.3) -- (0,6.5) node[below left]{$y$};
\draw[dashed,popmuted] (4.5,0) -- (4.5,6.2) node[above]{$x=\bar{x}$};
\draw[dashed,popmuted] (0,3.2) -- (8.2,3.2) node[right]{$y=\bar{y}$};
\fill[poppurple] (4.5,3.2) circle (2.5pt) node[below right=1pt]{$G(\bar{x}\,;\,\bar{y})$};
% point en haut à droite : rectangle positif
\fill[popgreen!35] (4.5,3.2) rectangle (7,5.2);
\draw[popgreen,thick] (4.5,3.2) rectangle (7,5.2);
\fill[popdark] (7,5.2) circle (2pt) node[above right]{$M_i$};
\node[popdark,font=\small] at (5.75,4.2) {$+$};
% point en bas à gauche : rectangle positif
\fill[popgreen!35] (4.5,3.2) rectangle (1.5,1.2);
\draw[popgreen,thick] (4.5,3.2) rectangle (1.5,1.2);
\fill[popdark] (1.5,1.2) circle (2pt) node[below left]{$M_j$};
\node[popdark,font=\small] at (3,2.2) {$+$};
% point en haut à gauche : rectangle négatif
\fill[poporange!40] (4.5,3.2) rectangle (2,5.5);
\draw[poporange,thick] (4.5,3.2) rectangle (2,5.5);
\fill[popdark] (2,5.5) circle (2pt) node[above left]{$M_k$};
\node[popdark,font=\small] at (3.25,4.35) {$-$};
\end{tikzpicture}
\end{center}
\legende{Lecture géométrique de la covariance. Le point moyen $G$ partage le plan en quatre zones. Pour chaque point $M_i$, le produit $(x_i-\bar{x})(y_i-\bar{y})$ est l'aire \emph{signée} du rectangle construit sur $G$ et $M_i$ : positive dans les zones « haut-droite » et « bas-gauche », négative dans les deux autres. La covariance est la moyenne de ces aires signées.}
\end{popfigure}

\courssub{Le coefficient de corrélation linéaire de Pearson}

Pour supprimer l'effet des unités, on \cle{normalise} la covariance en la divisant par le produit des écarts types.

\begin{definition}[Coefficient de corrélation linéaire]
Soit une série double $(X, Y)$ avec $\sigma_X > 0$ et $\sigma_Y > 0$. Le \cle{coefficient de corrélation linéaire} (ou coefficient de Pearson) est :
\[ r = \frac{\operatorname{cov}(X,Y)}{\sigma_X \, \sigma_Y}. \]
C'est un nombre \textbf{sans unité}.
\end{definition}

\begin{propriete}[Encadrement de $r$ (admise)]
Une conséquence de l'inégalité de Cauchy-Schwarz (admise à ce niveau) garantit que :
\[ -1 \leq r \leq 1. \]
De plus :
\begin{itemize}
\item $r = 1$ si et seulement si tous les points sont alignés sur une droite \textbf{croissante} ;
\item $r = -1$ si et seulement si tous les points sont alignés sur une droite \textbf{décroissante}.
\end{itemize}
\end{propriete}

\begin{aretenir}[Interprétation de $r$]
\begin{itemize}
\item Le \textbf{signe} de $r$ donne le \textbf{sens} de la liaison : $r>0$ liaison croissante, $r<0$ liaison décroissante.
\item La \textbf{valeur absolue} donne la \textbf{force} : $|r|$ proche de $1$ $\Rightarrow$ liaison linéaire forte (nuage resserré autour d'une droite) ; $|r|$ proche de $0$ $\Rightarrow$ liaison linéaire faible ou inexistante.
\item En pratique, on considère qu'un ajustement linéaire est pertinent lorsque $|r| \geq 0{,}87$ environ (soit $|r|$ proche de $1$).
\end{itemize}
\end{aretenir}

\begin{definition}[Coefficient de détermination]
Le \cle{coefficient de détermination} est $R^2 = r^2$. Il s'interprète comme la \textbf{part de la variance de $Y$ expliquée par la liaison linéaire avec $X$}, exprimée en pourcentage. Par exemple, $R^2 = 0{,}76$ signifie que $76\,\%$ de la variabilité de $Y$ s'explique par sa relation linéaire avec $X$ (les $24\,\%$ restants viennent d'autres facteurs).
\end{definition}

\courssub{Exemple chiffré complet : pluie et production agricole}

\begin{exoresolu}[Calcul manuel de la covariance et de $r$]
\textbf{Énoncé.} Un exploitant agricole de la région de l'Ouest a relevé, pendant 5 campagnes, la pluviométrie $x_i$ (en dizaines de mm) et la production de maïs $y_i$ (en tonnes) :
\begin{center}
\begin{tabular}{|l|ccccc|}
\hline
\rowcolor{popblueL} Campagne $i$ & 1 & 2 & 3 & 4 & 5 \\
\hline Pluie $x_i$ & 8 & 10 & 12 & 14 & 16 \\
\hline Production $y_i$ & 4 & 5 & 7 & 8 & 9 \\
\hline
\end{tabular}
\end{center}
Calculer la covariance, le coefficient de corrélation $r$, le coefficient $R^2$, puis interpréter.
\tcblower
\textbf{Solution détaillée.}

\textbf{Étape 1 : moyennes.}
\[ \bar{x} = \frac{8+10+12+14+16}{5} = \frac{60}{5} = 12, \qquad \bar{y} = \frac{4+5+7+8+9}{5} = \frac{33}{5} = 6{,}6. \]

\textbf{Étape 2 : covariance (formule de König-Huygens).} On calcule les produits $x_i y_i$ :
\[ \sum x_i y_i = 8\times 4 + 10\times 5 + 12\times 7 + 14\times 8 + 16\times 9 = 32+50+84+112+144 = 422. \]
\[ \operatorname{cov}(X,Y) = \frac{422}{5} - 12 \times 6{,}6 = 84{,}4 - 79{,}2 = 5{,}2. \]
La covariance est \textbf{positive} : plus il pleut, plus la production est élevée.

\textbf{Étape 3 : écarts types.}
\begin{align*}
\sigma_X^2 &= \frac{8^2+10^2+12^2+14^2+16^2}{5} - 12^2 = \frac{64+100+144+196+256}{5} - 144 = \frac{760}{5} - 144 = 152 - 144 = 8,\\
\sigma_Y^2 &= \frac{4^2+5^2+7^2+8^2+9^2}{5} - 6{,}6^2 = \frac{16+25+49+64+81}{5} - 43{,}56 = \frac{235}{5} - 43{,}56 = 47 - 43{,}56 = 3{,}44.
\end{align*}
Donc $\sigma_X = \sqrt{8} \approx 2{,}83$ et $\sigma_Y = \sqrt{3{,}44} \approx 1{,}85$.

\textbf{Étape 4 : coefficient de corrélation.}
\[ r = \frac{5{,}2}{2{,}83 \times 1{,}85} \approx \frac{5{,}2}{5{,}24} \approx 0{,}99. \]

\textbf{Étape 5 : interprétation.} $r \approx 0{,}99$ : la liaison linéaire est \textbf{forte et positive}. Le coefficient de détermination vaut $R^2 = r^2 \approx 0{,}98$ : environ $98\,\%$ de la variabilité de la production s'explique par la pluviométrie. Un ajustement linéaire (section suivante) est donc pleinement justifié.
\end{exoresolu}

\courssub{Obtention de $r$ à la calculatrice}

Au Probatoire et au Baccalauréat technique, le calcul à la main est exigible sur de petites séries, mais la calculatrice permet de \textbf{vérifier} et de traiter rapidement de grandes séries.

\begin{methode}[Calculer $r$, $\bar{x}$, $\bar{y}$, $\sigma_X$, $\sigma_Y$ à la calculatrice]
\textbf{Sur Casio (Graph 35/90) :}
\begin{enumerate}
\item Menu \textbf{STAT} : saisir les $x_i$ dans la List 1 et les $y_i$ dans la List 2.
\item Touche \textbf{CALC} (F2), puis \textbf{SET} (F6) : vérifier 2Var X : List1, 2Var Y : List2.
\item Choisir \textbf{REG} (F3) puis \textbf{X} (F1), mode $a+bx$.
\item Lire à l'écran : $r$, mais aussi $a$, $b$, et via \textbf{2VAR} : $\bar{x}$, $\bar{y}$, $\sigma_X$ ($x\sigma n$), $\sigma_Y$, $\sum xy$.
\end{enumerate}
\textbf{Sur NumWorks :}
\begin{enumerate}
\item Application \textbf{Régressions} : entrer les valeurs dans les colonnes X1 et Y1.
\item Onglet \textbf{Graph} : le nuage s'affiche ; onglet \textbf{Stats} : lire $\bar{x}$, $\bar{y}$, $\sigma_X$, $\sigma_Y$, $r$ et $r^2$ après avoir choisi le modèle \textbf{Linéaire} $ax+b$.
\end{enumerate}
\end{methode}

\begin{attention}[Erreurs fréquentes]
\begin{itemize}
\item \textbf{« $r = 0$ donc il n'y a pas de liaison » : FAUX.} $r \approx 0$ signifie seulement « pas de liaison \emph{linéaire} ». La série $Y = X^2$ sur $[-4\,;\,4]$ donne $r \approx 0$ alors que $Y$ dépend parfaitement de $X$ (voir le troisième nuage de la figure). Toujours \textbf{regarder le nuage} avant de conclure.
\item \textbf{Confondre corrélation et causalité.} $r \approx 0{,}99$ entre pluie et production ne prouve pas que la pluie \emph{seule} cause la production (engrais, variétés, main-d'œuvre interviennent), et surtout pas que la production provoque la pluie ! Une corrélation peut aussi venir d'une troisième variable cachée.
\item \textbf{Oublier les valeurs aberrantes.} Un seul point aberrant peut faire passer $r$ de $0{,}9$ à $0{,}4$. Toujours vérifier la saisie des données.
\item \textbf{Diviser par $\sigma_X \sigma_Y$ dans le mauvais sens} ou oublier que $r$ est sans unité alors que la covariance en a une.
\end{itemize}
\end{attention}

\begin{savaistu}[Francis Galton et la « régression vers la moyenne »]
C'est Francis Galton (1822--1911), cousin de Charles Darwin, qui a introduit ces notions. En étudiant la taille des parents et celle de leurs enfants, il observa un phénomène étonnant : les enfants de parents très grands sont grands, mais \emph{en moyenne moins grands que leurs parents} ; symétriquement pour les parents très petits. Il appela cela la « régression vers la médiocrité » — d'où le mot \cle{régression} encore utilisé aujourd'hui. Son élève Karl Pearson formalisa ensuite le coefficient $r$ qui porte son nom. Ces outils sont omniprésents : météorologie, contrôle qualité en usine, études de marché, médecine... et même l'agriculture camerounaise, où l'on corrèle pluviométrie et rendements pour planifier les campagnes.
\end{savaistu}

\begin{exercice}[À vous de jouer]
Un garage note pendant 6 mois le nombre $x_i$ de véhicules révisés et le chiffre d'affaires $y_i$ (en centaines de milliers de FCFA) : $(10\,;\,15)$, $(12\,;\,17)$, $(15\,;\,22)$, $(18\,;\,25)$, $(20\,;\,28)$, $(25\,;\,33)$.
\begin{enumerate}
\item Tracer le nuage de points et décrire visuellement la liaison.
\item Calculer $\bar{x}$, $\bar{y}$, $\operatorname{cov}(X,Y)$ puis $r$ (on donne $\sigma_X \approx 5{,}02$ et $\sigma_Y \approx 6{,}18$).
\item Calculer $R^2$ et conclure sur la pertinence d'un ajustement linéaire.
\end{enumerate}
\end{exercice}

\begin{corrige}[Exercice : garage]
\begin{enumerate}
\item Le nuage est allongé, montant de gauche à droite : liaison linéaire forte et positive.
\item $\bar{x} = \dfrac{10+12+15+18+20+25}{6} = \dfrac{100}{6} \approx 16{,}67$ ; $\bar{y} = \dfrac{15+17+22+25+28+33}{6} = \dfrac{140}{6} \approx 23{,}33$.\\
$\sum x_i y_i = 150+204+330+450+560+825 = 2519$, d'où $\operatorname{cov}(X,Y) = \dfrac{2519}{6} - 16{,}67 \times 23{,}33 \approx 419{,}83 - 388{,}89 = 30{,}94$.\\
$r = \dfrac{30{,}94}{5{,}02 \times 6{,}18} \approx \dfrac{30{,}94}{31{,}02} \approx 1{,}00$ (plus précisément $r \approx 0{,}997$).
\item $R^2 \approx 0{,}99$ : environ $99\,\%$ de la variabilité du chiffre d'affaires s'explique par le nombre de révisions. Comme $|r| \geq 0{,}87$, un ajustement linéaire est tout à fait pertinent.
\end{enumerate}
\end{corrige}

Nous savons désormais \textbf{voir} la liaison (nuage de points) et la \textbf{mesurer} (covariance, $r$, $R^2$). Reste à la \textbf{modéliser} : c'est l'objet de la section suivante, où nous tracerons la droite qui passe « au plus près » de tous les points — la droite de régression des moindres carrés — pour enfin faire des prévisions.

\coursec{Modélisation : Droite de régression des moindres carrés (y en x)}

Dans la section précédente, nous avons appris à \emph{visualiser} la liaison entre deux caractères statistiques grâce au nuage de points, puis à la \emph{mesurer} grâce au coefficient de corrélation linéaire $r$. Nous savons donc désormais répondre à la question : « Y a-t-il une liaison linéaire, et est-elle forte ? ». Mais un technicien, un agronome ou un gestionnaire veut aller plus loin : il veut \emph{prévoir}. Si la pluviométrie annoncée est de \SI{1200}{mm}, quelle production de maïs peut-on espérer ? Si l'atelier tourne 40 heures, quel sera le coût de maintenance ? Pour répondre, il faut remplacer le nuage de points par un \cle{modèle} simple : une droite. Toute cette section est consacrée à la construction rigoureuse de cette droite, appelée \cle{droite de régression} ou \cle{droite des moindres carrés}.

\courssub{Le problème : quelle droite « colle » le mieux au nuage ?}

\begin{definition}[Résidu d'un point par rapport à une droite]
Soit un nuage de $n$ points $M_i(x_i\,;\,y_i)$ et une droite $(D)$ d'équation $y = ax + b$. Pour chaque valeur $x_i$, la droite « prédit » la valeur $\widehat{y_i} = ax_i + b$, alors que la valeur réellement observée est $y_i$. On appelle \cle{résidu} (ou écart vertical) du point $M_i$ le nombre :
\[ e_i = y_i - (a x_i + b) = y_i - \widehat{y_i}. \]
Le résidu est positif si le point est au-dessus de la droite, négatif s'il est en dessous.
\end{definition}

L'idée naturelle est de chercher la droite qui rend les résidus « globalement petits ». Mais on ne peut pas simplement additionner les $e_i$ : un résidu de $+5$ et un résidu de $-5$ se compenseraient alors que la droite est mauvaise. On élève donc chaque résidu au carré (comme pour la variance !), ce qui supprime les signes et pénalise davantage les gros écarts.

\begin{definition}[Critère des moindres carrés]
La \cle{droite des moindres carrés} (ou droite de régression de $y$ en $x$) est la droite $y = ax + b$ qui rend minimale la somme des carrés des résidus :
\[ S(a,b) = \sum_{i=1}^{n} \bigl(y_i - a x_i - b\bigr)^2. \]
\end{definition}

\begin{attention}[Pourquoi des carrés ?]
Deux raisons à retenir : (1) le carré rend tous les écarts positifs, donc ils ne s'annulent pas entre eux ; (2) le carré « punit » fortement les grands écarts ($4^2 = 16$ contre $2^2 = 4$), ce qui force la droite à ne s'éloigner trop d'aucun point. Attention à ne pas confondre avec la somme des écarts simples, qui vaut toujours $0$ pour la droite des moindres carrés et ne permet donc pas de la trouver.
\end{attention}

\begin{popfigure}
\begin{center}
\begin{tikzpicture}
\begin{axis}[
  width=0.92\linewidth, height=7cm,
  xlabel={Pluie $x$ (mm)}, ylabel={Production $y$ (t)},
  xmin=500, xmax=2100, ymin=0, ymax=8,
  grid=both, grid style={popline!40},
  legend style={at={(0.02,0.98)}, anchor=north west, font=\small}
]
\addplot[popcurve, domain=500:2100, samples=50] {0.0045*x - 1.95};
\addlegendentry{Droite : $y = 0{,}0045x - 1{,}95$}
\addplot[only marks, mark=*, mark size=2.5pt, poporange] coordinates {
  (800,1.5) (1000,2.8) (1200,3.0) (1400,4.5) (1600,5.4) (1800,6.2)
};
\addlegendentry{Points observés $M_i$}
% Point moyen G
\addplot[only marks, mark=star, mark size=4pt, popgreen] coordinates {(1300,3.9)};
\addlegendentry{Point moyen $G(\bar{x}\,;\,\bar{y})$}
% Résidus verticaux pour 3 points
\draw[dashed, poppink, thick] (axis cs:1000,2.8) -- (axis cs:1000,2.55);
\draw[dashed, poppink, thick] (axis cs:1400,4.5) -- (axis cs:1400,4.35);
\draw[dashed, poppink, thick] (axis cs:1800,6.2) -- (axis cs:1800,6.15);
\node[poppink, font=\small] at (axis cs:1010,2.65) [right] {$e_2$};
\node[poppink, font=\small] at (axis cs:1410,4.42) [right] {$e_4$};
\node[poppink, font=\small] at (axis cs:1810,6.17) [right] {$e_6$};
\end{axis}
\end{tikzpicture}
\end{center}
\legende{Nuage de points, droite des moindres carrés, point moyen $G$ et trois résidus verticaux $e_i$ (pointillés). La droite minimise la somme des carrés de ces segments.}
\end{popfigure}

\courssub{Recherche du minimum : le système normal}

\begin{experience}[Démonstration guidée : trouver $a$ et $b$]
On cherche le couple $(a\,;\,b)$ qui minimise $S(a,b) = \sum_{i=1}^n (y_i - a x_i - b)^2$. Suivons les étapes. (Cette démonstration est donnée pour la culture et la compréhension ; elle n'est pas exigible au Baccalauréat technique.)

\textbf{Étape 1 : une condition nécessaire.} Pour qu'une quantité dépendant de deux nombres $a$ et $b$ soit minimale, il est nécessaire que ses « taux de variation » par rapport à $a$ et à $b$ s'annulent (idée admise : au fond d'une « cuvette », la pente est nulle dans toutes les directions). On calcule donc les dérivées de $S$ par rapport à $a$ (en considérant $b$ fixe) et par rapport à $b$ (en considérant $a$ fixe).

\textbf{Étape 2 : dériver par rapport à $b$.} En dérivant chaque carré (forme $u^2$, de dérivée $2u\,u'$) :
\[ \frac{\partial S}{\partial b} = \sum_{i=1}^n 2\bigl(y_i - a x_i - b\bigr)(-1) = -2\sum_{i=1}^n \bigl(y_i - a x_i - b\bigr). \]
L'annulation donne : $\sum y_i - a\sum x_i - n b = 0$, c'est-à-dire la \textbf{première équation normale} :
\[ \sum_{i=1}^n y_i = a\sum_{i=1}^n x_i + n\,b \]
En divisant par $n$ : $\bar{y} = a\bar{x} + b$. \emph{La droite passe donc par le point moyen $G(\bar{x}\,;\,\bar{y})$ !}

\textbf{Étape 3 : dériver par rapport à $a$.}
\[ \frac{\partial S}{\partial a} = \sum_{i=1}^n 2\bigl(y_i - a x_i - b\bigr)(-x_i) = -2\sum_{i=1}^n x_i\bigl(y_i - a x_i - b\bigr). \]
L'annulation donne la \textbf{deuxième équation normale} :
\[ \sum_{i=1}^n x_i y_i = a\sum_{i=1}^n x_i^2 + b\sum_{i=1}^n x_i \]

\textbf{Étape 4 : résoudre le système.} On remplace $b = \bar{y} - a\bar{x}$ (étape 2) dans la deuxième équation. En développant et en regroupant les termes en $a$ :
\[ a\left(\sum x_i^2 - n\bar{x}^2\right) = \sum x_i y_i - n\bar{x}\bar{y}. \]
Or on reconnaît les formules de König : $\sum x_i^2 - n\bar{x}^2 = n\,\mathrm{Var}(X)$ et $\sum x_i y_i - n\bar{x}\bar{y} = n\,\mathrm{cov}(X,Y)$. D'où :
\[ a = \frac{\mathrm{cov}(X,Y)}{\mathrm{Var}(X)}. \]
Le dénominateur $\mathrm{Var}(X)$ est strictement positif (les $x_i$ ne sont pas tous égaux), donc la solution existe et est unique : c'est bien un minimum car $S$ est une somme de carrés, quantité toujours positive qui « explose » dès que la droite s'éloigne du nuage.
\end{experience}

\begin{propriete}[Théorème : droite des moindres carrés]
Soit une série statistique double $(x_i\,;\,y_i)$ de $n$ points, avec $\mathrm{Var}(X) \neq 0$. La droite qui minimise $S(a,b) = \sum (y_i - a x_i - b)^2$ est la \cle{droite de régression de $y$ en $x$}, d'équation $y = ax + b$ avec :
\[ a = \frac{\mathrm{cov}(X,Y)}{\mathrm{Var}(X)} = r\,\frac{\sigma_Y}{\sigma_X} \qquad \text{et} \qquad b = \bar{y} - a\,\bar{x}. \]
\textbf{Propriété fondamentale :} cette droite passe par le \cle{point moyen} $G(\bar{x}\,;\,\bar{y})$. Elle admet donc aussi l'\cle{équation réduite} :
\[ y - \bar{y} = a\,(x - \bar{x}). \]
(Au Baccalauréat technique, on demande surtout de \emph{savoir utiliser} les formules et de \emph{justifier} le passage par $G$ ; l'égalité $a = r\,\sigma_Y/\sigma_X$ se déduit de $r = \dfrac{\mathrm{cov}(X,Y)}{\sigma_X \sigma_Y}$.)
\end{propriete}

\begin{aretenir}[Interprétation du coefficient directeur $a$]
Le nombre $a$ s'interprète concrètement : \textbf{quand $X$ augmente d'une unité, $Y$ varie en moyenne de $a$ unités}. Le signe de $a$ est celui de $r$ (car $\sigma_Y/\sigma_X > 0$) : $a > 0$ pour une liaison positive, $a < 0$ pour une liaison négative. L'équation réduite $y - \bar{y} = a(x - \bar{x})$ dit la même chose autrement : l'écart de $y$ à sa moyenne est proportionnel à l'écart de $x$ à la sienne.
\end{aretenir}

\begin{popfigure}
\begin{center}
\begin{tikzpicture}[scale=1.05, >=Stealth]
% Axes
\draw[->, popdark] (0,0) -- (7.2,0) node[below left]{$x$};
\draw[->, popdark] (0,0) -- (0,4.6) node[below left]{$y$};
% Droite de régression
\draw[very thick, popblue] (0.4,0.75) -- (6.8,3.95) node[above right, font=\small]{$(D) : y = ax + b$};
% Point moyen G
\coordinate (G) at (3.6,2.35);
\fill[popgreen] (G) circle (2.5pt) node[above right, popdark]{$G(\bar{x}\,;\,\bar{y})$};
% Projections de G
\draw[dashed, popmuted] (G) -- (3.6,0) node[below]{$\bar{x}$};
\draw[dashed, popmuted] (G) -- (0,2.35) node[left]{$\bar{y}$};
% Barycentre : points du nuage répartis autour
\fill[poporange] (1.2,1.5) circle (2pt);
\fill[poporange] (2.0,0.9) circle (2pt);
\fill[poporange] (2.8,2.1) circle (2pt);
\fill[poporange] (4.4,2.7) circle (2pt);
\fill[poporange] (5.2,3.4) circle (2pt);
\fill[poporange] (6.0,3.1) circle (2pt);
% Vecteurs vers G (idée barycentre)
\draw[->, poppink, thin] (1.2,1.5) -- (G);
\draw[->, poppink, thin] (6.0,3.1) -- (G);
\node[poppink, font=\small, align=center] at (3.6,4.3) {Les écarts des points autour de $G$\\ se compensent : $G$ est le barycentre du nuage};
\end{tikzpicture}
\end{center}
\legende{Idée géométrique du passage par $G$ : le point moyen est le barycentre (centre de gravité) du nuage de points. La première équation normale $\bar{y} = a\bar{x} + b$ impose à la droite de passer par ce centre d'équilibre.}
\end{popfigure}

\courssub{Méthode pratique et exemple chiffré}

\begin{methode}[Déterminer l'équation de la droite de régression de $y$ en $x$]
\begin{enumerate}
  \item À la calculatrice (menu \textbf{STAT} puis \textbf{REG LIN}), saisir les listes $(x_i)$ et $(y_i)$ et lire : $\bar{x}$, $\bar{y}$, $\sigma_X$, $\sigma_Y$ et $r$.
  \item Calculer le coefficient directeur : $a = r \times \dfrac{\sigma_Y}{\sigma_X}$ (ou directement $a = \dfrac{\mathrm{cov}(X,Y)}{\mathrm{Var}(X)}$).
  \item Calculer l'ordonnée à l'origine : $b = \bar{y} - a\bar{x}$.
  \item Écrire l'équation : $y = ax + b$, en remplaçant $a$ et $b$ par leurs valeurs arrondies avec soin.
  \item \textbf{Vérifier} que la droite passe par $G$ : calculer $a\bar{x} + b$ et contrôler qu'on retrouve bien $\bar{y}$.
\end{enumerate}
\end{methode}

\begin{exemplebox}[Pluie et production agricole]
Reprenons les données des sections précédentes : pour six parcelles, on a relevé la pluviométrie annuelle $x$ (en mm) et la production de maïs $y$ (en tonnes). La calculatrice a donné :
\[ \bar{x} = 1300 \text{ mm}, \quad \bar{y} = 3{,}9 \text{ t}, \quad \sigma_X \approx 346{,}4, \quad \sigma_Y \approx 1{,}56, \quad r \approx 0{,}99. \]
\textbf{Étape 2 :} $a = r\,\dfrac{\sigma_Y}{\sigma_X} \approx 0{,}99 \times \dfrac{1{,}56}{346{,}4} \approx 0{,}0045$ t/mm.
\textbf{Étape 3 :} $b = \bar{y} - a\bar{x} \approx 3{,}9 - 0{,}0045 \times 1300 \approx 3{,}9 - 5{,}85 = -1{,}95$ t (arrondi).
\textbf{Étape 4 :} l'équation du modèle s'écrit :
\[ \text{Production} = 0{,}0045 \times \text{Pluie} - 1{,}95. \]
\textbf{Étape 5 (vérification) :} $0{,}0045 \times 1300 - 1{,}95 = 5{,}85 - 1{,}95 = 3{,}9 = \bar{y}$. La droite passe bien par $G(1300\,;\,3{,}9)$.
\textbf{Interprétation :} $a = 0{,}0045$ t/mm signifie que lorsque la pluviométrie augmente de \SI{100}{mm}, la production augmente \emph{en moyenne} de $100 \times 0{,}0045 = 0{,}45$ tonne. L'ordonnée à l'origine $b = -1{,}95$ n'a pas de sens agricole direct (une production négative !) : c'est un paramètre mathématique du modèle, valable uniquement dans la plage des pluies observées.
\end{exemplebox}

\begin{exoresolu}[Atelier : heures de machine et coût de maintenance]
Un atelier mécanique de Douala relève sur 5 mois le nombre d'heures de fonctionnement $x$ d'une machine et le coût de maintenance $y$ (en milliers de FCFA) :
\begin{center}
\begin{tabular}{|l|ccccc|}\hline
\rowcolor{popblueL} Mois & 1 & 2 & 3 & 4 & 5 \\ \hline
$x$ (heures) & 100 & 120 & 140 & 160 & 180 \\ \hline
$y$ (milliers FCFA) & 22 & 26 & 27 & 32 & 33 \\ \hline
\end{tabular}
\end{center}
On donne : $\bar{x} = 140$, $\bar{y} = 28$, $\sum x_i y_i = \num{20160}$, $\sum x_i^2 = \num{102000}$.
Déterminer l'équation de la droite de régression de $y$ en $x$, puis prévoir le coût pour 200 heures.
\tcblower
\textbf{Solution détaillée.}
\textbf{1. Covariance :} $\mathrm{cov}(X,Y) = \dfrac{1}{n}\sum x_i y_i - \bar{x}\bar{y} = \dfrac{\num{20160}}{5} - 140 \times 28 = 4032 - 3920 = 112$.
\textbf{2. Variance de $X$ :} $\mathrm{Var}(X) = \dfrac{1}{n}\sum x_i^2 - \bar{x}^2 = \dfrac{\num{102000}}{5} - 140^2 = \num{20400} - \num{19600} = 800$.
\textbf{3. Coefficient directeur :} $a = \dfrac{\mathrm{cov}(X,Y)}{\mathrm{Var}(X)} = \dfrac{112}{800} = 0{,}14$ millier de FCFA par heure.
\textbf{4. Ordonnée à l'origine :} $b = \bar{y} - a\bar{x} = 28 - 0{,}14 \times 140 = 28 - 19{,}6 = 8{,}4$.
\textbf{5. Équation :} $y = 0{,}14x + 8{,}4$. Vérification du passage par $G$ : $0{,}14 \times 140 + 8{,}4 = 19{,}6 + 8{,}4 = 28 = \bar{y}$. \checkmark
\textbf{6. Prévision :} pour $x = 200$ heures : $y = 0{,}14 \times 200 + 8{,}4 = 28 + 8{,}4 = 36{,}4$, soit un coût prévisionnel de \num{36400} FCFA. Comme 200 heures est légèrement au-delà des données observées (100 à 180), on signalera qu'il s'agit d'une \emph{extrapolation}, à utiliser avec prudence.
\textbf{Interprétation :} chaque heure de fonctionnement supplémentaire coûte en moyenne $0{,}14$ millier, soit 140 FCFA de maintenance.
\end{exoresolu}

\courssub{Ne pas confondre : régression de y en x et de x en y}

\begin{definition}[Les deux droites de régression]
Il existe \textbf{deux} droites des moindres carrés, qui répondent à deux questions différentes :
\begin{itemize}
  \item la droite de régression de \cle{$y$ en $x$} : $y = ax + b$, avec $a = \dfrac{\mathrm{cov}(X,Y)}{\mathrm{Var}(X)}$, qui minimise les écarts \emph{verticaux} ; on l'utilise pour \textbf{prédire $Y$ connaissant $X$} ;
  \item la droite de régression de \cle{$x$ en $y$} : $x = a'y + b'$, avec $a' = \dfrac{\mathrm{cov}(X,Y)}{\mathrm{Var}(Y)}$, qui minimise les écarts \emph{horizontaux} ; on l'utilise pour \textbf{prédire $X$ connaissant $Y$}.
\end{itemize}
Les deux droites passent toutes les deux par $G(\bar{x}\,;\,\bar{y})$, mais leurs coefficients directeurs sont différents en général : on a $a \times a' = r^2$, donc $a' = a$ si et seulement si $r = \pm 1$. Elles sont confondues si et seulement si $r = \pm 1$ (corrélation parfaite : tous les points sont alignés). Plus $|r|$ est proche de 1, plus les deux droites sont proches l'une de l'autre.
\end{definition}

\begin{attention}[Erreurs fréquentes à éviter]
\begin{itemize}
  \item \textbf{Inverser $X$ et $Y$.} Pour prévoir une production ($Y$) à partir de la pluie ($X$), il faut la droite de $y$ en $x$, avec $a = \dfrac{\mathrm{cov}(X,Y)}{\mathrm{Var}(X)}$. Utiliser la droite de $x$ en $y$ « à l'envers » (en isolant $y$) ne donne \emph{pas} la bonne droite : c'est l'erreur la plus classique à l'examen. Repère : au dénominateur, on met la variance de la variable \emph{connue}.
  \item \textbf{Extrapoler sans le dire.} Le modèle n'est validé que sur la plage des valeurs observées de $X$. Prévoir la production pour \SI{3000}{mm} de pluie alors que les données vont de 800 à \SI{1800}{mm} est une extrapolation risquée : si vous le faites, signalez-le explicitement dans votre conclusion.
  \item \textbf{Oublier les unités.} $a$ a toujours une unité : celle de $Y$ divisée par celle de $X$ (ici t/mm). Une interprétation sans unité est une copie incomplète.
  \item \textbf{Confondre $r$ et $a$.} $r$ est sans unité et compris entre $-1$ et $1$ ; $a$ a une unité et peut prendre n'importe quelle valeur. Ils ont le même signe, mais ce ne sont pas les mêmes nombres.
\end{itemize}
\end{attention}

\begin{methode}[À la calculatrice]
La droite de régression s'obtient dans le même menu que le coefficient $r$ : mode \textbf{STAT}, saisie des deux listes, puis \textbf{REG LIN} (régression linéaire, parfois notée $y = ax + b$ ou $y = A + Bx$ selon les modèles). La calculatrice affiche directement $a$ (ou $B$) et $b$ (ou $A$), ainsi que $r$. \textbf{Vigilance :} vérifiez sur votre modèle si la formule affichée est $y = ax + b$ ou $y = A + Bx$ — dans le second cas, c'est $B$ qui est le coefficient directeur et $A$ l'ordonnée à l'origine. Notez toujours les valeurs avec au moins 3 chiffres significatifs avant d'écrire l'équation finale.
\end{methode}

\begin{savaistu}[Legendre, Gauss et l'astéroïde Cérès]
La méthode des moindres carrés a été publiée en 1805 par le mathématicien français Adrien-Marie Legendre, puis revendiquée par Carl Friedrich Gauss en 1809, qui affirmait l'utiliser depuis 1795. L'enjeu était astronomique : en 1801, l'astéroïde Cérès avait été observé quelques semaines puis « perdu » derrière le Soleil. Grâce aux moindres carrés, Gauss ajusta une orbite aux quelques observations disponibles et prédit la position où Cérès fut effectivement retrouvé fin 1801. Aujourd'hui, cette même méthode est utilisée partout : prévision des ventes dans les entreprises camerounaises, étalonnage des capteurs en électricité, modèles météorologiques, et même entraînement des intelligences artificielles.
\end{savaistu}

\begin{exercice}[À vous de jouer]
Une menuiserie de Yaoundé relève le nombre $x$ de meubles produits par semaine et la consommation électrique $y$ (en kWh) :
\begin{center}
\begin{tabular}{|l|cccccc|}\hline
\rowcolor{popblueL} $x$ (meubles) & 10 & 15 & 20 & 25 & 30 & 35 \\ \hline
$y$ (kWh) & 120 & 150 & 175 & 200 & 240 & 265 \\ \hline
\end{tabular}
\end{center}
On donne $\bar{x} = 22{,}5$, $\bar{y} \approx 191{,}7$, $\sigma_X \approx 8{,}54$, $\sigma_Y \approx 51{,}3$, $r \approx 0{,}998$.
\begin{enumerate}
  \item Calculer $a$ et $b$, puis écrire l'équation de la droite de régression de $y$ en $x$.
  \item Vérifier que la droite passe par le point moyen $G$.
  \item Interpréter concrètement le coefficient $a$.
  \item Prévoir la consommation pour une production de 40 meubles. Ce résultat est-il une interpolation ou une extrapolation ?
\end{enumerate}
\end{exercice}

\begin{corrige}[Exercice : menuiserie]
\textbf{1.} $a = r\,\dfrac{\sigma_Y}{\sigma_X} \approx 0{,}998 \times \dfrac{51{,}3}{8{,}54} \approx 6{,}0$ kWh/meuble. Puis $b = \bar{y} - a\bar{x} \approx 191{,}7 - 6{,}0 \times 22{,}5 = 191{,}7 - 135 = 56{,}7$. Équation : $y = 6x + 56{,}7$.
\textbf{2.} $6 \times 22{,}5 + 56{,}7 = 135 + 56{,}7 = 191{,}7 = \bar{y}$ : la droite passe bien par $G(22{,}5\,;\,191{,}7)$.
\textbf{3.} Chaque meuble supplémentaire produit entraîne en moyenne une consommation additionnelle d'environ 6 kWh.
\textbf{4.} Pour $x = 40$ : $y = 6 \times 40 + 56{,}7 = 296{,}7$ kWh. Comme 40 dépasse la plus grande valeur observée (35), il s'agit d'une \emph{extrapolation} : la prévision est plausible (corrélation très forte, $r \approx 0{,}998$) mais doit être annoncée avec prudence.
\end{corrige}

En résumé, nous savons désormais construire le modèle linéaire $y = ax + b$, le justifier par le critère des moindres carrés et l'interpréter concrètement. Reste une question essentielle : ce modèle est-il \emph{fiable} ? Peut-on l'utiliser pour prévoir, et comment communiquer les résultats ? C'est l'objet de la section suivante, consacrée à la validation du modèle, à la prévision et à la communication.

\coursec{Validation du modèle, Prévision et Communication}

Dans la section précédente, nous avons appris à déterminer l'équation de la droite de régression des moindres carrés $y = ax + b$ à partir d'un nuage de points. Mais obtenir une équation ne suffit pas : la calculatrice ou le tableur fournit \emph{toujours} une droite, même lorsque le nuage n'a aucune forme allongée ! La question qui se pose maintenant est donc celle de l'expert : \textbf{ce modèle est-il de bonne qualité ? Peut-on l'utiliser pour prévoir ? Comment rédiger une conclusion professionnelle ?} C'est l'objet de cette section, qui transforme un simple calcul en une véritable démarche d'expertise statistique, telle qu'elle est pratiquée dans les entreprises, les cabinets d'études et les administrations camerounaises.

\courssub{Analyse des résidus : le modèle sous le microscope}

\textbf{Intuition.} La droite de régression passe « au milieu » du nuage, mais elle ne passe en général par \emph{aucun} point exactement. Pour chaque point observé $(x_i, y_i)$, il existe un petit écart vertical entre la valeur réellement observée $y_i$ et la valeur $\hat{y}_i = a x_i + b$ prédite par le modèle. Cet écart, appelé \cle{résidu}, est la « part d'erreur » du modèle sur ce point. Étudier les résidus, c'est examiner les erreurs du modèle une par une : si ces erreurs ressemblent à du pur hasard, le modèle est bon ; si elles dessinent une structure (une courbe, un entonnoir), c'est que le modèle linéaire rate quelque chose.

\begin{definition}[Résidu]
Pour une série statistique double $(x_i, y_i)$ ajustée par la droite $y = ax + b$, on appelle \cle{résidu} du point d'indice $i$ le nombre
\[ e_i = y_i - \hat{y}_i \quad \text{où} \quad \hat{y}_i = a x_i + b. \]
Le résidu est donc l'écart vertical, compté positivement si le point est au-dessus de la droite et négativement s'il est en dessous.
\end{definition}

\begin{propriete}[Propriété des résidus (admise)]
La somme des résidus d'un ajustement par les moindres carrés est toujours nulle :
\[ \sum_{i=1}^{n} e_i = 0. \]
Autrement dit, les erreurs positives compensent exactement les erreurs négatives : la droite ne « sur-estime » ni ne « sous-estime » globalement les valeurs.
\end{propriete}

Le critère de la somme nulle ne suffit pas à valider le modèle : une droite complètement inadaptée peut aussi avoir des résidus de somme nulle. C'est pourquoi on trace le \cle{nuage des résidus}, c'est-à-dire le nuage des points $(x_i, e_i)$, et on applique les trois critères suivants.

\begin{methode}[Valider un modèle linéaire]
Pour juger de la qualité d'un ajustement linéaire, on procède en quatre étapes :
\begin{enumerate}
\item Regarder le coefficient de corrélation $r$ et le coefficient de détermination $R^2$ : $|r|$ proche de 1 (en pratique $|r| \geq 0{,}87$, c'est-à-dire $R^2 \geq 0{,}75$) indique une liaison forte.
\item Tracer le nuage des résidus $(x_i, e_i)$ à la calculatrice ou au tableur.
\item Vérifier que les résidus sont \textbf{centrés sur 0}, répartis dans une \textbf{bande horizontale de dispersion constante} (homoscédasticité), et qu'ils ne présentent \textbf{aucune structure courbe} (pas de forme en « U » ou en « $\cap$ »).
\item Conclure : si les trois conditions sont remplies, le modèle linéaire est validé ; sinon, il faut envisager un autre modèle (voir section 5).
\end{enumerate}
\end{methode}

\begin{popfigure}
\begin{tikzpicture}
\begin{axis}[
  width=\linewidth, height=6.2cm,
  axis lines=middle,
  xlabel={$x$}, ylabel={$e$},
  xmin=0, xmax=11, ymin=-1.6, ymax=1.6,
  xtick={2,4,6,8,10}, ytick={-1,0,1},
  title={\textbf{Résidus idéaux : bande horizontale aléatoire}},
  title style={font=\small},
  every axis plot/.append style={only marks, mark=*, mark size=1.8pt, color=popblue}
]
\addplot coordinates {(1,0.4) (2,-0.7) (3,0.2) (4,0.8) (5,-0.3) (6,-0.9) (7,0.5) (8,-0.2) (9,0.9) (10,-0.5)};
\addplot[popcurve, domain=0:11, samples=2, color=popdark]{0};
\end{axis}
\end{tikzpicture}

\vspace{0.4cm}

\begin{tikzpicture}
\begin{axis}[
  width=\linewidth, height=6.2cm,
  axis lines=middle,
  xlabel={$x$}, ylabel={$e$},
  xmin=0, xmax=11, ymin=-1.6, ymax=1.6,
  xtick={2,4,6,8,10}, ytick={-1,0,1},
  title={\textbf{Résidus en « U » : le modèle linéaire est inadapté}},
  title style={font=\small},
  every axis plot/.append style={only marks, mark=*, mark size=1.8pt, color=poporange}
]
\addplot coordinates {(1,1.2) (2,0.6) (3,0.1) (4,-0.5) (5,-0.9) (6,-1.0) (7,-0.7) (8,-0.2) (9,0.5) (10,1.1)};
\addplot[popcurve, domain=0:11, samples=2, color=popdark]{0};
\end{axis}
\end{tikzpicture}
\legende{En haut : résidus répartis aléatoirement dans une bande horizontale autour de 0, le modèle linéaire est validé. En bas : les résidus forment une courbe en « U » (positifs aux extrémités, négatifs au centre) : la relation est en réalité courbe, la droite est un mauvais modèle.}
\end{popfigure}

\begin{attention}[Signaux d'alerte dans le nuage des résidus]
Deux configurations doivent vous alerter :
\begin{itemize}
\item \textbf{Structure courbe} (forme en « U » ou en « $\cap$ ») : la relation entre $x$ et $y$ n'est pas linéaire, même si $r$ semble correct. Le modèle sous-estime aux extrémités et surestime au centre (ou l'inverse).
\item \textbf{Dispersion non constante} (forme d'entonnoir, résidus de plus en plus étalés quand $x$ grandit) : la précision des prévisions dépend de la zone considérée ; le modèle est moins fiable pour les grandes valeurs de $x$.
\end{itemize}
Dans les deux cas, il est incorrect de valider l'ajustement linéaire.
\end{attention}

\courssub{Le coefficient de détermination $R^2$ : la note globale du modèle}

\textbf{Intuition.} La variable $Y$ varie : par exemple, la production d'une coopérative varie d'une année à l'autre. Une partie de cette variation est \emph{expliquée} par la variation de $X$ (la pluviométrie, via la droite de régression) ; le reste est dû à d'autres facteurs (maladies, engrais, main-d'œuvre) et se retrouve dans les résidus. Le coefficient $R^2$ mesure la \textbf{part de la variance de $Y$ expliquée par le modèle}.

\begin{propriete}[Décomposition de la variance (admise)]
Pour un ajustement par les moindres carrés, la variance totale de $Y$ se décompose en :
\[ \mathrm{Var}(Y) = \mathrm{Var}(\hat{Y}) + \mathrm{Var}(\text{résidus}), \]
où $\mathrm{Var}(\hat{Y})$ est la variance des valeurs prédites (variance \emph{expliquée} par le modèle) et $\mathrm{Var}(\text{résidus}) = \dfrac{1}{n}\sum e_i^2$ la variance \emph{résiduelle} (non expliquée). On en déduit :
\[ R^2 = \frac{\mathrm{Var}(\hat{Y})}{\mathrm{Var}(Y)} = 1 - \frac{\mathrm{Var}(\text{résidus})}{\mathrm{Var}(Y)} = r^2, \]
où $r$ est le coefficient de corrélation linéaire. On a toujours $0 \leq R^2 \leq 1$.
\end{propriete}

\begin{aretenir}[Interpréter $R^2$]
$R^2$ est la proportion de la variance de $Y$ expliquée par la régression. Ainsi :
\begin{itemize}
\item $R^2 = 0{,}92$ signifie : « 92 \% de la variabilité de $Y$ est expliquée par le modèle linéaire en $X$ » ;
\item en sciences techniques et sociales, on considère en général l'ajustement comme \textbf{bon} lorsque $R^2 > 0{,}75$ ;
\item $R^2$ proche de 0 signifie que la droite n'explique presque rien : le modèle linéaire est à rejeter.
\end{itemize}
\end{aretenir}

\courssub{Prévision : interpolation et extrapolation}

Une fois le modèle validé, on peut l'utiliser pour \cle{prévoir} la valeur de $Y$ pour une valeur $x_0$ non observée : il suffit de calculer $\hat{y}_0 = a x_0 + b$. Mais toutes les prévisions ne se valent pas, et la distinction suivante est capitale.

\begin{definition}[Interpolation et extrapolation]
Soit $[x_{\min} ; x_{\max}]$ l'intervalle des valeurs observées de $X$.
\begin{itemize}
\item On parle d'\cle{interpolation} lorsque $x_0 \in [x_{\min} ; x_{\max}]$ : la prévision s'appuie sur le modèle à l'intérieur de la zone où l'on dispose de données. Sa fiabilité dépend de $R^2$ et de la dispersion des résidus.
\item On parle d'\cle{extrapolation} lorsque $x_0 \notin [x_{\min} ; x_{\max}]$ : on prolonge la droite \emph{au-delà} des données. Rien ne garantit que la relation linéaire reste valable (saturation, seuils, effets nouveaux) : l'extrapolation est \textbf{risquée} et doit toujours être signalée comme telle.
\end{itemize}
\end{definition}

\begin{attention}[Deux erreurs fréquentes à l'examen]
\begin{itemize}
\item \textbf{Confondre corrélation et causalité.} On écrit « la pluie est \emph{associée à} la production » et non « la pluie \emph{cause} la production » : une forte corrélation peut cacher une troisième variable (qualité des sols, politique agricole) ou une coïncidence. Le vocabulaire juste est : liaison, association, dépendance statistique.
\item \textbf{Présenter une prévision comme une certitude.} L'intervalle de confiance n'est pas au programme, mais il faut \emph{qualifier} la prévision : « prévision approximative », « fiable » (interpolation avec $R^2$ élevé) ou « hasardeuse » (extrapolation). Ne jamais écrire « la production \emph{sera} de 11,4 t » mais « le modèle \emph{estime} la production à environ 11,4 t ».
\end{itemize}
\end{attention}

\courssub{Exemple complet et rédaction d'un compte-rendu}

Reprenons l'exemple fil rouge de la leçon : une coopérative agricole de l'Ouest du Cameroun a relevé, pendant 8 années, la pluviométrie $x$ (en mm) et la production de maïs $y$ (en tonnes). Les valeurs de $x$ vont de \SI{1200}{mm} à \SI{2400}{mm}. Les calculs des sections précédentes ont donné :
\[ r = 0{,}95, \qquad R^2 = 0{,}90, \qquad y = 0{,}0045\,x - 2{,}1. \]

\begin{exoresolu}[Prévisions et rédaction de la conclusion]
\textbf{Énoncé.} (1) Prévoir la production pour une année où il tomberait \SI{2000}{mm} de pluie. (2) Prévoir la production pour \SI{3000}{mm}. (3) Rédiger la conclusion complète de l'étude.
\tcblower
\textbf{Solution.}
(1) Pour $x_0 = 2000$ : cette valeur appartient à $[1200 ; 2400]$, c'est une \textbf{interpolation}.
\[ \hat{y}_0 = 0{,}0045 \times 2000 - 2{,}1 = 9 - 2{,}1 = 6{,}9. \]
Le modèle estime la production à environ \textbf{6,9 tonnes}. Comme $R^2 = 0{,}90$ et qu'il s'agit d'une interpolation, cette prévision est \textbf{fiable} (à la dispersion des résidus près).

(2) Pour $x_0 = 3000$ : cette valeur est \emph{hors} de l'intervalle observé, c'est une \textbf{extrapolation}.
\[ \hat{y}_0 = 0{,}0045 \times 3000 - 2{,}1 = 13{,}5 - 2{,}1 = 11{,}4. \]
Le calcul donne 11,4 tonnes, mais cette prévision est \textbf{hasardeuse} : au-delà d'un certain seuil, les sols se saturent d'eau, les maladies se développent et la production peut stagner, voire baisser. La relation linéaire ne tient probablement plus.

(3) \textbf{Conclusion type.} \\
\emph{Contexte.} On a étudié la liaison entre la pluviométrie annuelle $x$ (mm) et la production de maïs $y$ (t) d'une coopérative, sur 8 années ($x \in [1200 ; 2400]$). \\
\emph{Liaison.} Le coefficient de corrélation $r = 0{,}95$ indique une corrélation linéaire positive très forte : la pluviométrie et la production sont fortement associées, dans le sens où les années pluvieuses correspondent aux productions élevées. \\
\emph{Modèle.} La droite des moindres carrés $y = 0{,}0045\,x - 2{,}1$ explique 90 \% de la variance de la production ($R^2 = 0{,}90$) ; le nuage des résidus ne présente ni structure courbe ni dispersion anormale : le modèle linéaire est validé. \\
\emph{Prévision.} Pour \SI{2000}{mm} (interpolation), le modèle estime la production à environ 6,9 t, prévision jugée fiable. \\
\emph{Limites.} La prévision de 11,4 t pour \SI{3000}{mm} est une extrapolation, à manier avec prudence (saturation des sols, maladies). Enfin, la corrélation constatée ne prouve pas une causalité directe : d'autres facteurs (engrais, variétés semencières) interviennent.
\end{exoresolu}

\begin{methode}[Rédiger une conclusion statistique en 5 points]
Toute conclusion d'étude statistique à deux variables suit le plan du « rapport d'expertise » :
\begin{enumerate}
\item \textbf{Contexte} : quelles variables, quelle population, combien d'observations, quelle plage de valeurs ?
\item \textbf{Liaison} : valeur de $r$, sens (positive/négative), intensité (faible/forte).
\item \textbf{Modèle} : équation de la droite, valeur de $R^2$, résultat de l'analyse des résidus.
\item \textbf{Prévision chiffrée} : calcul de $\hat{y}_0$, en précisant interpolation ou extrapolation et le degré de fiabilité.
\item \textbf{Limites} : extrapolation, absence de preuve de causalité, influence possible de points aberrants (outliers).
\end{enumerate}
\end{methode}

\begin{popfigure}
\begin{tikzpicture}[
  node distance=0.35cm,
  box/.style={draw=popblue, fill=popblueL, rounded corners=2pt, text width=0.86\linewidth, align=center, font=\small, inner sep=5pt},
  arr/.style={-{Stealth[length=2.5mm]}, thick, color=popdark}
]
\node[box] (c) {\textbf{1. Contexte} : variables $X$ et $Y$, effectif, plage des données};
\node[box, below=of c] (a) {\textbf{2. Analyse} : nuage de points, coefficient $r$, sens et intensité};
\node[box, below=of a] (m) {\textbf{3. Modèle} : droite $y=ax+b$, $R^2$, validation par les résidus};
\node[box, below=of m] (p) {\textbf{4. Prévision} : $\hat{y}_0 = a x_0 + b$, interpolation ou extrapolation, fiabilité};
\node[box, below=of p, draw=poporange, fill=poporangeL] (r) {\textbf{5. Recommandations} : limites du modèle, prudence, décision};
\draw[arr] (c) -- (a);
\draw[arr] (a) -- (m);
\draw[arr] (m) -- (p);
\draw[arr] (p) -- (r);
\end{tikzpicture}
\legende{Structure d'un compte-rendu statistique type « rapport d'expertise » : du contexte vers la recommandation, chaque étape s'appuie sur la précédente.}
\end{popfigure}

\courssub{Cas particulier : et si la liaison n'est pas linéaire ?}

Il arrive que le nuage de points présente une forme nettement courbe (croissance rapide puis ralentissement, décroissance exponentielle, etc.). Dans ce cas, le coefficient $r$ peut être trompeur et le nuage des résidus montre une structure en « U » : l'ajustement linéaire direct est à rejeter.

\begin{aretenir}[Réflexe face à une liaison non linéaire]
Lorsque le nuage est courbe, une démarche possible consiste à \textbf{transformer} une des variables (poser $z = \ln y$ ou $z = \sqrt{y}$, par exemple) puis à vérifier si le nouveau nuage $(x_i, z_i)$ est rectiligne : on se ramène alors à un ajustement linéaire sur les données transformées. Cette technique de \cle{linéarisation} est étudiée en section 5 (complément culturel, hors programme strict). À l'examen, l'essentiel est de \emph{savoir détecter} la non-linéarité (nuage courbe, résidus structurés) et de \emph{refuser} l'ajustement linéaire dans ce cas.
\end{aretenir}

\begin{exercice}[Validation et prévision]
Un atelier de mécanique note, pour 6 machines, leur âge $x$ (en années) et leur coût annuel d'entretien $y$ (en milliers de FCFA). On obtient : $r = 0{,}89$, droite de régression $y = 15x + 20$, âges observés entre 1 et 9 ans.
\begin{enumerate}
\item Calculer $R^2$ et commenter la qualité de l'ajustement.
\item Prévoir le coût d'entretien d'une machine de 5 ans. Qualifier cette prévision.
\item Prévoir le coût pour une machine de 15 ans. Que dire de cette prévision ?
\item Le nuage des résidus présente une forme en « U ». Que conclure ?
\end{enumerate}
\end{exercice}

\begin{corrige}[Exercice : Validation et prévision]
\begin{enumerate}
\item $R^2 = r^2 = 0{,}89^2 = 0{,}7921 \approx 0{,}79 > 0{,}75$ : environ 79 \% de la variance du coût est expliquée par l'âge ; l'ajustement semble bon sur le critère de $R^2$.
\item $x_0 = 5 \in [1 ; 9]$ : interpolation. $\hat{y}_0 = 15 \times 5 + 20 = 95$ milliers de FCFA. Prévision a priori fiable (sous réserve du point 4).
\item $x_0 = 15 \notin [1 ; 9]$ : extrapolation. Le calcul donne $15 \times 15 + 20 = 245$ milliers de FCFA, mais cette prévision est hasardeuse : une machine très vieille peut être réformée ou voir son coût croître plus vite que linéairement.
\item Un nuage de résidus en « U » révèle une structure courbe non prise en compte : malgré $R^2 = 0{,}79$, le modèle linéaire est \textbf{inadapté} ; il faut envisager une transformation des variables ou un autre modèle. Cela montre que $R^2$ seul ne suffit jamais à valider un ajustement.
\end{enumerate}
\end{corrige}

\begin{savaistu}[La régression au service de l'État camerounais]
En économétrie camerounaise, l'Institut National de la Statistique (INS) et le Ministère des Finances (MINFI) utilisent quotidiennement des modèles de régression pour prévoir le PIB en fonction de l'investissement public, ou l'inflation en fonction de la masse monétaire en circulation. Ces prévisions, toujours assorties d'une analyse de fiabilité et de marges d'incertitude, servent à élaborer le budget de l'État et à fixer les objectifs de croissance. La démarche que tu viens d'apprendre — modéliser, valider, prévoir, conclure avec prudence — est exactement celle des experts statisticiens nationaux.
\end{savaistu}

\coursec{Complément utile : Ajustement non linéaire par linéarisation (Hors programme strict, culture Terminale)}

\noindent Dans les sections précédentes, nous avons vu comment valider un modèle linéaire, l'utiliser pour prévoir et communiquer nos résultats. Mais que faire lorsque le nuage de points ne s'aligne pas le long d'une droite, mais dessine une courbe nette — croissance rapide, décroissance lente, forme en parabole ? C'est précisément la situation que rencontrent les techniciens face à des phénomènes d'amortissement, de croissance démographique, de cinétique chimique ou de lois d'échelle. Cette section, bien que hors programme strict du Baccalauréat technique, vous donne les clés pour \emph{linéariser} ces relations et réutiliser toute la puissance de la régression des moindres carrés.

\courssub{Quand le nuage n'est pas linéaire : repérer la forme}

\noindent Avant toute modélisation, l'observation du nuage de points reste l'étape reine. Trois formes non linéaires reviennent sans cesse en technique :

\begin{itemize}
  \item \textbf{Croissance ou décroissance exponentielle} : les points suivent une courbe qui s'éloigne ou se rapproche de l'axe des abscisses de façon multiplicative (ex. : population bactérienne, charge/décharge d'un condensateur, amortissement d'une machine, intérêt composé).
  \item \textbf{Loi de puissance} : les points suivent une courbe de type $Y = a X^{b}$ (ex. : loi de Pareto en économie, loi de Zipf en linguistique, relation surface/volume en biophysique, résistance d'un fil en fonction de sa section, loi de Poiseuille en hydraulique).
  \item \textbf{Croissance logistique (sigmoïde)} : croissance qui ralentit vers un plafond (capacité d'accueil d'un étang, diffusion d'une innovation, courbe d'apprentissage). Ce cas demande des outils plus avancés (régression non linéaire itérative) et ne se linéarise pas simplement par un changement de variable élémentaire.
\end{itemize}

\begin{attention}[Piège visuel : nuage en « banane »]
Un nuage courbe peut aussi cacher une relation linéaire avec hétéroscédasticité (variance non constante) ou des valeurs aberrantes. Tracez toujours le nuage \textbf{avant} de choisir un modèle. Si le nuage ressemble à une banane large, une transformation de variable (racine carrée, logarithme) peut stabiliser la variance \emph{et} linéariser la relation.
\end{attention}

\courssub{Principe de la linéarisation : transformer pour retrouver une droite}

\noindent L'idée est simple et puissante : \textbf{appliquer une fonction bijective aux variables} pour transformer la relation courbe en relation affine. Si $Y = f(X)$ est le modèle supposé, on cherche une transformation $u = g(X)$, $v = h(Y)$ telle que $v = \alpha + \beta u$ soit une droite.

\begin{definition}[Linéarisation par transformation]
Soit un modèle $Y = \varphi(X; \theta)$ non linéaire en les paramètres $\theta$. Si l'on trouve des fonctions $g, h$ telles que $h(Y) = \alpha + \beta \, g(X)$ soit affine en les nouveaux paramètres $(\alpha, \beta)$, alors la régression linéaire sur les données transformées $(g(x_i), h(y_i))$ fournit des estimateurs de $\alpha, \beta$, d'où l'on déduit $\theta$.
\end{definition}

\noindent Les deux cas classiques au programme de culture technique sont :

\begin{propriete}[Transformations usuelles pour linéariser]
\begin{enumerate}
  \item \textbf{Modèle exponentiel} $Y = a \, b^{X}$ avec $a>0, b>0$.\\
        Poser $v = \ln Y$, $u = X$. Alors $v = \ln a + (\ln b) \, u$ : droite d'équation $v = \alpha + \beta u$ avec $\alpha = \ln a$, $\beta = \ln b$.
  \item \textbf{Modèle puissance} $Y = a \, X^{b}$ avec $a>0, X>0$.\\
        Poser $v = \ln Y$, $u = \ln X$. Alors $v = \ln a + b \, u$ : droite d'équation $v = \alpha + \beta u$ avec $\alpha = \ln a$, $\beta = b$.
\end{enumerate}
\end{propriete}

\begin{savaistu}[Log-log, lin-log, log-lin : le vocabulaire des ingénieurs]
En pratique, on parle de graphique \textbf{lin-log} (axe Y logarithmique, axe X linéaire) pour l'exponentielle, et de graphique \textbf{log-log} (les deux axes logarithmiques) pour la puissance. Sur papier millimétrique logarithmique, ces courbes deviennent des droites : c'était la méthode graphique avant les calculatrices. Aujourd'hui, on fait la transformation dans les listes de la calculatrice (ou Excel, Python) et on lance une \texttt{LinReg}.
\end{savaistu}

\courssub{Modèle exponentiel $Y = a b^{x}$ : linéarisation par le logarithme népérien}

\noindent Considérons un phénomène où la variation relative est constante : population qui double chaque période, condensateur qui se décharge, valeur d'une machine qui perd 15\% par an. Le modèle théorique est $Y = a b^{x}$.

\begin{experience}[Démonstration guidée : pourquoi $\ln Y$ linéarise l'exponentielle]
Partons de $Y = a b^{x}$ avec $a>0, b>0$. Appliquons $\ln$ des deux côtés :
\[
\ln Y = \ln(a b^{x}) = \ln a + \ln(b^{x}) = \ln a + x \ln b.
\]
Posons $Y' = \ln Y$, $\alpha = \ln a$, $\beta = \ln b$. L'équation devient $Y' = \alpha + \beta x$, qui est une \textbf{droite} dans le repère $(x, Y')$.
\begin{enumerate}
  \item Si on trace le nuage $(x_i, \ln y_i)$, les points s'alignent (si le modèle est bon).
  \item La régression linéaire $Y' = \hat{\alpha} + \hat{\beta} x$ donne $\hat{\alpha}, \hat{\beta}$.
  \item On retrouve $a = e^{\hat{\alpha}}$, $b = e^{\hat{\beta}}$.
\end{enumerate}
\end{experience}

\noindent Sur calculatrice (mode \texttt{STAT} $\to$ \texttt{REG} $\to$ \texttt{LinReg $ax+b$}) :
\begin{enumerate}
  \item Liste L1 = $x_i$ (abscisses inchangées).
  \item Liste L2 = $y_i$ (mesures brutes).
  \item Créer L3 = $\ln($L2$)$ $\to$ \texttt{LN(L2)} $\to$ \texttt{STO$\to$ L3}.
  \item Lancer \texttt{LinReg(ax+b) L1, L3}.
  \item Lire $a_{\text{lin}}$ (pente) et $b_{\text{lin}}$ (ordonnée à l'origine) affichés.
  \item Modèle final : $\boxed{Y = e^{b_{\text{lin}}} \cdot (e^{a_{\text{lin}}})^{x}}$.
\end{enumerate}

\begin{methode}[Linéariser un modèle exponentiel $Y = a b^{x}$]
\begin{enumerate}
  \item Vérifier que tous les $y_i > 0$ (logarithme défini).
  \item Dans la calculatrice : L1 $\leftarrow x_i$, L2 $\leftarrow y_i$, L3 $\leftarrow \ln(\text{L2})$.
  \item Faire \texttt{LinReg(ax+b) L1, L3} $\to$ récupérer $a_{\text{lin}}$ (pente) et $b_{\text{lin}}$ (ordonnée à l'origine).
  \item Calculer $a = \exp(b_{\text{lin}})$, $b = \exp(a_{\text{lin}})$.
  \item Écrire le modèle : $Y = a \cdot b^{x}$.
  \item \textbf{Valider} : tracer le nuage $(x_i, y_i)$ et la courbe $Y = a b^{x}$ superposée ; examiner les résidus $y_i - a b^{x_i}$.
\end{enumerate}
\end{methode}

\courssub{Modèle puissance $Y = a x^{b}$ : linéarisation log-log}

\noindent Ce modèle apparaît dès qu'une grandeur scale comme une puissance d'une autre : surface $\propto$ (longueur)$^2$, débit $\propto$ (diamètre)$^4$ (loi de Poiseuille), fréquence cardiaque $\propto$ (masse)$^{-1/4}$, loi de Pareto (richesse $\propto$ rang$^{-1}$).

\begin{experience}[Démonstration : linéarisation du modèle puissance]
Partons de $Y = a X^{b}$ ($a>0, X>0$). Appliquons $\ln$ :
\[
\ln Y = \ln(a X^{b}) = \ln a + b \ln X.
\]
Posons $Y' = \ln Y$, $X' = \ln X$, $\alpha = \ln a$, $\beta = b$. L'équation devient $Y' = \alpha + \beta X'$, une droite dans le repère $(\ln X, \ln Y)$.
\begin{enumerate}
  \item Transformer les données : $x'_i = \ln x_i$, $y'_i = \ln y_i$ (exige $x_i>0, y_i>0$).
  \item Régression linéaire sur $(x'_i, y'_i)$ $\to$ $\hat{\alpha}, \hat{\beta}$.
  \item $a = e^{\hat{\alpha}}$, $b = \hat{\beta}$.
  \item Modèle final : $\boxed{Y = e^{\hat{\alpha}} \cdot X^{\hat{\beta}}}$.
\end{enumerate}
\end{experience}

\begin{methode}[Linéariser un modèle puissance $Y = a X^{b}$]
\begin{enumerate}
  \item Vérifier $x_i > 0$ et $y_i > 0$ pour tout $i$.
  \item L1 $\leftarrow \ln(x_i)$, L2 $\leftarrow \ln(y_i)$.
  \item \texttt{LinReg(ax+b) L1, L2} $\to$ $\hat{\alpha}$ (ordonnée à l'origine), $\hat{\beta}$ (pente).
  \item $a = \exp(\hat{\alpha})$, $b = \hat{\beta}$.
  \item Modèle : $Y = a X^{b}$.
  \item Valider par tracé du nuage $(x_i, y_i)$ et courbe théorique ; résidus.
\end{enumerate}
\end{methode}

\courssub{Mise en œuvre sur calculatrice : exemple complet exponentiel}

\noindent Illustrons la procédure sur un cas concret simulé : croissance d'une population de poissons en étang.

\begin{exoresolu}[Croissance d'une population de poissons en étang]
\textbf{Énoncé.} Un pisciculteur suit la biomasse totale (en kg) d'une population de tilapias dans un étang de 500 m$^2$ sur 10 mois. Données relevées :

\begin{center}
\begin{tabular}{|c|cccccccccc|}\hline
Mois $x$ & 1 & 2 & 3 & 4 & 5 & 6 & 7 & 8 & 9 & 10 \\ \hline
Biomasse $y$ (kg) & 12 & 15 & 19 & 24 & 30 & 38 & 48 & 61 & 77 & 97 \\ \hline
\end{tabular}
\end{center}

On suspecte une croissance exponentielle (ressources non limitantes au début). Ajuster un modèle $Y = a b^{x}$ par linéarisation. Prévoir la biomasse au mois 12. Commenter la validité.

\textbf{Solution détaillée.}

\textit{Étape 1 : Observation.} Le nuage $(x, y)$ courbe vers le haut, croissance accélérée $\to$ hypothèse exponentielle plausible.

\textit{Étape 2 : Transformation.} Calculer $\ln y$ (arrondi à 4 décimales) :

\begin{center}
\begin{tabular}{|c|cccccccccc|}\hline
$x$ & 1 & 2 & 3 & 4 & 5 & 6 & 7 & 8 & 9 & 10 \\ \hline
$y$ & 12 & 15 & 19 & 24 & 30 & 38 & 48 & 61 & 77 & 97 \\ \hline
$\ln y$ & 2,4849 & 2,7081 & 2,9444 & 3,1781 & 3,4012 & 3,6376 & 3,8712 & 4,1109 & 4,3438 & 4,5747 \\ \hline
\end{tabular}
\end{center}

\textit{Étape 3 : Régression linéaire sur $(x, \ln y)$.}
Calculs (formules du cours ou calculatrice) :
\begin{align*}
\bar{x} &= 5,5, \quad \overline{\ln y} \approx 3,5255, \\
S_{xx} &= \sum (x_i - \bar{x})^2 = 82,5, \\
S_{x,\ln y} &= \sum (x_i - \bar{x})(\ln y_i - \overline{\ln y}) \approx 19,10.
\end{align*}
Pente $\hat{\beta} = \dfrac{S_{x,\ln y}}{S_{xx}} \approx \dfrac{19,10}{82,5} \approx 0,2316$.
Ordonnée à l'origine $\hat{\alpha} = \overline{\ln y} - \hat{\beta}\bar{x} \approx 3,5255 - 0,2316 \times 5,5 \approx 2,2519$.

\textit{Étape 4 : Retour au modèle exponentiel.}
$a = e^{\hat{\alpha}} \approx e^{2,2519} \approx 9,51$,
$b = e^{\hat{\beta}} \approx e^{0,2316} \approx 1,261$.

Modèle ajusté : $\boxed{Y \approx 9,51 \times 1,261^{x}}$.

\textit{Étape 5 : Prévision mois 12.}
$Y(12) \approx 9,51 \times 1,261^{12} \approx 9,51 \times 16,09 \approx 153,0$ kg.

\textit{Étape 6 : Validation / Commentaire.}
Le coefficient de détermination sur données transformées : $R^2_{\text{lin}} \approx 0,990$ (très fort alignement sur $\ln y$).
Mais \textbf{attention} : ce $R^2$ mesure la qualité d'ajustement de $\ln Y$, pas de $Y$ directement. Les résidus sur $Y$ (écarts $y_i - 9,51 \times 1,261^{x_i}$) sont :
$+0,02 ; -0,10 ; -0,03 ; +0,01 ; -0,24 ; -0,12 ; -0,05 ; +0,43 ; +0,65 ; +0,76$ (kg).
Ils sont globalement petits mais présentent une légère structure en « U » (négatifs au centre, positifs aux extrémités), suggérant que la croissance réelle est peut-être un peu plus lente que l'exponentielle pure à long terme (amorce de saturation). Le modèle reste acceptable \emph{à court terme}.
\emph{Prudence} : au mois 12, la prévision 153 kg suppose ressources illimitées. En réalité, l'étang a une capacité d'accueil (modèle logistique). Cette prévision n'est fiable que tant que la biomasse reste bien en deçà du plafond.
\end{exoresolu}

\courssub{Illustration graphique : nuage initial et nuage linéarisé}

\noindent La figure ci-dessous montre les deux étapes visuelles : le nuage brut (courbe exponentielle) et le nuage transformé $(\ln y$ en fonction de $x)$ aligné sur la droite de régression.

\begin{popfigure}
\begin{tikzpicture}
\begin{groupplot}[
  group style={group size=2 by 1, horizontal sep=2.5cm},
  width=0.45\linewidth,
  height=6cm,
  tick label style={font=\small},
  label style={font=\small},
  grid=major,
  grid style={popline!50},
]
% Premier graphique : nuage exponentiel brut
\nextgroupplot[
  title={Nuage initial $(x, y)$ : allure exponentielle},
  xlabel={Mois $x$},
  ylabel={Biomasse $y$ (kg)},
  xmin=0, xmax=12,
  ymin=0, ymax=170,
  xtick={1,2,...,10},
  ytick={0,20,40,60,80,100,120,140,160},
]
% Points de données
\addplot[only marks, mark=*, mark size=2.5, popblue] coordinates {
  (1,12) (2,15) (3,19) (4,24) (5,30) (6,38) (7,48) (8,61) (9,77) (10,97)
};
% Courbe modèle ajusté Y = 9.51 * 1.261^x
\addplot[popcurve, poporange, domain=0:12, samples=100] {9.51 * 1.261^x};
\addlegendentry{Données}
\addlegendentry{Modèle $Y=9,51\times1,261^x$}

% Deuxième graphique : nuage linéarisé (x, ln y)
\nextgroupplot[
  title={Nuage linéarisé $(x, \ln y)$ : alignement linéaire},
  xlabel={Mois $x$},
  ylabel={$\ln(y)$},
  xmin=0, xmax=11,
  ymin=2, ymax=5,
  xtick={1,2,...,10},
  ytick={2,2.5,3,3.5,4,4.5,5},
]
% Points transformés
\addplot[only marks, mark=*, mark size=2.5, popblue] coordinates {
  (1,2.4849) (2,2.7081) (3,2.9444) (4,3.1781) (5,3.4012) 
  (6,3.6376) (7,3.8712) (8,4.1109) (9,4.3438) (10,4.5747)
};
% Droite de régression : ln y = 2.2519 + 0.2316 x
\addplot[popcurve, poporange, domain=0:11, samples=2] {2.2519 + 0.2316*x};
\addlegendentry{Données $(\ln y)$}
\addlegendentry{Droite $\ln y = 2,252 + 0,232 x$}
\end{groupplot}
\end{tikzpicture}
\legende{À gauche : nuage brut $(x, y)$ et courbe exponentielle ajustée. À droite : nuage transformé $(x, \ln y)$ et droite de régression des moindres carrés. La linéarisation rend l'ajustement visuellement évident.}
\end{popfigure}

\courssub{Validation et limites : attention au $R^2$ et aux résidus}

\noindent C'est le point le plus subtil et le plus souvent mal compris.

\begin{attention}[Erreur fréquente : confondre $R^2$ transformé et $R^2$ réel]
Le coefficient de détermination $R^2$ affiché par la calculatrice après \texttt{LinReg} sur les données transformées porte sur la variable \textbf{transformée} ($\ln Y$ ou $\ln X, \ln Y$), \textbf{pas} sur $Y$ directement.
\begin{itemize}
  \item Un $R^2_{\text{lin}} = 0,99$ sur $\ln Y$ ne garantit \textbf{pas} que le modèle exponentiel explique 99\% de la variance de $Y$.
  \item Les résidus sur $Y$ (échelle originale) ne sont \textbf{plus homogènes} : la variance des erreurs croît souvent avec $Y$ (hétéroscédasticité). La régression sur $\ln Y$ minimise la somme des carrés des écarts \emph{relatifs} (en log), pas des écarts absolus.
  \item \textbf{Deux méthodes donnent des résultats différents :}
  \begin{enumerate}
    \item Linéarisation + \texttt{LinReg} (minimise $\sum (\ln y_i - \ln \hat{y}_i)^2$).
    \item Régression exponentielle native \texttt{ExpReg} de la calculatrice (minimise $\sum (y_i - \hat{y}_i)^2$ sur échelle originale, par algorithme itératif).
  \end{enumerate}
  Les paramètres $a, b$ diffèrent légèrement. Pour un technicien, la méthode par linéarisation est transparente, pédagogique et suffit pour l'essentiel. La méthode \texttt{ExpReg} est statistiquement plus rigoureuse sur l'échelle d'origine mais « boîte noire ».
\end{itemize}
\end{attention}

\begin{propriete}[Critères de validation après linéarisation]
Pour valider un modèle obtenu par linéarisation, on examine :
\begin{enumerate}
  \item \textbf{Nuage transformé} : alignement visuel sur une droite (graphique $(x, \ln y)$ ou $(\ln x, \ln y)$).
  \item \textbf{Résidus sur échelle originale} : $e_i = y_i - \hat{y}_i$. Ils doivent être centrés sur 0, sans structure (pas de courbe résiduelle), et idéalement homoscédastiques (même dispersion pour tous $x$).
  \item \textbf{Coefficient de détermination « pseudo-$R^2$ » sur $Y$} : $R^2_Y = 1 - \dfrac{\sum (y_i - \hat{y}_i)^2}{\sum (y_i - \bar{y})^2}$. Il se calcule à la main ou dans un tableur. C'est lui qui compte pour la prévision.
  \item \textbf{Sens physique} : les paramètres $a, b$ ont-ils du sens ? (ex. $b>1$ pour croissance, $0<b<1$ pour décroissance).
\end{enumerate}
\end{propriete}

\courssub{Autre cas classique : modèle puissance et loi de Pareto}

\noindent La méthode est identique, avec double logarithme. Illustrons par un exemple culturel.

\begin{savaistu}[Loi de Pareto (80/20) et loi de Zipf : la puissance du log-log]
\begin{itemize}
  \item \textbf{Loi de Pareto} (économie, gestion de stock) : 80\% des effets viennent de 20\% des causes. Modèle : $Y = a X^{b}$ avec $b \approx -1,16$ (car $0,8 = 0,2^{b} \Rightarrow b = \ln 0,8 / \ln 0,2 \approx -1,16$). En log-log : $\ln Y = \ln a - 1,16 \ln X$. Une droite de pente $-1,16$.
  \item \textbf{Loi de Zipf} (linguistique, traitement du signal) : fréquence du $r^{\text{ème}}$ mot $\propto 1/r$. Modèle $Y = a r^{-1}$. En log-log : pente $-1$.
  \item \textbf{Échelle allométrique} (biologie, génie civil) : $Y = a M^{b}$ où $M$ est la masse. Métabolisme basal $\propto M^{3/4}$ (loi de Kleiber). Résistance d'un os $\propto$ (diamètre)$^2$.
\end{itemize}
Dans tous ces cas, tracer $\ln Y$ vs $\ln X$ révèle une droite. La pente donne l'exposant $b$, l'ordonnée à l'origine donne $\ln a$. C'est l'outil universel des lois d'échelle.
\end{savaistu}

\courssub{Résumé de la boîte à outils : quand et comment linéariser}

\begin{center}
\begin{tabularx}{\linewidth}{|l|X|X|X|}\hline
\textbf{Modèle suspecté} & \textbf{Transformation} & \textbf{Régression sur} & \textbf{Paramètres finaux} \\ \hline
Exponentiel $Y = a b^{x}$ & $Y' = \ln Y$ & $(x, Y')$ : LinReg $Y' = \alpha + \beta x$ & $a = e^{\alpha},\; b = e^{\beta}$ \\ \hline
Puissance $Y = a x^{b}$ & $X' = \ln X,\; Y' = \ln Y$ & $(X', Y')$ : LinReg $Y' = \alpha + \beta X'$ & $a = e^{\alpha},\; b = \beta$ \\ \hline
Exponentiel base $e$ : $Y = a e^{kx}$ & $Y' = \ln Y$ & $(x, Y')$ : LinReg $Y' = \alpha + k x$ & $a = e^{\alpha},\; k = \beta$ \\ \hline
\end{tabularx}
\end{center}

\begin{aretenir}[À retenir pour la culture technique]
\begin{itemize}
  \item Face à un nuage courbe, \textbf{ne pas forcer une droite}. Cherchez la transformation qui linéarise.
  \item Exponentiel $\to$ $\ln Y$ vs $X$ (lin-log). Puissance $\to$ $\ln Y$ vs $\ln X$ (log-log).
  \item Faites la régression \texttt{LinReg} sur données transformées, puis \textbf{retransformez} les paramètres.
  \item \textbf{Vérifiez toujours} les résidus sur l'échelle d'origine ($Y$). Le $R^2$ affiché par la calculatrice ne s'applique qu'à la variable transformée.
  \item Méthode transparente, reproductible, sans « boîte noire » : idéale pour le technicien qui doit justifier son modèle.
\end{itemize}
\end{aretenir}

\noindent Cette section clôture l'unité « Statistiques à deux variables ». Vous disposez maintenant d'une chaîne complète : \textbf{organiser} (tableau double entrée), \textbf{visualiser} (nuage, $r$), \textbf{modéliser} (droite de régression), \textbf{valider} (résidus, $R^2$, intervalle), \textbf{prévoir} — et, en bonus, \textbf{linéariser} pour aller au-delà du linéaire. Ces compétences sont transverses : électricien (caractéristiques diode), mécanicien (courbe d'usure), gestionnaire (prévisions de ventes), chimiste (cinétique). Le langage commun, c'est la méthode des moindres carrés et l'esprit critique face aux résidus.

\newpage

\coursec{Activité d'intégration}

\coursec{Statistiques à deux variables}

\courssub{Objectifs de l'activité}
Cette activité d'intégration vise à mobiliser l'ensemble des compétences de l'unité \og Statistiques à deux variables \fg dans un contexte professionnel réaliste : la gestion d'un \cle{call-box} (cabine téléphonique publique) à Douala. L'élève devra :
\begin{itemize}
    \item Organiser des données brutes en un \cle{tableau à double entrée} avec choix des classes ;
    \item Visualiser la liaison statistique via un \cle{nuage de points} et quantifier son intensité par le \cle{coefficient de corrélation linéaire $r$} et le \cle{coefficient de détermination $R^2$} ;
    \item Déterminer l'équation de la \cle{droite de régression des moindres carrés} ($y$ en $x$) et l'utiliser pour des prévisions ;
    \item Valider le modèle linéaire par l'\cle{analyse des résidus} ;
    \item Distinguer \cle{interpolation} et \cle{extrapolation} et en comprendre les risques ;
    \item Communiquer les résultats de manière synthétique et professionnelle (SMS) ;
    \item Réfléchir sur la distinction fondamentale entre \cle{corrélation} et \cle{causalité}.
\end{itemize}

\courssub{Situation}
Monsieur \textsc{Kamga} gère un call-box situé au carrefour \textsc{Deido} à Douala. Pour optimiser sa trésorerie et anticiper ses recettes, il a consigné pendant \textbf{20 jours ouvrés consécutifs} le nombre de clients servis ($X$) et la recette journalière en Francs CFA ($Y$). Les données brutes sont reprises dans les tableaux ci-dessous.

\begin{center}
\begin{tabularx}{\linewidth}{|c|*{10}{X|}}\hline
\textbf{Jour} & 1 & 2 & 3 & 4 & 5 & 6 & 7 & 8 & 9 & 10 \\ \hline
$X$ (Clients) & 55 & 62 & 68 & 70 & 75 & 78 & 80 & 82 & 85 & 88 \\ \hline
$Y$ (FCFA) & \num{27000} & \num{30000} & \num{32000} & \num{33000} & \num{35000} & \num{36000} & \num{37000} & \num{38000} & \num{39000} & \num{40000} \\ \hline
\end{tabularx}
\end{center}
\vspace{0.3cm}
\begin{center}
\begin{tabularx}{\linewidth}{|c|*{10}{X|}}\hline
\textbf{Jour} & 11 & 12 & 13 & 14 & 15 & 16 & 17 & 18 & 19 & 20 \\ \hline
$X$ (Clients) & 90 & 95 & 98 & 102 & 105 & 110 & 115 & 120 & 125 & 130 \\ \hline
$Y$ (FCFA) & \num{41000} & \num{43000} & \num{44000} & \num{46000} & \num{47000} & \num{49000} & \num{51000} & \num{53000} & \num{55000} & \num{57000} \\ \hline
\end{tabularx}
\end{center}

M. \textsc{Kamga} souhaite :
\begin{enumerate}
    \item Avoir une vue synthétique de la relation entre fréquentation et recette.
    \item Savoir s'il peut prévoir la recette de demain s'il attend \textbf{120 clients} (interpolation).
    \item Savoir s'il peut prévoir la recette pour un jour de fête (ex. \textsc{Fête du Travail}) s'il attend \textbf{200 clients} (extrapolation).
\end{enumerate}

\courssub{Consignes}
\begin{enumerate}
    \item \textbf{Tableau à double entrée :} Regrouper les données dans un tableau à double entrée avec \textbf{4 classes pour $X$} (amplitude 20, première classe $[50;70[$) et \textbf{3 classes pour $Y$} (amplitude 10\,000, première classe $[25\,000;35\,000[$). Calculer les \textbf{moyennes conditionnelles} de $Y$ pour chaque classe de $X$. Commenter.
    \item \textbf{Nuage de points :} Tracer le nuage de points $(x_i; y_i)$ sur papier millimétré (ou GeoGebra/calculatrice). Estimer visuellement le coefficient de corrélation $r$. La liaison semble-t-elle linéaire ?
    \item \textbf{Régression linéaire :} À l'aide de la calculatrice (mode \texttt{STAT} $\to$ \texttt{REG} $\to$ \texttt{Lin}), calculer :
    \begin{itemize}
        \item Le coefficient de corrélation linéaire $r$ et le coefficient de détermination $R^2$. Interpréter.
        \item L'équation de la droite de régression $y = ax + b$ (donner $a$ et $b$ arrondis à $10^{-2}$ près).
    \end{itemize}
    Tracer cette droite sur le nuage de points.
    \item \textbf{Validation du modèle (Résidus) :} Calculer les résidus $e_i = y_i - (ax_i + b)$ (liste \texttt{RESID} sur calculatrice). Tracer le nuage des résidus $(x_i; e_i)$. Que constatez-vous ? Le modèle linéaire est-il validé ?
    \item \textbf{Prévisions :}
    \begin{itemize}
        \item Prévision pour $x = 120$ clients (interpolation). Calculer $\hat{y}$.
        \item Prévision pour $x = 200$ clients (extrapolation). Calculer $\hat{y}$.
    \end{itemize}
    \item \textbf{Communication professionnelle :} Rédiger un \textbf{SMS} (maximum 160 caractères) adressé à M. \textsc{Kamga} donnant la prévision pour 120 clients et la mise en garde pour 200 clients.
    \item \textbf{Débat oral (Causalité vs Corrélation) :} \og Peut-on affirmer que l'augmentation du nombre de clients \textbf{CAUSE} l'augmentation de la recette ? \fg Citer au moins deux \textbf{facteurs perturbateurs} (variables confondantes) propres au métier de call-box (ex. : panne réseau, promo opérateur, durée moyenne d'appel).
\end{enumerate}

\courssub{Ressources à mobiliser}
\begin{propriete}[Formules de la régression linéaire (Moindres Carrés)]
Pour une série statistique à deux variables $(x_i; y_i)_{i=1..n}$ de effectif total $N$ :
\begin{itemize}
    \item \textbf{Moyennes :} $\bar{x} = \frac{\sum x_i}{N}$, $\bar{y} = \frac{\sum y_i}{N}$.
    \item \textbf{Variances et Covariance :}
    $s_x^2 = \frac{\sum (x_i - \bar{x})^2}{N} = \frac{\sum x_i^2}{N} - \bar{x}^2$,
    $s_y^2 = \frac{\sum (y_i - \bar{y})^2}{N}$,
    $s_{xy} = \frac{\sum (x_i - \bar{x})(y_i - \bar{y})}{N} = \frac{\sum x_i y_i}{N} - \bar{x}\bar{y}$.
    \item \textbf{Coefficient de corrélation linéaire (Pearson) :} $r = \frac{s_{xy}}{s_x s_y}$ \quad ($-1 \le r \le 1$).
    \item \textbf{Coefficient de détermination :} $R^2 = r^2$ (part de variance de $Y$ expliquée par $X$).
    \item \textbf{Droite de régression $y$ en $x$ :} $y = ax + b$ avec
    \[ a = \frac{s_{xy}}{s_x^2} = r \frac{s_y}{s_x} \qquad \text{et} \qquad b = \bar{y} - a\bar{x}. \]
    Cette droite passe par le point moyen $G(\bar{x}; \bar{y})$.
    \item \textbf{Résidu :} $e_i = y_i - \hat{y}_i = y_i - (ax_i + b)$. $\sum e_i = 0$.
\end{itemize}
\textbf{Critères de validation :} $|r| \ge 0,8$ (liaison forte) ; $R^2 \ge 0,65$ ; nuage des résidus sans structure apparente (aléatoire, centré sur 0).
\end{propriete}

\begin{attention}[Pièges fréquents]
\begin{itemize}
    \item Confondre droite de régression $y$ en $x$ (prévoir $y$ connaissant $x$) et $x$ en $y$.
    \item Oublier que l'extrapolation (prévoir hors intervalle des $x$ observés) est hasardeuse.
    \item Ne pas vérifier les résidus : un $r$ proche de 1 ne garantit pas la linéarité (ex. relation parabolique symétrique).
    \item Écrire $r = 0,95$ au lieu de $r \approx 0,95$ (valeur approchée calculatrice).
    \item Confondre corrélation (lien statistique) et causalité (lien cause-effet).
\end{itemize}
\end{attention}

\courssub{Démarche et corrigé commenté}
\begin{exoresolu}[Corrigé détaillé de l'activité d'intégration]
\textbf{1. Tableau à double entrée et moyennes conditionnelles}\par
Choix des classes :
\begin{itemize}
    \item $X$ (Clients) : $[50;70[$, $[70;90[$, $[90;110[$, $[110;130]$ (amplitude 20, dernière classe fermée pour inclure 130).
    \item $Y$ (Recette) : $[25\,000;35\,000[$, $[35\,000;45\,000[$, $[45\,000;60\,000[$ (amplitude 10\,000, dernière classe élargie pour inclure 57\,000).
\end{itemize}

Effectifs croisés (tableau de contingence) :
\begin{center}
\begin{tabular}{|c|ccc|c|}\hline
$X \backslash Y$ & $[25k;35k[$ & $[35k;45k[$ & $[45k;60k[$ & \textbf{Total} \\ \hline
$[50;70[$ & 3 & 0 & 0 & 3 \\
$[70;90[$ & 1 & 6 & 0 & 7 \\
$[90;110[$ & 0 & 3 & 2 & 5 \\
$[110;130]$ & 0 & 0 & 5 & 5 \\ \hline
\textbf{Total} & 4 & 9 & 7 & 20 \\ \hline
\end{tabular}
\end{center}

\textbf{Moyennes conditionnelles de $Y$ sachant la classe de $X$ :}
\begin{align*}
[50;70[ &: \quad \bar{y}_{[50;70[} = \frac{27\,000 + 30\,000 + 32\,000}{3} = \frac{89\,000}{3} \approx \textbf{29\,667 FCFA} \\
[70;90[ &: \quad \bar{y}_{[70;90[} = \frac{33\,000 + 35\,000 + 36\,000 + 37\,000 + 38\,000 + 39\,000 + 40\,000}{7} = \frac{258\,000}{7} \approx \textbf{36\,857 FCFA} \\
[90;110[ &: \quad \bar{y}_{[90;110[} = \frac{41\,000 + 43\,000 + 44\,000 + 46\,000 + 47\,000}{5} = \frac{221\,000}{5} = \textbf{44\,200 FCFA} \\
[110;130] &: \quad \bar{y}_{[110;130]} = \frac{49\,000 + 51\,000 + 53\,000 + 55\,000 + 57\,000}{5} = \frac{265\,000}{5} = \textbf{53\,000 FCFA}
\end{align*}
\textbf{Commentaire :} Les moyennes conditionnelles augmentent régulièrement avec les classes de $X$ (29\,667 $\to$ 36\,857 $\to$ 44\,200 $\to$ 53\,000). Cela confirme une liaison positive forte entre fréquentation et recette. Les points moyens $(\bar{x}_{\text{classe}}; \bar{y}_{\text{classe}})$ sont quasi alignés.

\textbf{2. Nuage de points et estimation visuelle de $r$}\par
Le nuage de points (voir figure ci-dessous) montre des points très groupés autour d'une droite montante. Visuellement, la liaison est \textbf{linéaire, positive et très forte}. On estime $r \approx 0,99$.

\begin{popfigure}
\begin{tikzpicture}[scale=0.65]
\begin{axis}[
    width=\linewidth, height=8cm,
    xlabel={Nombre de clients $X$}, ylabel={Recette $Y$ (FCFA)},
    xmin=50, xmax=135, ymin=20000, ymax=60000,
    xtick={50,70,90,110,130}, ytick={20000,30000,40000,50000,60000},
    yticklabel style={/pgf/number format/fixed, /pgf/number format/precision=0, /pgf/number format/set thousands separator={\,}},
    grid=major, grid style={popline}, minor tick num=1,
    axis lines=middle, enlargelimits=0.05,
    legend style={at={(0.02,0.98)}, anchor=north west, fill=white, draw=popmuted}
]
% Data points
\addplot[only marks, mark=*, mark size=2.5, popblue] coordinates {
    (55,27000) (62,30000) (68,32000) (70,33000) (75,35000)
    (78,36000) (80,37000) (82,38000) (85,39000) (88,40000)
    (90,41000) (95,43000) (98,44000) (102,46000) (105,47000)
    (110,49000) (115,51000) (120,53000) (125,55000) (130,57000)
};
\addlegendentry{Données observées}

% Regression line y = 400x + 5000
\addplot[poporange, thick, domain=50:135, samples=2] {400*x + 5000};
\addlegendentry{Régression $y=400x+5000$}

% Mean point G
\addplot[only marks, mark=*, mark size=4, popgreen] coordinates {(90,41000)};
\addlegendentry{Point moyen $G(90;41000)$}
\end{axis}
\end{tikzpicture}
\legende{Nuage de points et droite de régression des moindres carrés.}
\end{popfigure}

\textbf{3. Calculs sur calculatrice (Mode STAT $\to$ REG $\to$ Lin)}\par
Saisie des listes $L1$ ($X$) et $L2$ ($Y$).
Résultats exacts (données quasi parfaites $Y = 400X + 5000$ avec micro-arrondis) :
\begin{itemize}
    \item $\bar{x} = 90$ \quad $\bar{y} = 41\,000$
    \item $s_x \approx 20,73$ \quad $s_y \approx 8\,295,6$
    \item $s_{xy} \approx 171\,975$
    \item \textbf{Coefficient de corrélation :} $r = \frac{s_{xy}}{s_x s_y} \approx \mathbf{0,9997}$
    \item \textbf{Coefficient de détermination :} $R^2 = r^2 \approx \mathbf{0,9994}$
    \item \textbf{Pente :} $a = \frac{s_{xy}}{s_x^2} \approx \mathbf{400,00}$ (arrondi au centime)
    \item \textbf{Ordonnée à l'origine :} $b = \bar{y} - a\bar{x} = 41\,000 - 400 \times 90 = \mathbf{5\,000}$
\end{itemize}
\textbf{Équation de la droite de régression :} $\boxed{\hat{y} = 400,00\,x + 5\,000}$.
\textbf{Interprétation :}
\begin{itemize}
    \item $r \approx 0,9997 \approx 1$ : liaison linéaire positive \textbf{quasi parfaite}.
    \item $R^2 \approx 99,94\%$ : la variable $X$ (nombre de clients) explique \textbf{99,94\%} de la variabilité de la recette $Y$. Le modèle est excellent.
    \item $a = 400$ : chaque client supplémentaire apporte en moyenne \textbf{400 FCFA} de recette (prix moyen unitaire effectif).
    \item $b = 5\,000$ : recette théorique pour 0 client (coûts fixes, abonnements, pourboires...). Attention : extrapolation hors champ ($x=0$ non observé).
\end{itemize}

\textbf{4. Analyse des résidus (Validation du modèle)}\par
Calcul des résidus $e_i = y_i - (400 x_i + 5000)$ :
\begin{center}
\begin{tabular}{|c|c|c|c|c|}\hline
$i$ & $x_i$ & $y_i$ & $\hat{y}_i$ & $e_i$ \\ \hline
1 & 55 & \num{27000} & \num{27000} & 0 \\
2 & 62 & \num{30000} & \num{29800} & +200 \\
3 & 68 & \num{32000} & \num{32200} & -200 \\
4 & 70 & \num{33000} & \num{33000} & 0 \\
5 & 75 & \num{35000} & \num{35000} & 0 \\
6 & 78 & \num{36000} & \num{36200} & -200 \\
7 & 80 & \num{37000} & \num{37000} & 0 \\
8 & 82 & \num{38000} & \num{37800} & +200 \\
9 & 85 & \num{39000} & \num{39000} & 0 \\
10 & 88 & \num{40000} & \num{40200} & -200 \\ \hline
11 & 90 & \num{41000} & \num{41000} & 0 \\
12 & 95 & \num{43000} & \num{43000} & 0 \\
13 & 98 & \num{44000} & \num{44200} & -200 \\
14 & 102 & \num{46000} & \num{45800} & +200 \\
15 & 105 & \num{47000} & \num{47000} & 0 \\
16 & 110 & \num{49000} & \num{49000} & 0 \\
17 & 115 & \num{51000} & \num{51000} & 0 \\
18 & 120 & \num{53000} & \num{53000} & 0 \\
19 & 125 & \num{55000} & \num{55000} & 0 \\
20 & 130 & \num{57000} & \num{57000} & 0 \\ \hline
\end{tabular}
\end{center}

\begin{popfigure}
\begin{tikzpicture}[scale=0.7]
\begin{axis}[
    width=\linewidth, height=6cm,
    xlabel={$X$ (Clients)}, ylabel={Résidus $e_i$ (FCFA)},
    xmin=50, xmax=135, ymin=-500, ymax=500,
    xtick={50,70,90,110,130}, ytick={-400,-200,0,200,400},
    grid=major, grid style={popline}, axis lines=middle,
    legend style={at={(0.98,0.98)}, anchor=north east, fill=white, draw=popmuted}
]
\addplot[only marks, mark=*, mark size=2.5, poppurple] coordinates {
    (55,0) (62,200) (68,-200) (70,0) (75,0) (78,-200) (80,0) (82,200) (85,0) (88,-200)
    (90,0) (95,0) (98,-200) (102,200) (105,0) (110,0) (115,0) (120,0) (125,0) (130,0)
};
\addlegendentry{Résidus}
\addplot[popred, dashed, thick, domain=50:135] {0};
\addlegendentry{Zéro résidu}
\end{axis}
\end{tikzpicture}
\legende{Nuage des résidus : dispersion aléatoire autour de 0, amplitude faible ($\pm 200$ FCFA).}
\end{popfigure}

\textbf{Analyse :} Les résidus sont \textbf{centrés sur 0} (somme nulle), de \textbf{faible amplitude} ($\pm 200$ FCFA soit $\pm 0,4\%$ de la recette moyenne) et \textbf{sans structure systématique} (pas de courbe, pas d'entonnoir). Le modèle linéaire est \textbf{parfaitement validé}.

\textbf{5. Prévisions}\par
\begin{itemize}
    \item \textbf{Interpolation ($x_0 = 120$ clients) :} $120 \in [55; 130]$ (intervalle des observations).
    $\hat{y}(120) = 400 \times 120 + 5\,000 = \textbf{53\,000 FCFA}$.
    \textit{Fiabilité : TRÈS ÉLEVÉE (interpolation, $R^2 \approx 1$).}
    \item \textbf{Extrapolation ($x_0 = 200$ clients) :} $200 \notin [55; 130]$ (bien au-delà).
    $\hat{y}(200) = 400 \times 200 + 5\,000 = \textbf{85\,000 FCFA}$.
    \textit{Fiabilité : FAIBLE. Risque de non-linéarité (saturation de la capacité, pannes, tarifs dégressifs, horaires d'ouverture).}
\end{itemize}

\textbf{6. SMS professionnel (160 caractères max)}\par
\begin{center}
\begin{tabular}{|p{0.9\linewidth}|}\hline
\texttt{M. Kamga, prev 120 clients: 53 000 FCFA (fiable). Attention 200 clients: 85 000 FCFA extrapole risqué (saturation/reseau). Cordialement.} \\
\hline
\end{tabular}
\end{center}
\textit{Comptage : 158 caractères. Contient : prévision chiffrée interpolation, alerte extrapolation, cause technique.}

\textbf{7. Débat : Corrélation vs Causalité}\par
\textbf{Non, on ne peut pas affirmer que plus de clients \textbf{CAUSE} plus de recette au sens strict de causalité univoque.}
La corrélation ($r \approx 1$) montre un lien statistique fort, mais des \textbf{facteurs perturbateurs (variables confondantes)} influencent $Y$ pour un $X$ donné :
\begin{enumerate}
    \item \textbf{Durée moyenne des appels :} 120 clients de 2 min rapportent moins que 100 clients de 5 min (facturation à la minute).
    \item \textbf{État du réseau / Pannes opérateurs :} Une panne MTN/Orange réduit la recette par client (échecs d'appel, raccourcis).
    \item \textbf{Promotions opérateurs (ex. "Appel gratuit 20h-22h") :} Augmente $X$ mais baisse drastiquement le prix unitaire $\to$ $Y$ n'augmente pas proportionnellement.
    \item \textbf{Capacité physique du call-box :} Au-delà d'un seuil (ex. 150 clients/jour), la file d'attente fait fuir des clients $\to$ relation non linéaire (plafonnement).
\end{enumerate}
\cle{Conclusion} : La régression est un outil de \textbf{prévision statistique}, pas de \textbf{détermination causale}. Le gérant doit croiser ces données avec la qualité de service et les offres commerciales.
\end{exoresolu}

\courssub{Bilan de l'activité}
Cette activité a permis de mettre en évidence, dans un contexte métier authentique (call-box Douala), la \textbf{chaîne complète de l'analyse bivariée} :
\begin{enumerate}
    \item \textbf{Organisation :} Le tableau à double entrée synthétise 200 données brutes en 12 cases lisibles ; les moyennes conditionnelles révèlent la tendance sans calculateur.
    \item \textbf{Visualisation :} Le nuage de points confirme l'hypothèse linéaire avant tout calcul.
    \item \textbf{Modélisation :} La calculatrice fournit $r$, $R^2$, $a$, $b$ instantanément. L'interprétation économique ($a$ = panier moyen, $b$ = fixe) donne du sens aux paramètres.
    \item \textbf{Validation critique :} L'analyse des résidus (souvent oubliée) est \textbf{indispensable} pour légitimer l'usage de la droite. Ici, résidus aléatoires et minuscules $\to$ modèle validé.
    \item \textbf{Prévision raisonnée :} Distinction nette \textbf{interpolation} (fiable, dans l'intervalle des données) vs \textbf{extrapolation} (hasardeuse, hors intervalle). Le chiffre 85\,000 FCFA pour 200 clients est un \textbf{repère théorique}, pas une certitude.
    \item \textbf{Communication technique :} Le SMS force la synthèse : chiffre clé + niveau de confiance + alerte risque.
    \item \textbf{Esprit critique :} Le débat final ancre la distinction \cle{Corrélation $\neq$ Causalité}, compétence transversale majeure (sciences, économie, vie citoyenne).
\end{enumerate}
L'élève termine cette leçon capable de \textbf{traiter un dossier statistique réel de A à Z} : données brutes $\to$ tableau $\to$ graphique $\to$ modèle $\to$ validation $\to$ décision $\to$ communication.

\courssub{Complément : Ajustement non linéaire par linéarisation}
\begin{savaistu}[Culture Terminale -- Hors programme strict]
Dans de nombreux métiers techniques, la liaison entre deux variables n'est pas linéaire (ex. : loi de décharge d'un condensateur $u = U_0 e^{-t/RC}$, loi de puissance d'une résistance $P = RI^2$, croissance bactérienne). On peut parfois \textbf{linéariser} la relation par un changement de variable pour utiliser les outils de la régression linéaire.

\textbf{Exemples classiques :}
\begin{itemize}
    \item \textbf{Modèle exponentiel :} $y = a e^{bx}$ $\xrightarrow{\ln}$ $\ln(y) = \ln(a) + bx$ (affine en $x$). On régresse $\ln(y)$ sur $x$.
    \item \textbf{Modèle puissance :} $y = a x^b$ $\xrightarrow{\ln}$ $\ln(y) = \ln(a) + b \ln(x)$ (affine en $\ln(x)$). On régresse $\ln(y)$ sur $\ln(x)$.
    \item \textbf{Modèle inverse :} $y = \frac{a}{x} + b$ $\xrightarrow{X=1/x}$ $y = a X + b$ (affine en $X=1/x$).
\end{itemize}

\textbf{Méthode :}
\begin{enumerate}
    \item Identifier le type de liaison suspecté (nuage de points, connaissance physique).
    \item Transformer les variables ($Y' = \ln Y$, $X' = \ln X$, etc.).
    \item Effectuer la régression linéaire sur les variables transformées.
    \item Revenir au modèle initial et \textbf{valider sur les résidus du modèle initial} (pas sur les résidus du modèle linéarisé !).
\end{enumerate}

\cle{Attention :} Le coefficient de détermination $R^2$ obtenu sur le modèle linéarisé ne s'interprète pas directement sur le modèle initial. La validation par les résidus du modèle non linéaire est impérative.
\end{savaistu}

\newpage

\coursec{Méthodes et exercices résolus}

\coursec{Statistiques à deux variables}

\courssub{Construire et exploiter un tableau à double entrée}

\begin{methode}[Lire et compléter un tableau à double entrée]
Pour organiser les données d'une série statistique à deux variables :
\begin{enumerate}
\item Identifier les deux caractères étudiés (exemple : nombre d'heures de formation et nombre de pièces produites par ouvrier).
\item Placer les modalités du premier caractère en ligne, celles du second en colonne.
\item Remplir les cases intérieures avec les effectifs conjoints $n_{ij}$.
\item Ajouter une colonne et une ligne « Total » : les effectifs marginaux (totaux par ligne et par colonne).
\item Vérifier la cohérence : la somme des totaux en ligne = somme des totaux en colonne = effectif total $N$.
\item Pour répondre à une question, repérer la case, la ligne ou la colonne concernée et calculer la fréquence : $f = \dfrac{\text{effectif}}{\text{effectif total}}$.
\end{enumerate}
\textbf{Réflexe :} toujours vérifier les totaux avant de conclure ; une erreur de total fausse toutes les fréquences.
\end{methode}

\begin{exoresolu}[Production d'un atelier de menuiserie à Douala]
Dans un atelier employant 50 ouvriers, on relève le nombre d'années d'expérience $X$ et le nombre de pièces produites par jour $Y$. Les résultats sont :
\begin{itemize}
\item 12 ouvriers ont moins de 5 ans d'expérience et produisent moins de 10 pièces ;
\item 8 ouvriers ont moins de 5 ans d'expérience et produisent 10 pièces ou plus ;
\item 18 ouvriers ont 5 ans ou plus d'expérience et produisent 10 pièces ou plus ;
\item les autres ont 5 ans ou plus d'expérience et produisent moins de 10 pièces.
\end{itemize}
\begin{enumerate}
\item Construire le tableau à double entrée complet avec les totaux.
\item Quelle proportion d'ouvriers a 5 ans ou plus d'expérience ?
\item Parmi les ouvriers expérimentés (5 ans ou plus), quelle proportion produit 10 pièces ou plus ?
\end{enumerate}
\tcblower
\textbf{1. Construction du tableau.} L'effectif total est $N = 50$. La case manquante vaut :
\[ 50 - (12 + 8 + 18) = 50 - 38 = 12. \]
\begin{center}
\begin{tabularx}{\linewidth}{|l|X|X|X|}
\hline
\rowcolor{popblueL} & $Y < 10$ pièces & $Y \geq 10$ pièces & Total \\
\hline
$X < 5$ ans & 12 & 8 & 20 \\
\hline
$X \geq 5$ ans & 12 & 18 & 30 \\
\hline
Total & 24 & 26 & 50 \\
\hline
\end{tabularx}
\end{center}
Vérification : $20 + 30 = 50$ et $24 + 26 = 50$. Le tableau est cohérent.

\textbf{2. Proportion d'ouvriers expérimentés.} On lit le total de la ligne « $X \geq 5$ ans » : 30 ouvriers.
\[ f = \frac{30}{50} = 0{,}6. \]
\textbf{Conclusion :} 60 \% des ouvriers ont 5 ans ou plus d'expérience.

\textbf{3. Proportion conditionnelle.} On se restreint à la ligne des 30 ouvriers expérimentés ; parmi eux, 18 produisent 10 pièces ou plus :
\[ f = \frac{18}{30} = 0{,}6. \]
\textbf{Conclusion :} parmi les ouvriers expérimentés, 60 \% produisent au moins 10 pièces par jour.
\end{exoresolu}

\courssub{Représenter un nuage de points et le point moyen}

\begin{methode}[Tracer un nuage de points et placer le point moyen $G$]
\begin{enumerate}
\item Choisir deux axes gradués adaptés aux valeurs de $x$ et de $y$ (on peut commencer les graduations ailleurs qu'en 0).
\item Placer chaque point $M_i(x_i\,;\,y_i)$ du tableau.
\item Calculer les coordonnées du \cle{point moyen} :
\[ \bar{x} = \frac{x_1 + x_2 + \cdots + x_n}{n} \qquad \text{et} \qquad \bar{y} = \frac{y_1 + y_2 + \cdots + y_n}{n}. \]
\item Placer $G(\bar{x}\,;\,\bar{y})$ d'une couleur différente et le nommer.
\item Observer la forme du nuage : s'il est allongé autour d'une droite, un ajustement linéaire est envisageable.
\end{enumerate}
\textbf{Réflexe :} au Probatoire, le point moyen $G$ est presque toujours demandé : le calculer avec soin et le placer sur le graphique.
\end{methode}

\begin{exoresolu}[Dépenses publicitaires et chiffre d'affaires]
Une entreprise de Yaoundé relève, sur 6 mois, ses dépenses publicitaires $x_i$ (en centaines de milliers de FCFA) et son chiffre d'affaires $y_i$ (en millions de FCFA) :
\begin{center}
\begin{tabular}{|c|cccccc|}
\hline
\rowcolor{popblueL} $x_i$ & 1 & 2 & 3 & 4 & 5 & 6 \\
\hline
$y_i$ & 12 & 15 & 17 & 20 & 24 & 26 \\
\hline
\end{tabular}
\end{center}
\begin{enumerate}
\item Calculer les coordonnées du point moyen $G$.
\item Représenter le nuage de points et placer $G$.
\item Un ajustement linéaire semble-t-il justifié ?
\end{enumerate}
\tcblower
\textbf{1. Point moyen.} On a $n = 6$ points.
\[ \bar{x} = \frac{1+2+3+4+5+6}{6} = \frac{21}{6} = 3{,}5. \]
\[ \bar{y} = \frac{12+15+17+20+24+26}{6} = \frac{114}{6} = 19. \]
Donc $G(3{,}5\,;\,19)$.

\textbf{2. Nuage de points.}
\begin{popfigure}
\begin{tikzpicture}
\begin{axis}[width=0.85\linewidth, height=6cm, xlabel={$x$ (dépenses)}, ylabel={$y$ (chiffre d'affaires)}, xmin=0, xmax=7, ymin=8, ymax=30, grid=both, grid style={popline!40}, xtick={1,2,3,4,5,6}]
\addplot[only marks, mark=*, mark size=2.5pt, popblue] coordinates {(1,12) (2,15) (3,17) (4,20) (5,24) (6,26)};
\addplot[only marks, mark=star, mark size=4pt, poporange] coordinates {(3.5,19)};
\node[poporange, anchor=south west] at (axis cs:3.6,19) {$G$};
\end{axis}
\end{tikzpicture}
\legende{Nuage de points et point moyen $G(3{,}5\,;\,19)$}
\end{popfigure}

\textbf{3. Justification de l'ajustement.} Les points du nuage sont sensiblement alignés : quand $x$ augmente, $y$ augmente de façon quasi régulière. Le nuage est allongé autour d'une droite.
\textbf{Conclusion :} un ajustement linéaire est justifié.
\end{exoresolu}

\courssub{Calculer le coefficient de corrélation linéaire}

\begin{methode}[Calculer et interpréter le coefficient $r$]
\begin{enumerate}
\item Calculer les moyennes $\bar{x}$ et $\bar{y}$.
\item Calculer la \cle{covariance} :
\[ \operatorname{cov}(x,y) = \frac{1}{n}\sum_{i=1}^{n} x_i y_i - \bar{x}\,\bar{y}. \]
\item Calculer les \cle{variances} :
\[ V(x) = \frac{1}{n}\sum_{i=1}^{n} x_i^2 - \bar{x}^2 \qquad \text{et} \qquad V(y) = \frac{1}{n}\sum_{i=1}^{n} y_i^2 - \bar{y}^2. \]
\item En déduire le \cle{coefficient de corrélation linéaire} :
\[ r = \frac{\operatorname{cov}(x,y)}{\sqrt{V(x)\,V(y)}}. \]
\item Interpréter : $-1 \leq r \leq 1$. Plus $|r|$ est proche de 1, plus la corrélation linéaire est forte. En pratique, on considère l'ajustement linéaire pertinent si $|r| \geq 0{,}75$ environ. Si $r > 0$ la corrélation est positive, si $r < 0$ elle est négative.
\end{enumerate}
\textbf{Réflexe :} présenter les calculs dans un tableau à colonnes $x_i$, $y_i$, $x_i y_i$, $x_i^2$, $y_i^2$ avec une ligne de totaux.
\end{methode}

\begin{exoresolu}[Coefficient de corrélation de la série publicité / chiffre d'affaires]
On reprend la série de l'exercice précédent :
\begin{center}
\begin{tabular}{|c|cccccc|}
\hline
\rowcolor{popblueL} $x_i$ & 1 & 2 & 3 & 4 & 5 & 6 \\
\hline
$y_i$ & 12 & 15 & 17 & 20 & 24 & 26 \\
\hline
\end{tabular}
\end{center}
\begin{enumerate}
\item Compléter le tableau de calcul et déterminer $\operatorname{cov}(x,y)$, $V(x)$ et $V(y)$.
\item Calculer le coefficient de corrélation $r$ (arrondi à $10^{-3}$).
\item L'ajustement linéaire est-il pertinent ?
\end{enumerate}
\tcblower
\textbf{1. Tableau de calcul.} On rappelle $\bar{x} = 3{,}5$ et $\bar{y} = 19$.
\begin{center}
\begin{tabular}{|c|c|c|c|c|}
\hline
\rowcolor{popblueL} $x_i$ & $y_i$ & $x_i y_i$ & $x_i^2$ & $y_i^2$ \\
\hline
1 & 12 & 12 & 1 & 144 \\
2 & 15 & 30 & 4 & 225 \\
3 & 17 & 51 & 9 & 289 \\
4 & 20 & 80 & 16 & 400 \\
5 & 24 & 120 & 25 & 576 \\
6 & 26 & 156 & 36 & 676 \\
\hline
\textbf{Totaux} & & \textbf{449} & \textbf{91} & \textbf{2310} \\
\hline
\end{tabular}
\end{center}
Covariance :
\[ \operatorname{cov}(x,y) = \frac{449}{6} - 3{,}5 \times 19 = 74{,}833\ldots - 66{,}5 = \frac{25}{3} \approx 8{,}333. \]
Variances :
\[ V(x) = \frac{91}{6} - (3{,}5)^2 = 15{,}166\ldots - 12{,}25 = \frac{35}{12} \approx 2{,}917. \]
\[ V(y) = \frac{2310}{6} - 19^2 = 385 - 361 = 24. \]

\textbf{2. Coefficient de corrélation.}
\[ r = \frac{8{,}333}{\sqrt{2{,}917 \times 24}} = \frac{8{,}333}{\sqrt{70}} = \frac{8{,}333}{8{,}367} \approx 0{,}996. \]

\textbf{3. Pertinence.} On a $r \approx 0{,}996$, très proche de 1.
\textbf{Conclusion :} la corrélation linéaire est très forte et positive ; l'ajustement linéaire est pleinement justifié.
\end{exoresolu}

\courssub{Déterminer la droite de régression des moindres carrés}

\begin{methode}[Déterminer l'équation de la droite de régression de $y$ en $x$]
\begin{enumerate}
\item Calculer $\bar{x}$, $\bar{y}$, $\operatorname{cov}(x,y)$ et $V(x)$ (comme pour le coefficient $r$).
\item Le coefficient directeur de la droite des moindres carrés est :
\[ a = \frac{\operatorname{cov}(x,y)}{V(x)}. \]
\item La droite passe par le point moyen $G(\bar{x}\,;\,\bar{y})$, donc :
\[ b = \bar{y} - a\,\bar{x}. \]
\item Écrire l'équation : $y = ax + b$.
\item Pour tracer la droite : placer $G$ et un second point calculé avec l'équation (par exemple pour la plus petite et la plus grande valeur de $x$).
\end{enumerate}
\begin{aretenir}[Propriété fondamentale]
La droite de régression de $y$ en $x$ passe toujours par le point moyen $G(\bar{x}\,;\,\bar{y})$. Cette propriété sert à calculer $b$ et à vérifier le tracé.
\end{aretenir}
\end{methode}

\begin{exoresolu}[Droite de régression publicité / chiffre d'affaires]
On reprend la même série avec les résultats déjà établis : $\bar{x} = 3{,}5$, $\bar{y} = 19$, $\operatorname{cov}(x,y) = \dfrac{25}{3}$, $V(x) = \dfrac{35}{12}$.
\begin{enumerate}
\item Déterminer l'équation de la droite de régression de $y$ en $x$ (coefficients arrondis à $10^{-2}$).
\item Tracer cette droite sur le nuage de points.
\end{enumerate}
\tcblower
\textbf{1. Équation de la droite.} Coefficient directeur :
\[ a = \frac{\operatorname{cov}(x,y)}{V(x)} = \frac{\dfrac{25}{3}}{\dfrac{35}{12}} = \frac{25}{3} \times \frac{12}{35} = \frac{300}{105} = \frac{20}{7} \approx 2{,}86. \]
Ordonnée à l'origine, en utilisant le passage par $G$ :
\[ b = \bar{y} - a\bar{x} = 19 - \frac{20}{7} \times 3{,}5 = 19 - 10 = 9. \]
\textbf{Conclusion :} la droite de régression a pour équation $y = 2{,}86x + 9$ (valeur exacte : $y = \dfrac{20}{7}x + 9$).

\textbf{2. Tracé.} La droite passe par $G(3{,}5\,;\,19)$. Pour $x = 1$ : $y = \dfrac{20}{7} + 9 \approx 11{,}86$. Pour $x = 6$ : $y = \dfrac{120}{7} + 9 \approx 26{,}14$.
\begin{popfigure}
\begin{tikzpicture}
\begin{axis}[width=0.85\linewidth, height=6cm, xlabel={$x$}, ylabel={$y$}, xmin=0, xmax=7, ymin=8, ymax=30, grid=both, grid style={popline!40}, xtick={1,2,3,4,5,6}]
\addplot[only marks, mark=*, mark size=2.5pt, popblue] coordinates {(1,12) (2,15) (3,17) (4,20) (5,24) (6,26)};
\addplot[popcurve, domain=0.5:6.5] {20/7*x + 9};
\addplot[only marks, mark=star, mark size=4pt, poporange] coordinates {(3.5,19)};
\node[poporange, anchor=south west] at (axis cs:3.6,19) {$G$};
\end{axis}
\end{tikzpicture}
\legende{Droite de régression $y = 2{,}86x + 9$ passant par $G$}
\end{popfigure}
La droite passe bien par $G$ et suit l'allure du nuage : le tracé est cohérent.
\end{exoresolu}

\courssub{Valider le modèle et effectuer une prévision}

\begin{methode}[Estimer une valeur par interpolation ou extrapolation]
\begin{enumerate}
\item Vérifier d'abord que l'ajustement linéaire est justifié (nuage allongé, $|r|$ proche de 1).
\item Pour estimer $y$ à partir d'une valeur $x_0$, remplacer $x$ par $x_0$ dans l'équation : $y_0 = a x_0 + b$.
\item Préciser la nature de la prévision : \cle{interpolation} si $x_0$ est à l'intérieur de la plage des données (fiable), \cle{extrapolation} si $x_0$ est en dehors (à utiliser avec prudence).
\item Rédiger la conclusion avec les unités de la situation.
\end{enumerate}
\textbf{Réflexe :} une prévision n'a de sens que si le modèle a été validé ; toujours citer la justification dans la réponse.
\end{methode}

\begin{exoresolu}[Prévision de chiffre d'affaires]
On utilise le modèle $y = \dfrac{20}{7}x + 9$ obtenu précédemment ($x$ en centaines de milliers de FCFA, $y$ en millions de FCFA).
\begin{enumerate}
\item Estimer le chiffre d'affaires pour des dépenses publicitaires de \num{450000} FCFA.
\item Estimer le chiffre d'affaires pour des dépenses de \num{1000000} FCFA. Commenter la fiabilité de chaque prévision.
\end{enumerate}
\tcblower
\textbf{1. Dépenses de \num{450000} FCFA.} Cela correspond à $x_0 = 4{,}5$ (centaines de milliers).
\[ y_0 = \frac{20}{7} \times 4{,}5 + 9 = \frac{90}{7} + 9 = \frac{90 + 63}{7} = \frac{153}{7} \approx 21{,}86. \]
\textbf{Conclusion :} le chiffre d'affaires estimé est d'environ 21,86 millions de FCFA.

\textbf{2. Dépenses de \num{1000000} FCFA.} Cela correspond à $x_0 = 10$.
\[ y_0 = \frac{20}{7} \times 10 + 9 = \frac{200}{7} + 9 = \frac{263}{7} \approx 37{,}57. \]
Le chiffre d'affaires estimé est d'environ 37,57 millions de FCFA.

\textbf{Commentaire de fiabilité :}
\begin{itemize}
\item Pour $x_0 = 4{,}5$ : la valeur est à l'intérieur de la plage observée $[1\,;\,6]$ ; c'est une \cle{interpolation}, la prévision est fiable car $r \approx 0{,}996$.
\item Pour $x_0 = 10$ : la valeur est en dehors de la plage observée ; c'est une \cle{extrapolation}. Rien ne garantit que la tendance linéaire se poursuive (marché saturé, concurrence...) : la prévision doit être utilisée avec prudence.
\end{itemize}
\end{exoresolu}

\courssub{Appliquer les statistiques à une situation concrète complète}

\begin{methode}[Conduire une étude complète type Probatoire]
Face à un problème complet de statistique à deux variables :
\begin{enumerate}
\item Lire l'énoncé et identifier les variables, les unités et l'effectif $n$.
\item Représenter le nuage de points et calculer le point moyen $G$.
\item Justifier (ou non) l'ajustement linéaire : allure du nuage et/ou coefficient $r$.
\item Déterminer l'équation de la droite de régression par les moindres carrés.
\item Tracer la droite (en vérifiant le passage par $G$).
\item Répondre aux questions de prévision en rédigeant des conclusions avec les unités et en distinguant interpolation et extrapolation.
\end{enumerate}
\textbf{Réflexe rédaction :} chaque question appelle une phrase de conclusion ; les formules utilisées doivent être écrites avant d'être appliquées numériquement.
\end{methode}

\begin{exoresolu}[Sujet type Probatoire : consommation d'un groupe électrogène]
Un atelier de soudure à Bafoussam relève la puissance utilisée $x_i$ (en kW) et la consommation de carburant $y_i$ (en litres par heure) de son groupe électrogène :
\begin{center}
\begin{tabular}{|c|ccccc|}
\hline
\rowcolor{popblueL} $x_i$ (kW) & 2 & 4 & 5 & 7 & 9 \\
\hline
$y_i$ (L/h) & 1{,}5 & 2{,}3 & 2{,}8 & 3{,}7 & 4{,}6 \\
\hline
\end{tabular}
\end{center}
On donne : $\sum x_i y_i = 94{,}1$ ; $\sum x_i^2 = 175$ ; $\sum y_i^2 = 53{,}23$.
\begin{enumerate}
\item Calculer les coordonnées du point moyen $G$.
\item Calculer le coefficient de corrélation $r$ (arrondi à $10^{-3}$). Un ajustement linéaire est-il justifié ?
\item Déterminer l'équation de la droite de régression de $y$ en $x$.
\item Estimer la consommation pour une puissance de 6 kW, puis pour 15 kW. Commenter.
\end{enumerate}
\tcblower
\textbf{1. Point moyen.} On a $n = 5$.
\[ \bar{x} = \frac{2+4+5+7+9}{5} = \frac{27}{5} = 5{,}4 \qquad ; \qquad \bar{y} = \frac{1{,}5+2{,}3+2{,}8+3{,}7+4{,}6}{5} = \frac{14{,}9}{5} = 2{,}98. \]
Donc $G(5{,}4\,;\,2{,}98)$.

\textbf{2. Coefficient de corrélation.}
\[ \operatorname{cov}(x,y) = \frac{94{,}1}{5} - 5{,}4 \times 2{,}98 = 18{,}82 - 16{,}092 = 2{,}728. \]
\[ V(x) = \frac{175}{5} - (5{,}4)^2 = 35 - 29{,}16 = 5{,}84. \]
\[ V(y) = \frac{53{,}23}{5} - (2{,}98)^2 = 10{,}646 - 8{,}8804 = 1{,}7656. \]
\[ r = \frac{2{,}728}{\sqrt{5{,}84 \times 1{,}7656}} = \frac{2{,}728}{\sqrt{10{,}311\ldots}} = \frac{2{,}728}{3{,}211} \approx 0{,}850. \]
On a $r \approx 0{,}850 > 0{,}75$.
\textbf{Conclusion :} la corrélation linéaire est forte et positive ; l'ajustement linéaire est justifié.

\textbf{3. Droite de régression.}
\[ a = \frac{\operatorname{cov}(x,y)}{V(x)} = \frac{2{,}728}{5{,}84} \approx 0{,}467. \]
\[ b = \bar{y} - a\bar{x} = 2{,}98 - 0{,}467 \times 5{,}4 = 2{,}98 - 2{,}522 = 0{,}458. \]
\textbf{Conclusion :} la droite de régression a pour équation $y = 0{,}467x + 0{,}458$.

\textbf{4. Prévisions.}
Pour $x_0 = 6$ kW :
\[ y_0 = 0{,}467 \times 6 + 0{,}458 = 2{,}802 + 0{,}458 = 3{,}26 \text{ L/h}. \]
C'est une interpolation ($6 \in [2\,;\,9]$) : prévision fiable. La consommation estimée est d'environ 3,26 L/h.

Pour $x_0 = 15$ kW :
\[ y_0 = 0{,}467 \times 15 + 0{,}458 = 7{,}005 + 0{,}458 \approx 7{,}46 \text{ L/h}. \]
C'est une extrapolation (15 est hors de la plage observée) : le résultat, environ 7,46 L/h, est une estimation à prendre avec prudence, d'autant que le groupe peut ne pas supporter une telle puissance.
\end{exoresolu}

\begin{attention}[Les 5 erreurs qui coûtent des points au Probatoire]
\begin{enumerate}
\item \textbf{Confondre les formules de covariance et de variance.} La covariance utilise $\sum x_i y_i$ et le produit $\bar{x}\,\bar{y}$ ; la variance utilise $\sum x_i^2$ et $\bar{x}^2$. Mélanger les deux fausse $r$ et $a$.
\item \textbf{Oublier que la droite de régression passe par $G$.} Calculer $b$ autrement que par $b = \bar{y} - a\bar{x}$, ou tracer une droite qui ne passe pas par $G$, est sanctionné. Vérifier toujours le passage par $G$ sur le graphique.
\item \textbf{Interpréter $r$ sans le comparer à 1.} Dire « $r = 0{,}6$ donc l'ajustement est bon » est faux : il faut justifier en précisant que $|r|$ doit être proche de 1 (en pratique $|r| \geq 0{,}75$).
\item \textbf{Se tromper dans les unités lors des prévisions.} Si $x$ est en centaines de milliers de FCFA, une dépense de \num{450000} FCFA correspond à $x = 4{,}5$ et non 450. Convertir avant de remplacer dans l'équation.
\item \textbf{Présenter une extrapolation comme une certitude.} Toute prévision hors de la plage des données doit être accompagnée d'une réserve ; oublier ce commentaire fait perdre le point de justification. De même, ne pas conclure chaque question par une phrase rédigée fait perdre des points de présentation.
\end{enumerate}
\end{attention}

\courssub{Complément : linéarisation pour ajustement non linéaire (Culture)}

\begin{savaistu}[Ajustement non linéaire par linéarisation]
Le programme de Terminale technique se limite à l'ajustement \textbf{linéaire}. Cependant, en pratique technique (électrotechnique, mécanique, chimie), de nombreuses lois sont non linéaires. On peut parfois les \emph{linéariser} par un changement de variable pour réutiliser les outils vus ci-dessus.

\begin{center}
\begin{tabularx}{\linewidth}{|l|X|X|X|}
\hline
\rowcolor{popblueL} \textbf{Modèle} & \textbf{Équation} & \textbf{Changement de variable} & \textbf{Droite obtenue} \\
\hline
Puissance & $y = a x^b$ & $X = \ln x$, $Y = \ln y$ & $Y = b X + \ln a$ \\
\hline
Exponentielle & $y = a e^{bx}$ & $X = x$, $Y = \ln y$ & $Y = b X + \ln a$ \\
\hline
Inverse & $y = \dfrac{a}{x} + b$ & $X = \dfrac{1}{x}$, $Y = y$ & $Y = a X + b$ \\
\hline
\end{tabularx}
\end{center}

\textbf{Démarche :}
\begin{enumerate}
\item Appliquer la transformation sur les données $(x_i, y_i) \to (X_i, Y_i)$.
\item Tracer le nuage $(X_i, Y_i)$ et calculer $r$ : si $|r| \approx 1$, la linéarisation est validée.
\item Déterminer la droite de régression $Y = a X + b$ des moindres carrés.
\item Revenir au modèle initial en identifiant les paramètres.
\end{enumerate}

\begin{attention}[Piège]
Le coefficient de corrélation $r$ calculé sur les variables transformées ($X, Y$) mesure la qualité de l'ajustement \textbf{linéarisé}, pas celle du modèle initial sur les données brutes. De plus, la droite des moindres carrés sur $(X,Y)$ ne correspond pas exactement à la courbe des moindres carrés sur $(x,y)$ (critère minimisé différent). Cette méthode reste un outil d'approche commode en Terminale.
\end{attention}
\end{savaistu}

\newpage

\coursec{Exercices}

\courssub{Je vérifie mes connaissances}

\begin{exercice}[QCM : Tableau à double entrée et corrélation]
Parmi les affirmations suivantes, une seule est exacte. La recopier en justifiant brièvement.
\begin{enumerate}
\item Dans un tableau à double entrée, la somme des effectifs marginaux de la variable $X$ est toujours strictement supérieure à la somme des effectifs marginaux de la variable $Y$.
\item Le coefficient de corrélation linéaire $r$ est toujours compris entre $-1$ et $1$. Si $r = 0$, il n'existe \emph{aucune} relation entre $X$ et $Y$.
\item La droite de régression des moindres carrés $y$ en $x$ passe toujours par le point moyen $(\bar{x}, \bar{y})$.
\item La pente de la droite de régression $y = ax + b$ a le signe opposé à celui du coefficient de corrélation $r$.
\end{enumerate}
\end{exercice}

\begin{exercice}[Vrai / Faux justifié : Propriétés de la régression]
Dire si chaque affirmation est vraie ou fausse. Justifier chaque réponse par une phrase ou un contre-exemple.
\begin{enumerate}
\item Si le nuage de points est aligné parfaitement selon une droite de pente négative, alors $r = -1$.
\item La droite de régression $x$ en $y$ a pour équation $x = a'y + b'$ avec $a' = \frac{s_{xy}}{s_y^2}$. Les deux droites de régression ($y$ en $x$ et $x$ en $y$) sont toujours distinctes.
\item Le coefficient de détermination $R^2$ représente la part de la variance de $Y$ expliquée par le modèle linéaire. Il est toujours positif.
\item Pour prévoir une valeur de $Y$ pour une valeur de $X$ située \emph{en dehors} de l'intervalle des $x_i$ observés, on utilise l'équation de la droite de régression sans risque d'erreur.
\end{enumerate}
\end{exercice}

\begin{exercice}[Calculs directs : Covariance, variance, coefficient de corrélation]
On donne la série statistique à deux variables $(x_i, y_i)_{i=1..5}$ suivante :
\[
\begin{array}{c|ccccc}
x_i & 2 & 4 & 6 & 8 & 10 \\ \hline
y_i & 5 & 7 & 10 & 12 & 16
\end{array}
\]
\begin{enumerate}
\item Calculer les moyennes $\bar{x}$ et $\bar{y}$.
\item Calculer les variances $s_x^2$, $s_y^2$ et la covariance $s_{xy}$ (utiliser les formules algorithmiques $\frac{1}{n}\sum x_i^2 - \bar{x}^2$, etc.).
\item En déduire le coefficient de corrélation linéaire $r$. Arrondir à $10^{-3}$.
\item Commenter le sens et l'intensité de la liaison linéaire.
\end{enumerate}
\end{exercice}

\begin{exercice}[Lecture d'un tableau à double entrée]
Le tableau ci-dessous récapitule les notes (sur 20) en Mathématiques ($X$) et en Physique ($Y$) de 30 élèves d'une classe de Terminale.
\[
\begin{array}{|c|ccc|c|}\hline
X \backslash Y & [6;10[ & [10;14[ & [14;18[ & \text{Total} \\ \hline
\relax [6;10[ & 2 & 3 & 0 & 5 \\
\relax [10;14[ & 1 & 8 & 4 & 13 \\
\relax [14;18[ & 0 & 5 & 7 & 12 \\ \hline
\text{Total} & 3 & 16 & 11 & 30 \\ \hline
\end{array}
\]
\begin{enumerate}
\item Vérifier les effectifs marginaux. Quelle est la fréquence de la modalité $[10;14[$ pour $X$ ?
\item On prend comme valeurs caractéristiques les milieux des classes. Calculer une valeur approchée de la moyenne conditionnelle de $Y$ sachant que $X \in [14;18[$.
\item Sans calculer $r$, que pouvez-vous dire du sens de la liaison entre $X$ et $Y$ ? Justifier en observant le tableau.
\end{enumerate}
\end{exercice}

\courssub{Je m'entraîne}

\begin{exercice}[Application complète : Données brutes (Consommation électrique)]
Un technicien relève la consommation électrique (en kWh) $Y$ d'un atelier en fonction de la température extérieure moyenne (en $^\circ$C) $X$ sur 6 jours.
\[
\begin{array}{c|cccccc}
X & 12 & 15 & 18 & 20 & 22 & 25 \\ \hline
Y & 180 & 165 & 150 & 140 & 130 & 115
\end{array}
\]
\begin{enumerate}
\item Représenter le nuage de points dans un repère orthonormé (échelle adaptée).
\item Calculer $\bar{x}$, $\bar{y}$, $s_x^2$, $s_y^2$, $s_{xy}$ et le coefficient de corrélation $r$. Interpréter.
\item Déterminer l'équation de la droite de régression des moindres carrés $y$ en $x$ (coefficients arrondis à $10^{-2}$).
\item Tracer cette droite sur le nuage de points.
\item Prédire la consommation pour une température de $16^\circ$C. Cette prévision est-elle fiable ? Justifier.
\end{enumerate}
\end{exercice}

\begin{exercice}[Tableau à double entrée : Effectifs partiels (Usinage)]
Dans un atelier d'usinage, on étudie la liaison entre le diamètre $X$ (en mm) d'une pièce et sa rugosité $Y$ (en $\mu$m). Le tableau à double entrée suivant a été établi sur 50 pièces.
\[
\begin{array}{|c|ccc|c|}\hline
X \backslash Y & [1.0; 1.5[ & [1.5; 2.0[ & [2.0; 2.5[ & \text{Total} \\ \hline
\relax [8.0; 8.5[ & 4 & 6 & 0 & 10 \\
\relax [8.5; 9.0[ & 2 & 10 & 8 & 20 \\
\relax [9.0; 9.5[ & 0 & 5 & 15 & 20 \\ \hline
\text{Total} & 6 & 21 & 23 & 50 \\ \hline
\end{array}
\]
\begin{enumerate}
\item Compléter les effectifs marginaux (déjà faits). Calculer les fréquences marginales de $Y$.
\item En prenant les milieux des classes, calculer les moyennes $\bar{x}$ et $\bar{y}$.
\item Construire le tableau de calcul pour déterminer $s_{xy}$, $s_x^2$, $s_y^2$ et $r$. (On pourra regrouper les calculs par cellules).
\item Déduire l'équation de la droite de régression $y$ en $x$.
\item Pour un diamètre de $8.7$ mm, estimer la rugosité attendue.
\end{enumerate}
\end{exercice}

\begin{exercice}[Calculs à partir des sommes (Gestion de stock)]
Pour gérer son stock, un magasinier a relevé pendant $n=12$ mois le nombre de commandes $X$ et le niveau de stock $Y$ (en unités). Il a obtenu les sommes suivantes :
\[
\sum x_i = 360,\quad \sum y_i = 240,\quad \sum x_i^2 = 11\,800,\quad \sum y_i^2 = 5\,200,\quad \sum x_i y_i = 6\,800.
\]
\begin{enumerate}
\item Calculer $\bar{x}$, $\bar{y}$, $s_x^2$, $s_y^2$, $s_{xy}$.
\item Calculer le coefficient de corrélation $r$. Conclure sur la liaison.
\item Déterminer l'équation de la droite de régression $y$ en $x$.
\item Interpréter le coefficient directeur $a$ dans le contexte de l'exercice.
\item Si le magasinier prévoit 35 commandes le mois prochain, quel niveau de stock doit-il anticiper ?
\end{enumerate}
\end{exercice}

\begin{exercice}[Interprétation et validation (Maintenance préventive)]
Une entreprise de maintenance suit le nombre de pannes $Y$ survenues sur des machines en fonction de leur âge $X$ (en années). L'étude statistique sur 40 machines donne :
$r = 0,89$, $\bar{x} = 5,2$, $\bar{y} = 12,5$, $s_x = 2,1$, $s_y = 4,8$.
La droite de régression $y$ en $x$ est $y = 2,04x + 1,89$.
\begin{enumerate}
\item Calculer le coefficient de détermination $R^2$ et l'interpréter en pourcentage.
\item Une machine a 7 ans. Prédire son nombre de pannes.
\item Une machine neuve ($x=0$) aurait-elle $-1,89$ pannes selon le modèle ? Expliquer pourquoi cette prévision est absurde (notion d'\emph{extrapolation}).
\item Le nuage de points présente une courbure marquée (forme de banane). Le modèle linéaire est-il vraiment adapté ? Que proposez-vous ?
\end{enumerate}
\end{exercice}

\courssub{Je me prépare au Probatoire}

\begin{exercice}[Sujet type Probatoire 1 : Poids et Taille d'élèves (Données brutes)]
On a mesuré le poids $X$ (en kg) et la taille $Y$ (en cm) de 10 élèves d'une classe de Terminale.
\[
\begin{array}{c|cccccccccc}
\text{Élève} & 1 & 2 & 3 & 4 & 5 & 6 & 7 & 8 & 9 & 10 \\ \hline
X & 52 & 58 & 63 & 65 & 70 & 72 & 75 & 78 & 82 & 85 \\
Y & 160 & 165 & 168 & 170 & 172 & 175 & 178 & 180 & 183 & 186
\end{array}
\]
\begin{enumerate}
\item Construire un tableau à double entrée en regroupant les poids par classes de largeur 10 ($[50;60[$, $[60;70[$, $[70;80[$, $[80;90[$) et les tailles par classes de largeur 5 ($[155;165[$, $[165;175[$, $[175;185[$, $[185;195[$).
\item Calculer les moyennes $\bar{x}$ et $\bar{y}$ à partir des données brutes.
\item Calculer les variances $s_x^2$, $s_y^2$ et la covariance $s_{xy}$. En déduire le coefficient de corrélation linéaire $r$ (arrondi à $10^{-3}$).
\item Déterminer l'équation de la droite de régression des moindres carrés $y$ en $x$ (coefficients arrondis à $10^{-2}$).
\item Un nouvel élève pèse 65 kg. Estimer sa taille à l'aide du modèle.
\item Commenter la fiabilité de cette prévision en vous basant sur la valeur de $r$ et sur la position de $x=65$ par rapport à l'intervalle des observations.
\end{enumerate}
\end{exercice}

\begin{exercice}[Sujet type Probatoire 2 : Tableau double entrée (Effectifs partiels - Électricité)]
Un électricien étudie la liaison entre l'intensité $X$ (en A) traversant une résistance et la tension $Y$ (en V) à ses bornes. Il obtient le tableau à double entrée suivant (effectifs partiels) sur 60 mesures.
\[
\begin{array}{|c|cccc|c|}\hline
X \backslash Y & [10;12[ & [12;14[ & [14;16[ & [16;18[ & \text{Total} \\ \hline
\relax [1;2[ & 5 & 3 & 0 & 0 & 8 \\
\relax [2;3[ & 0 & 10 & 12 & 0 & 22 \\
\relax [3;4[ & 0 & 0 & 15 & 15 & 30 \\ \hline
\text{Total} & 5 & 13 & 27 & 15 & 60 \\ \hline
\end{array}
\]
\begin{enumerate}
\item Vérifier les totaux marginaux. Calculer les fréquences marginales de $X$ et $Y$.
\item En prenant les milieux des classes, calculer les moyennes $\bar{x}$ et $\bar{y}$.
\item Construire le tableau de calcul complet (colonnes : $x_i$, $y_j$, $n_{ij}$, $x_i n_{ij}$, $y_j n_{ij}$, $x_i^2 n_{ij}$, $y_j^2 n_{ij}$, $x_i y_j n_{ij}$) pour déterminer $s_x^2$, $s_y^2$, $s_{xy}$.
\item Calculer le coefficient de corrélation $r$. Préciser le sens et l'intensité de la liaison.
\item Déterminer l'équation de la droite de régression $y$ en $x$. Rappeler la loi physique (Loi d'Ohm) qui justifie ce modèle linéaire passant par l'origine (théoriquement). Le modèle trouvé passe-t-il par l'origine ? Commenter.
\end{enumerate}
\end{exercice}

\begin{exercice}[Sujet type Probatoire 3 : Problème ouvert - Optimisation du Chiffre d'Affaires (PME Bafoussam)]
Une PME de Bafoussam produit et vend un jus de fruit local. Le gérant veut fixer le prix de vente optimal pour maximiser son Chiffre d'Affaires (CA). Il a relevé les ventes des 8 derniers mois :
\[
\begin{array}{c|cccccccc}
\text{Prix } X \text{ (en FCFA)} & 500 & 550 & 600 & 650 & 700 & 750 & 800 & 850 \\ \hline
\text{Qté vendue } Y \text{ (en litres)} & 1200 & 1100 & 1020 & 930 & 850 & 760 & 680 & 600
\end{array}
\]
On admet que la liaison est linéaire et que la droite de régression $y$ en $x$ a pour équation : $y = -1,6x + 2000$ (coefficients arrondis). Le coefficient de corrélation est $r = -0,98$.
\begin{enumerate}
\item Interpréter le signe de $r$ et sa valeur proche de $-1$ dans le contexte commercial.
\item Calculer le coefficient de détermination $R^2$ et l'interpréter.
\item Exprimer le Chiffre d'Affaires mensuel $CA(x)$ (en FCFA) en fonction du prix de vente $x$, en utilisant le modèle linéaire pour la quantité vendue.
\item Étudier les variations de la fonction $CA(x)$ sur l'intervalle $[500; 850]$. (On rappelle que $CA(x) = x \times y_{\text{modèle}}$).
\item Déterminer le prix de vente $x_{\text{opt}}$ qui maximise le Chiffre d'Affaires et donner la valeur maximale $CA_{\text{max}}$.
\item Le gérant propose d'augmenter le prix à 1000 FCFA pour "gagner plus". En utilisant le modèle, calculer le CA prévu. Pourquoi cette prévision est-elle risquée ? (Notion d'extrapolation).
\item \textbf{Communication :} Rédiger une courte note (5-6 lignes) à l'attention du gérant pour lui recommander le prix optimal, en justifiant par les calculs et en précisant les limites du modèle.
\end{enumerate}
\end{exercice}

\newpage

\coursec{Corrigés des exercices}

\coursec{Corrigés détaillés — Statistiques à deux variables}

\begin{corrige}[Exercice 1 — QCM : Tableau à double entrée et corrélation]
\emph{Analyse de chaque proposition :}

\begin{enumerate}
\item \textbf{Fausse.} Dans un tableau à double entrée, la somme des effectifs marginaux de la variable $X$ est égale à l'effectif total $N$, qui est aussi la somme des effectifs marginaux de la variable $Y$. Elles sont donc égales, non strictement supérieures.
\item \textbf{Fausse.} La première partie est exacte : $r \in [-1; 1]$. En revanche, $r = 0$ signifie qu'il n'y a \emph{aucune liaison linéaire} entre $X$ et $Y$ ; il peut exister une liaison non linéaire (parabolique, exponentielle, etc.).
\item \textbf{Vraie.} C'est une propriété fondamentale de la droite de régression des moindres carrés $y$ en $x$ : elle passe toujours par le point moyen $(\bar{x}, \bar{y})$.
\item \textbf{Fausse.} La pente de la droite de régression $y$ en $x$ est $a = \frac{s_{xy}}{s_x^2}$. Le coefficient de corrélation est $r = \frac{s_{xy}}{s_x s_y}$. Comme les variances $s_x^2$ et $s_y^2$ sont positives, $a$ et $r$ ont \emph{le même signe} (celui de la covariance $s_{xy}$).
\end{enumerate}
\textbf{Bonne réponse :} La proposition 3.
\end{corrige}

\begin{corrige}[Exercice 2 — Vrai / Faux justifié : Propriétés de la régression]
\begin{enumerate}
\item \textbf{Vrai.} Si les points sont parfaitement alignés sur une droite de pente négative, la liaison linéaire est parfaite et négative, donc $r = -1$.
\item \textbf{Faux.} Les deux droites de régression ($y$ en $x$ et $x$ en $y$) coïncident \emph{si et seulement si} $|r| = 1$ (liaison linéaire parfaite). Si $|r| < 1$, elles sont distinctes et se coupent en $(\bar{x}, \bar{y})$.
\item \textbf{Vrai.} Le coefficient de détermination $R^2 = r^2$ est toujours compris entre $0$ et $1$. Il mesure la proportion de la variance de $Y$ expliquée par le modèle linéaire.
\item \textbf{Faux.} Prédire en dehors de l'intervalle des $x_i$ observés s'appelle de l'\emph{extrapolation}. C'est risqué car la relation linéaire observée sur l'échantillon n'est pas garantie en dehors de cet intervalle ; l'erreur de prévision peut être très grande.
\end{enumerate}
\end{corrige}

\begin{corrige}[Exercice 3 — Calculs directs : Covariance, variance, coefficient de corrélation]
Données : $n=5$, $(x_i) = (2, 4, 6, 8, 10)$, $(y_i) = (5, 7, 10, 12, 16)$.

\begin{enumerate}
\item \textbf{Moyennes :}
\[
\bar{x} = \frac{2+4+6+8+10}{5} = \frac{30}{5} = 6,
\qquad
\bar{y} = \frac{5+7+10+12+16}{5} = \frac{50}{5} = 10.
\]

\item \textbf{Variances et covariance (formules algorithmiques) :}
\[
\begin{aligned}
\sum x_i^2 &= 4 + 16 + 36 + 64 + 100 = 220, &
s_x^2 &= \frac{220}{5} - 6^2 = 44 - 36 = 8. \\[4pt]
\sum y_i^2 &= 25 + 49 + 100 + 144 + 256 = 574, &
s_y^2 &= \frac{574}{5} - 10^2 = 114,8 - 100 = 14,8. \\[4pt]
\sum x_i y_i &= 10 + 28 + 60 + 96 + 160 = 354, &
s_{xy} &= \frac{354}{5} - 6 \times 10 = 70,8 - 60 = 10,8.
\end{aligned}
\]

\item \textbf{Coefficient de corrélation linéaire $r$ :}
\[
r = \frac{s_{xy}}{\sqrt{s_x^2 s_y^2}} = \frac{10,8}{\sqrt{8 \times 14,8}} = \frac{10,8}{\sqrt{118,4}} \approx \frac{10,8}{10,881} \approx 0,993 \quad (\text{arrondi à } 10^{-3}).
\]

\item \textbf{Interprétation :} $r \approx 0,993$ est très proche de $+1$. Il existe une liaison linéaire \textbf{très forte} et \textbf{positive} entre $X$ et $Y$ : lorsque $X$ augmente, $Y$ augmente de manière presque proportionnelle.
\end{enumerate}
\end{corrige}

\begin{corrige}[Exercice 4 — Lecture d'un tableau à double entrée]
\begin{enumerate}
\item \textbf{Vérification des marginaux :}
Lignes : $2+3+0=5$ ; $1+8+4=13$ ; $0+5+7=12$. Total $5+13+12=30$.
Colonnes : $2+1+0=3$ ; $3+8+5=16$ ; $0+4+7=11$. Total $3+16+11=30$. \textbf{Correct.}
Fréquence de la modalité $[10;14[$ pour $X$ : $\frac{13}{30} \approx 0,433$ (soit $43,3\%$).

\item \textbf{Moyenne conditionnelle de $Y$ sachant $X \in [14;18[$ :}
Milieux des classes de $Y$ : $8$, $12$, $16$.
Effectifs de la ligne $X \in [14;18[$ (milieu $16$) : $n_{31}=0$, $n_{32}=5$, $n_{33}=7$. Total $12$.
\[
\bar{y}_{|X\in[14;18[} = \frac{0\times 8 + 5\times 12 + 7\times 16}{12} = \frac{60 + 112}{12} = \frac{172}{12} \approx 14,33.
\]

\item \textbf{Sens de la liaison :} On observe une concentration des effectifs sur la diagonale « bas-gauche / haut-droite » : les petites valeurs de $X$ (ligne 1) sont associées aux petites valeurs de $Y$ (colonne 1), et les grandes valeurs de $X$ (ligne 3) aux grandes valeurs de $Y$ (colonne 3). La liaison est \textbf{positive} : les notes en Maths et en Physique évoluent dans le même sens.
\end{enumerate}
\end{corrige}

\begin{corrige}[Exercice 5 — Application complète : Consommation électrique]
Données : $n=6$.
$X$ : 12, 15, 18, 20, 22, 25.
$Y$ : 180, 165, 150, 140, 130, 115.

\begin{enumerate}
\item \textbf{Nuage de points :} Repère orthonormé, $X$ en abscisse (échelle 1 cm pour 2°C), $Y$ en ordonnée (échelle 1 cm pour 10 kWh). Points alignés visuellement sur une droite décroissante.

\item \textbf{Calculs :}
\[
\begin{aligned}
\sum x_i &= 112, & \bar{x} &= \frac{112}{6} = \frac{56}{3} \approx 18,67. \\
\sum y_i &= 880, & \bar{y} &= \frac{880}{6} = \frac{440}{3} \approx 146,67. \\
\sum x_i^2 &= 2202, & s_x^2 &= \frac{2202}{6} - \left(\frac{56}{3}\right)^2 = 367 - \frac{3136}{9} = \frac{167}{9} \approx 18,56. \\
\sum y_i^2 &= 131850, & s_y^2 &= \frac{131850}{6} - \left(\frac{440}{3}\right)^2 = 21975 - \frac{193600}{9} = \frac{4175}{9} \approx 463,89. \\
\sum x_i y_i &= 15870, & s_{xy} &= \frac{15870}{6} - \frac{56}{3}\times\frac{440}{3} = 2645 - \frac{24640}{9} = -\frac{835}{9} \approx -92,78.
\end{aligned}
\]
\[
r = \frac{s_{xy}}{\sqrt{s_x^2 s_y^2}} = \frac{-835/9}{\sqrt{(167/9)(4175/9)}} = \frac{-835}{\sqrt{167 \times 4175}}.
\]
Or $4175 = 25 \times 167$, donc $\sqrt{167 \times 4175} = \sqrt{25 \times 167^2} = 5 \times 167 = 835$.
\[
\boxed{r = -1}.
\]
\textbf{Interprétation :} Corrélation linéaire parfaite négative. Les points sont exactement alignés sur une droite décroissante. La consommation diminue linéairement quand la température augmente.

\item \textbf{Droite de régression $y$ en $x$ :}
\[
a = \frac{s_{xy}}{s_x^2} = \frac{-835/9}{167/9} = -5, \qquad
b = \bar{y} - a\bar{x} = \frac{440}{3} - (-5)\times\frac{56}{3} = \frac{440+280}{3} = 240.
\]
Équation : $\boxed{y = -5x + 240}$ (coefficients exacts, arrondi $10^{-2}$ inchangé).

\item \textbf{Tracé :} La droite passe par $(12;180)$ et $(25;115)$ (ou par le point moyen $(18,67; 146,67)$). Elle recouvre exactement les points du nuage.

\item \textbf{Prévision pour $x=16^\circ$C :} $y = -5\times 16 + 240 = 160$ kWh.
\textbf{Fiabilité :} Excellente. C'est une \emph{interpolation} ($16 \in [12; 25]$) et le modèle est parfait ($r=-1$, $R^2=1$).
\end{enumerate}
\end{corrige}

\begin{corrige}[Exercice 6 — Tableau à double entrée : Effectifs partiels (Usinage)]
$n=50$. Milieux : $X$ : 8,25 ; 8,75 ; 9,25. $Y$ : 1,25 ; 1,75 ; 2,25.

\begin{enumerate}
\item \textbf{Fréquences marginales de $Y$ :}
$f_{Y_1} = 6/50 = 0,12$ ; $f_{Y_2} = 21/50 = 0,42$ ; $f_{Y_3} = 23/50 = 0,46$.

\item \textbf{Moyennes :}
\[
\bar{x} = \frac{8,25\times10 + 8,75\times20 + 9,25\times20}{50} = \frac{82,5 + 175 + 185}{50} = \frac{442,5}{50} = 8,85 \text{ mm}.
\]
\[
\bar{y} = \frac{1,25\times6 + 1,75\times21 + 2,25\times23}{50} = \frac{7,5 + 36,75 + 51,75}{50} = \frac{96}{50} = 1,92 \ \mu\text{m}.
\]

\item \textbf{Tableau de calcul (sommes par cellules) :}
\begin{center}
\begin{tabular}{|c|c|c|c|c|c|c|c|c|}\hline
$x_i$ & $y_j$ & $n_{ij}$ & $x_i n_{ij}$ & $y_j n_{ij}$ & $x_i^2 n_{ij}$ & $y_j^2 n_{ij}$ & $x_i y_j n_{ij}$ \\ \hline
8,25 & 1,25 & 4 & 33,00 & 5,00 & 272,25 & 6,25 & 41,25 \\
8,25 & 1,75 & 6 & 49,50 & 10,50 & 408,375 & 18,375 & 86,625 \\
8,75 & 1,25 & 2 & 17,50 & 2,50 & 153,125 & 3,125 & 21,875 \\
8,75 & 1,75 & 10 & 87,50 & 17,50 & 765,625 & 30,625 & 153,125 \\
8,75 & 2,25 & 8 & 70,00 & 18,00 & 612,50 & 40,50 & 157,50 \\
9,25 & 1,75 & 5 & 46,25 & 8,75 & 427,8125 & 15,3125 & 80,9375 \\
9,25 & 2,25 & 15 & 138,75 & 33,75 & 1283,4375 & 75,9375 & 312,1875 \\ \hline
\textbf{Total} & & \textbf{50} & \textbf{442,5} & \textbf{96} & \textbf{3923,125} & \textbf{190,125} & \textbf{853,5} \\ \hline
\end{tabular}
\end{center}

\[
\begin{aligned}
s_x^2 &= \frac{3923,125}{50} - 8,85^2 = 78,4625 - 78,3225 = 0,14. \\
s_y^2 &= \frac{190,125}{50} - 1,92^2 = 3,8025 - 3,6864 = 0,1161. \\
s_{xy} &= \frac{853,5}{50} - 8,85 \times 1,92 = 17,07 - 16,992 = 0,078. \\
r &= \frac{0,078}{\sqrt{0,14 \times 0,1161}} = \frac{0,078}{\sqrt{0,016254}} \approx \frac{0,078}{0,1275} \approx 0,612.
\end{aligned}
\]

\item \textbf{Droite de régression $y$ en $x$ :}
\[
a = \frac{s_{xy}}{s_x^2} = \frac{0,078}{0,14} \approx 0,5571, \qquad
b = \bar{y} - a\bar{x} = 1,92 - 0,5571 \times 8,85 \approx 1,92 - 4,930 \approx -3,01.
\]
Équation : $\boxed{y = 0,557x - 3,01}$ (arrondi $10^{-3}$ pour $a$, $10^{-2}$ pour $b$).

\item \textbf{Estimation pour $x = 8,7$ mm :}
$y = 0,557 \times 8,7 - 3,01 \approx 4,846 - 3,01 = 1,836 \ \mu\text{m}$.
Rugosité attendue : $\approx 1,84 \ \mu\text{m}$.
\end{enumerate}
\end{corrige}

\begin{corrige}[Exercice 7 — Calculs à partir des sommes (Gestion de stock)]
$n=12$, $\sum x_i=360$, $\sum y_i=240$, $\sum x_i^2=11800$, $\sum y_i^2=5200$, $\sum x_i y_i=6800$.

\begin{enumerate}
\item \textbf{Paramètres :}
\[
\bar{x} = \frac{360}{12} = 30, \quad \bar{y} = \frac{240}{12} = 20.
\]
\[
s_x^2 = \frac{11800}{12} - 30^2 = 983,\overline{3} - 900 = \frac{250}{3} \approx 83,33.
\]
\[
s_y^2 = \frac{5200}{12} - 20^2 = 433,\overline{3} - 400 = \frac{100}{3} \approx 33,33.
\]
\[
s_{xy} = \frac{6800}{12} - 30 \times 20 = 566,\overline{6} - 600 = -\frac{100}{3} \approx -33,33.
\]

\item \textbf{Coefficient de corrélation :}
\[
r = \frac{-100/3}{\sqrt{(250/3)(100/3)}} = \frac{-100}{\sqrt{25000}} = \frac{-100}{50\sqrt{10}} = -\frac{2}{\sqrt{10}} \approx -0,632.
\]
\textbf{Conclusion :} Liaison linéaire \textbf{négative} et \textbf{modérée} (plus de commandes $\rightarrow$ stock plus bas).

\item \textbf{Droite de régression $y$ en $x$ :}
\[
a = \frac{s_{xy}}{s_x^2} = \frac{-100/3}{250/3} = -0,4, \qquad
b = \bar{y} - a\bar{x} = 20 - (-0,4)\times 30 = 32.
\]
Équation : $\boxed{y = -0,4x + 32}$.

\item \textbf{Interprétation de $a = -0,4$ :} En moyenne, chaque commande supplémentaire entraîne une diminution du niveau de stock de $0,4$ unité (soit 2 unités pour 5 commandes).

\item \textbf{Prévision pour 35 commandes :} $y = -0,4 \times 35 + 32 = -14 + 32 = 18$ unités.
\end{enumerate}
\end{corrige}

\begin{corrige}[Exercice 8 — Interprétation et validation (Maintenance préventive)]
Données : $r=0,89$, $\bar{x}=5,2$, $\bar{y}=12,5$, $s_x=2,1$, $s_y=4,8$, droite $y = 2,04x + 1,89$.

\begin{enumerate}
\item \textbf{Coefficient de détermination :} $R^2 = r^2 = 0,89^2 = 0,7921 = 79,21\%$.
\textbf{Interprétation :} $79,21\%$ de la variabilité du nombre de pannes ($Y$) est expliquée par l'âge des machines ($X$) via le modèle linéaire. Les $20,79\%$ restants sont dus à d'autres facteurs (conditions d'utilisation, maintenance, hasard...).

\item \textbf{Prédiction pour $x=7$ ans :} $y = 2,04 \times 7 + 1,89 = 14,28 + 1,89 = 16,17$ pannes.

\item \textbf{Machine neuve ($x=0$) :} Le modèle donne $y = 2,04 \times 0 + 1,89 = 1,89$ pannes (et non $-1,89$).
\textbf{Absurdité de la prévision :} C'est une \emph{extrapolation} (l'âge 0 est hors de l'intervalle des âges observés, probablement $>1$ an). De plus, une machine neuve devrait avoir 0 panne ; l'ordonnée à l'origine $b=1,89$ n'a pas de sens physique ici, elle est un artefact mathématique de l'ajustement linéaire sur l'intervalle observé.

\item \textbf{Adéquation du modèle :} Si le nuage présente une courbure marquée (« forme de banane »), le modèle \textbf{linéaire n'est pas adapté} (biais de spécification). Les résidus présenteront une structure systématique.
\textbf{Proposition :} Essayer un ajustement non linéaire (modèle quadratique $y = ax^2+bx+c$, ou exponentiel, ou linéarisation par transformation de variables, ex: $\ln y$ en fonction de $x$).
\end{enumerate}
\end{corrige}

\begin{corrige}[Exercice 9 — Sujet type Probatoire 1 : Poids et Taille]
Données brutes $n=10$.
$X$ (kg) : 52, 58, 63, 65, 70, 72, 75, 78, 82, 85. $\sum X = 700$, $\bar{x}=70$.
$Y$ (cm) : 160, 165, 168, 170, 172, 175, 178, 180, 183, 186. $\sum Y = 1737$, $\bar{y}=173,7$.
$\sum X^2 = 50004$, $\sum Y^2 = 302327$, $\sum XY = 122370$.

\begin{enumerate}
\item \textbf{Tableau à double entrée (classes $X$: larg. 10 ; $Y$: larg. 5) :}
\begin{center}
\begin{tabular}{|c|cccc|c|}\hline
$X \backslash Y$ & $[155;165[$ & $[165;175[$ & $[175;185[$ & $[185;195[$ & Total \\ \hline
$[50;60[$ & 1 (52;160) & 1 (58;165) & 0 & 0 & 2 \\
$[60;70[$ & 0 & 2 (63;168; 65;170) & 0 & 0 & 2 \\
$[70;80[$ & 0 & 1 (70;172) & 3 (72;175; 75;178; 78;180) & 0 & 4 \\
$[80;90[$ & 0 & 0 & 1 (82;183) & 1 (85;186) & 2 \\ \hline
Total & 1 & 4 & 4 & 1 & 10 \\ \hline
\end{tabular}
\end{center}
\emph{Note : la valeur 175 appartient à la classe $[175;185[$ (convention intervalle fermé à gauche).}

\item \textbf{Moyennes (données brutes) :} $\bar{x} = 70$ kg, $\bar{y} = 173,7$ cm.

\item \textbf{Variances, covariance, $r$ :}
\[
s_x^2 = \frac{50004}{10} - 70^2 = 5000,4 - 4900 = 100,4.
\]
\[
s_y^2 = \frac{302327}{10} - 173,7^2 = 30232,7 - 30171,69 = 61,01.
\]
\[
s_{xy} = \frac{122370}{10} - 70 \times 173,7 = 12237 - 12159 = 78.
\]
\[
r = \frac{78}{\sqrt{100,4 \times 61,01}} = \frac{78}{\sqrt{6125,404}} \approx \frac{78}{78,265} \approx \boxed{0,997}.
\]
Liaison linéaire très forte, positive.

\item \textbf{Droite de régression $y$ en $x$ :}
\[
a = \frac{78}{100,4} \approx 0,7769, \quad b = 173,7 - 0,7769 \times 70 \approx 119,32.
\]
Équation : $\boxed{y = 0,78x + 119,32}$ (arrondi $10^{-2}$).

\item \textbf{Estimation pour $x=65$ kg :} $y = 0,78 \times 65 + 119,32 = 50,7 + 119,32 = \boxed{170,02 \text{ cm}}$.

\item \textbf{Fiabilité :} Très bonne. $r \approx 0,997$ (liaison quasi parfaite). $x=65$ est à l'intérieur de l'intervalle d'observation $[52; 85]$ (interpolation). Le modèle est fiable ici.
\end{enumerate}

\textbf{Barème indicatif (sur 20) :}
Tableau double entrée (3 pts), Moyennes (2 pts), Variances/Covariance/$r$ (5 pts), Droite régression (3 pts), Prévision (2 pts), Commentaire fiabilité (3 pts), Présentation/Rigueur (2 pts).

\textbf{Points de vigilance :} Utiliser les données brutes pour les moyennes/variances (pas les classes). Arrondis demandés. Justifier la fiabilité par $r$ ET position de $x$ (interpolation).
\end{corrige}

\begin{corrige}[Exercice 10 — Sujet type Probatoire 2 : Tableau double entrée (Électricité)]
$n=60$. Milieux $X$: 1,5 ; 2,5 ; 3,5. Milieux $Y$: 11 ; 13 ; 15 ; 17.

\begin{enumerate}
\item \textbf{Marginaux vérifiés.} Fréquences $X$ : $8/60 \approx 0,133$ ; $22/60 \approx 0,367$ ; $30/60 = 0,5$.
Fréquences $Y$ : $5/60 \approx 0,083$ ; $13/60 \approx 0,217$ ; $27/60 = 0,45$ ; $15/60 = 0,25$.

\item \textbf{Moyennes :}
\[
\bar{x} = \frac{1,5\times8 + 2,5\times22 + 3,5\times30}{60} = \frac{12 + 55 + 105}{60} = \frac{172}{60} \approx 2,867 \text{ A}.
\]
\[
\bar{y} = \frac{11\times5 + 13\times13 + 15\times27 + 17\times15}{60} = \frac{55 + 169 + 405 + 255}{60} = \frac{884}{60} \approx 14,733 \text{ V}.
\]

\item \textbf{Tableau de calcul complet :}
\begin{center}
\begin{tabular}{|c|c|c|c|c|c|c|c|c|}\hline
$x_i$ & $y_j$ & $n_{ij}$ & $x_i n_{ij}$ & $y_j n_{ij}$ & $x_i^2 n_{ij}$ & $y_j^2 n_{ij}$ & $x_i y_j n_{ij}$ \\ \hline
1,5 & 11 & 5 & 7,5 & 55 & 11,25 & 605 & 82,5 \\
1,5 & 13 & 3 & 4,5 & 39 & 6,75 & 507 & 58,5 \\
2,5 & 13 & 10 & 25 & 130 & 62,5 & 1690 & 325 \\
2,5 & 15 & 12 & 30 & 180 & 75 & 2700 & 450 \\
3,5 & 15 & 15 & 52,5 & 225 & 183,75 & 3375 & 787,5 \\
3,5 & 17 & 15 & 52,5 & 255 & 183,75 & 4335 & 892,5 \\ \hline
\textbf{Total} & & \textbf{60} & \textbf{172} & \textbf{884} & \textbf{523} & \textbf{13212} & \textbf{2596} \\ \hline
\end{tabular}
\end{center}

\[
s_x^2 = \frac{523}{60} - \left(\frac{172}{60}\right)^2 \approx 8,7167 - 8,2178 = 0,4989.
\]
\[
s_y^2 = \frac{13212}{60} - \left(\frac{884}{60}\right)^2 = 220,2 - 217,071 = 3,1289.
\]
\[
s_{xy} = \frac{2596}{60} - \frac{172}{60}\times\frac{884}{60} \approx 43,2667 - 42,3022 = 0,9645.
\]

\item \textbf{Coefficient de corrélation :}
\[
r = \frac{0,9645}{\sqrt{0,4989 \times 3,1289}} \approx \frac{0,9645}{1,2494} \approx \boxed{0,772}.
\]
Liaison \textbf{positive} et \textbf{forte} (proche de 1).

\item \textbf{Droite de régression $y$ en $x$ :}
\[
a = \frac{0,9645}{0,4989} \approx 1,933, \quad b = 14,733 - 1,933 \times 2,867 \approx 9,19.
\]
Équation : $\boxed{y = 1,93x + 9,19}$ (arrondi $10^{-2}$).

\textbf{Loi d'Ohm :} $U = R I$ (tension = résistance $\times$ intensité). Théoriquement, la droite passe par l'origine $(0,0)$ (proportionnalité).
\textbf{Modèle trouvé :} $y = 1,93x + 9,19$. L'ordonnée à l'origine $9,19 \neq 0$. Le modèle \textbf{ne passe pas par l'origine}.
\textbf{Commentaire :} Cela s'explique par les erreurs de mesure, les résistances de contact, les fils non parfaits, ou le fait que l'ajustement linéaire par moindres carrés minimise les erreurs verticales sur l'intervalle observé sans contraindre le passage par l'origine. Une régression forcée par l'origine donnerait une pente $R \approx \bar{y}/\bar{x} \approx 5,14 \Omega$, différente de $1,93 \Omega$.
\end{enumerate}

\textbf{Barème indicatif (sur 20) :}
Marginaux/Fréquences (3 pts), Moyennes (3 pts), Tableau calcul complet (5 pts), $r$ + interprétation (3 pts), Droite régression (3 pts), Loi d'Ohm + commentaire origine (3 pts).

\textbf{Points de vigilance :} Tableau de calcul avec toutes les colonnes demandées. Utiliser les milieux. Préciser « liaison positive forte ». Citer explicitement la loi d'Ohm $U=RI$. Expliquer l'écart à l'origine par les imperfections réelles / méthode des moindres carrés non contrainte.
\end{corrige}

\begin{corrige}[Exercice 11 — Sujet type Probatoire 3 : Optimisation CA (PME Bafoussam)]
Modèle admis : $y = -1,6x + 2000$, $r = -0,98$, $x \in [500; 850]$.

\begin{enumerate}
\item \textbf{Interprétation de $r$ :} $r = -0,98$ indique une liaison linéaire \textbf{très forte} et \textbf{négative}. Dans le contexte : quand le prix de vente augmente, la quantité vendue diminue de manière quasi proportionnelle. Les clients sont très sensibles au prix.

\item \textbf{Coefficient de détermination :} $R^2 = r^2 = 0,9604 = 96,04\%$.
\textbf{Interprétation :} $96,04\%$ de la variabilité des quantités vendues est expliquée par la variation du prix via ce modèle linéaire. Le modèle est très performant sur l'intervalle observé.

\item \textbf{Chiffre d'Affaires $CA(x)$ :}
$CA(x) = x \times y_{\text{modèle}} = x(-1,6x + 2000) = \boxed{-1,6x^2 + 2000x}$ (en FCFA).

\item \textbf{Variations de $CA(x)$ sur $[500; 850]$ :}
$CA(x)$ est un trinôme du second degré, coefficient de $x^2$ négatif ($-1,6$), donc parabole concave, maximum au sommet.
Dérivée : $CA'(x) = -3,2x + 2000$.
$CA'(x) = 0 \Leftrightarrow x = \frac{2000}{3,2} = 625$.
$CA'(x) > 0$ sur $[500; 625[$ (croissant), $CA'(x) < 0$ sur $]625; 850]$ (décroissant).
\begin{center}
\begin{tabular}{|c|ccccc|}\hline
$x$ & 500 & & 625 & & 850 \\ \hline
$CA'(x)$ & & $+$ & 0 & $-$ & \\ \hline
$CA(x)$ & 600\,000 & $\nearrow$ & 625\,000 & $\searrow$ & 544\,000 \\ \hline
\end{tabular}
\end{center}

\item \textbf{Optimum :} Prix optimal $x_{\text{opt}} = \boxed{625 \text{ FCFA}}$.
CA maximal $CA_{\text{max}} = CA(625) = -1,6 \times 625^2 + 2000 \times 625 = 625 \times (-1000 + 2000) = \boxed{625\,000 \text{ FCFA}}$.

\item \textbf{Extrapolation à $x=1000$ :}
$y = -1,6 \times 1000 + 2000 = 400$ litres.
$CA(1000) = 1000 \times 400 = 400\,000 \text{ FCFA}$.
\textbf{Risque :} $x=1000$ est \textbf{hors de l'intervalle d'observation} $[500; 850]$ (extrapolation). La relation linéaire prix-demande n'est pas garantie au-delà : le marché peut saturer, la demande devenir nulle, ou la concurrence réagir. Le modèle surestime probablement la demande (et le CA) à ce prix.

\item \textbf{Communication (Note au gérant) :}
\begin{quote}
\textbf{Objet : Recommandation de prix de vente optimal pour le jus de fruit.}

Monsieur le Gérant,

L'analyse statistique des 8 derniers mois montre une relation linéaire très forte ($r=-0,98$) entre le prix et la quantité vendue. Le modèle $q = -1,6p + 2000$ permet de modéliser le Chiffre d'Affaires $CA(p) = -1,6p^2 + 2000p$. Ce CA est maximal pour un prix de \textbf{625 FCFA}, atteignant \textbf{625 000 FCFA/mois}.

Augmenter le prix à 1000 FCFA ferait chuter le CA prévu à 400 000 FCFA, mais cette prévision est peu fiable car elle sort de la plage de données observées (extrapolation). Je vous recommande donc de fixer le prix à \textbf{625 FCFA} pour maximiser votre revenu, tout en surveillant la réaction du marché.
\end{quote}
\end{enumerate}

\textbf{Barème indicatif (sur 20) :}
Interprétation $r$ (2 pts), $R^2$ (2 pts), Expression $CA(x)$ (2 pts), Étude variations + tableau (4 pts), Optimum + $CA_{max}$ (3 pts), Extrapolation $x=1000$ (3 pts), Note de synthèse (4 pts).

\textbf{Points de vigilance :} Ne pas confondre $r$ et $R^2$. Tableau de variations obligatoire avec 3 lignes, 3 colonnes (valeurs) + 2 colonnes (intervalles) = 5 colonnes ou format standard. Justifier l'optimum par la dérivée ou sommet parabole. Note : ton professionnel, chiffres clés, limite (extrapolation) citée.
\end{corrige}

\newpage

\coursec{Fiche bilan}

\begin{popfigure}
\begin{center}\tcbox[colback=popblue,colframe=popblue,coltext=white,arc=9pt,boxsep=4pt,left=6pt,right=6pt]{\bfseries\small Statistiques à deux variables}\end{center}
\vspace{-2pt}
\begin{tcbraster}[raster columns=3,raster equal height=rows,raster column skip=6pt,raster row skip=6pt,enhanced,blankest]
\begin{tcolorbox}[enhanced,colback=poporange,colframe=poporange,coltext=white,boxrule=0.9pt,arc=3pt,left=4pt,right=4pt,top=3pt,bottom=3pt,boxsep=1pt,halign=left]\bfseries\scriptsize Tableau à double entrée\end{tcolorbox}
\begin{tcolorbox}[enhanced,colback=popgreen,colframe=popgreen,coltext=white,boxrule=0.9pt,arc=3pt,left=4pt,right=4pt,top=3pt,bottom=3pt,boxsep=1pt,halign=left]\bfseries\scriptsize Nuage de points et point moyen $G$\end{tcolorbox}
\begin{tcolorbox}[enhanced,colback=poppurple,colframe=poppurple,coltext=white,boxrule=0.9pt,arc=3pt,left=4pt,right=4pt,top=3pt,bottom=3pt,boxsep=1pt,halign=left]\bfseries\scriptsize Coefficient de corrélation $r$\end{tcolorbox}
\begin{tcolorbox}[enhanced,colback=popgold,colframe=popgold,coltext=white,boxrule=0.9pt,arc=3pt,left=4pt,right=4pt,top=3pt,bottom=3pt,boxsep=1pt,halign=left]\bfseries\scriptsize Droite des moindres carrés $y=ax+b$\end{tcolorbox}
\begin{tcolorbox}[enhanced,colback=poppink,colframe=poppink,coltext=white,boxrule=0.9pt,arc=3pt,left=4pt,right=4pt,top=3pt,bottom=3pt,boxsep=1pt,halign=left]\bfseries\scriptsize Prévisions et validation\end{tcolorbox}
\begin{tcolorbox}[enhanced,colback=popmuted,colframe=popmuted,coltext=white,boxrule=0.9pt,arc=3pt,left=4pt,right=4pt,top=3pt,bottom=3pt,boxsep=1pt,halign=left]\bfseries\scriptsize Linéarisation (culture)\end{tcolorbox}
\end{tcbraster}

\legende{Carte mentale de la leçon : statistiques à deux variables}
\end{popfigure}

\begin{bilan}
\textbf{Définitions clés}
\begin{itemize}
  \item \cle{Série statistique à deux variables} : on étudie simultanément deux caractères $x$ et $y$ sur une même population de $n$ individus ; les données forment des couples $(x_i\,;\,y_i)$.
  \item \cle{Tableau à double entrée} : tableau croisant les valeurs (ou classes) de $x$ en lignes et celles de $y$ en colonnes ; la case $(i,j)$ contient l'effectif $n_{ij}$.
  \item \cle{Nuage de points} : ensemble des points $M_i(x_i\,;\,y_i)$ placés dans un repère.
  \item \cle{Point moyen} : $G(\bar{x}\,;\,\bar{y})$ avec $\bar{x}=\dfrac{1}{n}\sum x_i$ et $\bar{y}=\dfrac{1}{n}\sum y_i$.
  \item \cle{Ajustement linéaire} : remplacer le nuage par une droite passant « au plus près » des points, pour modéliser et prévoir.
\end{itemize}

\textbf{Formules à connaître par cœur}
\begin{center}
\begin{tabularx}{\linewidth}{|l|X|}
\hline
\rowcolor{popblueL} \textbf{Grandeur} & \textbf{Formule} \\
\hline
Covariance & $\mathrm{cov}(x,y)=\dfrac{1}{n}\sum x_i y_i-\bar{x}\,\bar{y}$ \\
\hline
Variances & $V(x)=\dfrac{1}{n}\sum x_i^2-\bar{x}^2$ ; \quad $V(y)=\dfrac{1}{n}\sum y_i^2-\bar{y}^2$ \\
\hline
Coefficient de corrélation & $r=\dfrac{\mathrm{cov}(x,y)}{\sqrt{V(x)\,V(y)}}$, avec $-1\le r\le 1$ \\
\hline
Droite de régression de $y$ en $x$ & $y=ax+b$ avec $a=\dfrac{\mathrm{cov}(x,y)}{V(x)}$ et $b=\bar{y}-a\bar{x}$ \\
\hline
\end{tabularx}
\end{center}

\textbf{Méthodes express}
\begin{itemize}
  \item \cle{Lecture d'un tableau à double entrée} : effectif marginal d'une ligne = somme de la ligne ; effectif total $n$ = somme de tous les $n_{ij}$.
  \item \cle{Interpréter $r$} : $|r|$ proche de $1$ (en pratique $|r|\ge 0{,}87$) $\Rightarrow$ forte corrélation linéaire, ajustement affine justifié ; $r$ proche de $0$ $\Rightarrow$ pas de liaison linéaire.
  \item \cle{Tracer la droite des moindres carrés} : elle passe toujours par le point moyen $G$ ; placer $G$ et un autre point calculé avec l'équation.
  \item \cle{Prévision} : remplacer $x$ par la valeur voulue dans $y=ax+b$ (interpolation fiable, extrapolation avec prudence).
  \item \cle{Linéarisation (culture)} : si le nuage n'est pas rectiligne, poser un changement de variable ($Y=\ln y$, $X=x^2$, ...) pour ramener à un nuage aligné.
\end{itemize}
\end{bilan}

\begin{aretenir}[Checklist avant le Probatoire]
\begin{itemize}
  \item[$\square$] Je sais lire et compléter un tableau à double entrée (effectifs, marges, total).
  \item[$\square$] Je sais représenter un nuage de points dans un repère bien choisi.
  \item[$\square$] Je sais calculer les moyennes $\bar{x}$ et $\bar{y}$ et placer le point moyen $G$.
  \item[$\square$] Je sais calculer la covariance $\mathrm{cov}(x,y)$ sans oublier de soustraire $\bar{x}\,\bar{y}$.
  \item[$\square$] Je sais calculer les variances $V(x)$ et $V(y)$ avec la formule de König-Huygens.
  \item[$\square$] Je sais calculer le coefficient de corrélation $r$ et vérifier que $-1\le r\le 1$.
  \item[$\square$] Je sais interpréter $r$ : décider si un ajustement affine est pertinent.
  \item[$\square$] Je sais déterminer l'équation $y=ax+b$ de la droite des moindres carrés.
  \item[$\square$] Je sais que la droite de régression passe par $G(\bar{x}\,;\,\bar{y})$.
  \item[$\square$] Je sais utiliser l'équation pour faire une prévision et rédiger une conclusion en langage courant.
  \item[$\square$] Je vérifie mes calculs à la calculatrice (mode statistiques à deux variables).
  \item[$\square$] Je rédige et justifie chaque étape de ma démarche, comme exigé à l'examen.
\end{itemize}
\end{aretenir}

\findecours

\end{document}
